% PDF version of the blueprint (`leanblueprint pdf`, LuaLaTeX via latexmkrc).
% The content is generated by scripts/blueprint_gen.py (documentation plan B1-B12).

\documentclass[a4paper]{report}

\usepackage{geometry}
\usepackage{expl3}
\usepackage{amssymb, amsthm, mathtools}
\usepackage[unicode,colorlinks=true,linkcolor=blue,urlcolor=magenta,citecolor=blue]{hyperref}
\usepackage[warnings-off={mathtools-colon,mathtools-overbracket}]{unicode-math}

% The Markdown text keeps its Unicode mathematics (as the Lean docstrings, D5); a fallback chain
% supplies the glyphs that the text and mono fonts lack. All fonts are loaded by file name from
% TeX Live (packages stix2-otf, dejavu, gnu-freefont), never from the system, so that the PDF
% builds the same on macOS and on the Linux runner of the Pages workflow.
\directlua{luaotfload.add_fallback("bpfb",
  {"file:STIXTwoMath-Regular.otf:mode=harf;", "file:DejaVuSans.ttf:mode=harf;",
   "file:FreeSerif.otf:mode=harf;", "file:FreeMono.otf:mode=harf;"})}
\setmainfont{STIXTwoText}[
  Extension=.otf, UprightFont=*-Regular, BoldFont=*-Bold,
  ItalicFont=*-Italic, BoldItalicFont=*-BoldItalic, RawFeature={fallback=bpfb}]
\setmathfont{STIXTwoMath-Regular.otf}
\setmonofont{DejaVuSansMono}[
  Extension=.ttf, UprightFont=*, BoldFont=*-Bold, ItalicFont=*-Oblique,
  BoldItalicFont=*-BoldOblique, Scale=0.85, RawFeature={fallback=bpfb}]
\emergencystretch=4em

% Macros for both the web and the pdf version (documentation plan B12).
% The chapters in ../chapters are generated by scripts/blueprint_gen.py; edit this file by hand.

\newtheorem{theorem}{Theorem}
\newtheorem{proposition}[theorem]{Proposition}
\newtheorem{lemma}[theorem]{Lemma}
\newtheorem{corollary}[theorem]{Corollary}

\theoremstyle{definition}
\newtheorem{definition}[theorem]{Definition}

% Macros only for the pdf version; each has a counterpart in macros/web.tex.
%
% Dummy macros for commands that only make sense in the web version (from the leanblueprint
% template).
\newcommand{\lean}[1]{}
\newcommand{\discussion}[1]{}
\newcommand{\leanok}{}
\newcommand{\mathlibok}{}
\newcommand{\notready}{}
% Arguments of \uses and \proves must be real labels: invisible refs (LaTeX warns if a label
% is missing, nothing is shown in the pdf).
\ExplSyntaxOn
\NewDocumentCommand{\uses}{m}
 {\clist_map_inline:nn{#1}{\vphantom{\ref{##1}}}%
  \ignorespaces}
\NewDocumentCommand{\proves}{m}
 {\clist_map_inline:nn{#1}{\vphantom{\ref{##1}}}%
  \ignorespaces}
\ExplSyntaxOff

% Deviation sign of the Markdown blueprints.
\newcommand{\warn}{⚠}

% Unicode characters that no installed font has (LuaLaTeX logs "Missing character: There is no
% ... in font ..."): typeset from glyphs every font has. Add new ones here when the log shows them.
% The web version needs nothing: the browser finds a font.
\usepackage{newunicodechar}
\newunicodechar{ₗ}{\textsubscript{l}}
\newunicodechar{ₙ}{\textsubscript{n}}
\newunicodechar{ᶜ}{\textsuperscript{c}}
\newunicodechar{ȟ}{\v{h}}


\title{Block's theorem in Lean 4: blueprint}
\author{Project \texttt{cohesivemodules}}

\begin{document}
\maketitle
\tableofcontents
% Generated by scripts/blueprint_gen.py; do not edit.
% Generated by scripts/blueprint_gen.py from docs/formalization-summary.tex.
\chapter{Introduction}\label{chap:intro}
\section{The theorem}

\paragraph{Sources.} Block \cite[Thm~4.1.3]{B}; Antweiler \cite[Thm~7.2.20 and Thm~7.4.3]{A};
\cite{CHL} for the comparison theorem in the abstract setting. For the analytic lemmas:
\cite{Hor, GH, DK, BSW}; for twisted sheaves: \cite{Cal, Toda}; for the triangulated structure:
\cite{BK}. The formalization is built on Mathlib \cite{Mathlib}.

\paragraph{Informal statement.} Let $E$ be a finite-dimensional complex normed space and $X$ a
compact Hausdorff complex manifold modelled on $E$. Let $A={\mathcal{A}}_X(X)={\mathcal{A}}^{0,\bullet}(X)$ be the
Dolbeault algebra of $X$ (smooth $(0,q)$-forms with ${\bar\partial}$ and wedge product) and ${\mathcal{O}}_X$ the
sheaf of holomorphic functions. Then the functor
\[ J\colon {\mathrm{Ho}}(\mathcal P_A)\longrightarrow D({\mathcal{O}}_X),\qquad M\longmapsto \alpha(M), \]
from the homotopy category of cohesive modules over $A$ to the derived category of sheaves of
${\mathcal{O}}_X$-modules, is full and faithful, and its essential image is ${D_{\mathrm{perf}}}(X,{\mathcal{O}}_X)$. In other
words, $J$ corestricts to an equivalence ${\mathrm{Ho}}(\mathcal P_A)\simeq{D_{\mathrm{perf}}}(X)$.

\paragraph{History of the statement.} The project first proved the theorem under two analytic
hypotheses on $X$:
\begin{itemize}
\item[(H1)] \emph{Dolbeault--Grothendieck lemma}: ${\bar\partial}$-closed $(0,q{+}1)$-forms are locally
${\bar\partial}$-exact;
\item[(H2)] \emph{local gauge lemma, weak form}: every cohesive module over $A$ is, near each
point, homotopy equivalent to one coming from a bounded complex of finitely generated projective
${\mathcal{O}}$-modules.
\end{itemize}
Both are true on every complex manifold (H1 is classical; H2 is the content of the proof of
\cite[Lemma~4.1.5]{B}), and both are now theorems of the project (\S\ref{chap:Z4}, \S\ref{chap:Z5}). The
conditional statement is kept, as a corollary, under the name \texttt{MainStatementConditional}.

\subsection{The Lean statement}
The statement lives in \texttt{Cohesivemodules/Statements.lean}, which imports only files marked
\texttt{-- FROZEN} (definitions checked in excerpts at a review gate, which may not be changed)
and Mathlib:
\begin{verbatim}
def MainStatement : Prop :=
  ∀ (E : Type) [NormedAddCommGroup E] [NormedSpace ℂ E] [FiniteDimensional ℂ E]
    (X : Type) [TopologicalSpace X] [T2Space X] [CompactSpace X]
    [ChartedSpace E X] [IsManifold 𝓘(ℂ, E) ∞ X],
      letI := HasDerivedCategory.standard (Mod (ringOf (holSheaf E X)))
      (blockFunctor E X).Full ∧ (blockFunctor E X).Faithful ∧
        (blockFunctor E X).essImage = Perf (ringOf (holSheaf E X))
\end{verbatim}
and the theorem, in \texttt{MainH2.lean}:
\begin{verbatim}
theorem main_theorem : Cohesive.MainStatement
-- #print axioms main_theorem : [propext, Classical.choice, Quot.sound]
\end{verbatim}
Reading the statement line by line:
\begin{itemize}
\item \texttt{IsManifold 𝓘(ℂ, E) ∞ X}: $X$ has an atlas with complex-$C^\infty$ (that is,
holomorphic) transition maps. In Mathlib $\infty$ means $C^\infty$; $\top$ would mean
analytic ($C^\omega$). We checked the elaborated statement: the exponent is $\infty$.
\item \texttt{T2Space X}, \texttt{CompactSpace X}: compact Hausdorff; these are implicit in the
sources.
\item \texttt{blockFunctor E X : HoP (dolbeault E X) ⥤ DerivedCategory (Mod (ringOf (holSheaf E X)))}:
the functor $J$ (\S\ref{sec:J}).
\item \texttt{Perf R}: the objects of $D(R)$ isomorphic to $Q(K)$ for a perfect complex $K$
(\S\ref{sec:perf}).
\item \texttt{HasDerivedCategory.standard}: Mathlib's derived category is defined relative to an
instance (a choice of universe and of localization); the statement fixes the standard one,
once, for both \texttt{blockFunctor} and \texttt{Perf}.
\end{itemize}
The conditional form \texttt{MainStatementConditional} (\texttt{StatementsConditional.lean}, frozen
at gate~5) is the same statement with the premises
\texttt{DolbeaultGrothendieck E X → LocalGaugeLemma E X →} after the binders; it is proved in
\texttt{Main.lean} as \texttt{main\_theorem\_conditional}, and \texttt{main\_theorem} applies it to the
proofs \texttt{dolbeaultGrothendieck E X} and \texttt{localGaugeLemma E X}. (Until 6 November 2026
the conditional statement was called \texttt{MainStatement} and the present one
\texttt{MainStatementUnconditional}; the contents did not change.)

\paragraph{Why ``full, faithful, essential image'' and not ``an equivalence onto ${D_{\mathrm{perf}}}$''.}
To write a functor ${\mathrm{Ho}}(\mathcal P_A)\to{D_{\mathrm{perf}}}(X)$ one must already know that $\alpha(M)$ is
perfect, which uses H2 (step C2 below). The functor $J$ is therefore defined into the whole
derived category, without using H1 or H2. The check file \texttt{ChecksMain.lean} proves (K6)
that, for any functor $F$ and any isomorphism-closed property $P$,
\[ F \text{ full} \wedge F\text{ faithful}\wedge \operatorname{essImage}F=P \iff
\exists\, h\colon \forall X,\, P(F X),\ \ (P.\mathrm{lift}\,F\,h) \text{ is an equivalence,} \]
and applies it to \texttt{blockFunctor} and \texttt{Perf} (\texttt{Checks.blockFunctor\_iff}): the
statement is equivalent to ``$J$ corestricts to an equivalence ${\mathrm{Ho}}(\mathcal P_A)\simeq{D_{\mathrm{perf}}}(X)$''.

\paragraph{Sanity checks on the statement.} In complex dimension $0$
($E=\mathrm{Fin}\,0\to\mathbb C$) H1 and H2 are proved directly (K1), independently of their
general proofs, and the point is an example of the theorem (\texttt{blockPerf\_dim\_zero}).
Non-vacuity of every definition (cohesive modules exist, the Dolbeault data satisfy all fields of
the abstract structures, ${\bar\partial}\bar z=d\bar z$ on $\mathbb C$, ${\bar\partial}{\bar\partial}=0$ and Leibniz on
explicit forms on $\mathbb C^2$) is checked in \texttt{Checks*.lean}. Further checks of the
definitions behind the statement: ${\mathrm{Hom}}_{{\mathrm{Ho}}(\mathcal P_A)}({\mathcal{A}},{\mathcal{A}}[n])\cong H^{0,n}_{{\bar\partial}}(X)$
(K4, \S\ref{chap:Z2}), and agreement of our ${\bar\partial}$ with an independent definition by type
projections (K7).

\subsection{The twisted theorem}\label{sec:tw-statement}
For a ${\bar\partial}$-closed $h\in{\mathcal{A}}^{0,2}(X)$, cohesive modules over the \emph{curved} algebra
$({\mathcal{A}}_X(X),{\bar\partial},h)$ have ${\mathbb{X}}\circ{\mathbb{X}}=h\cdot(-)$, so $\alpha(M)$ is no longer a complex. On the
members $U_\alpha$ of a finite cover with $h|_{U_\alpha}={\bar\partial} h^0_\alpha$ the operator
${\mathbb{X}}-h^0_\alpha$ is flat, and the local complexes glue by units $u_{\alpha\beta}$ up to a Čech
2-cocycle ${\check h}$ of units of ${\mathcal{O}}_X$. \cite[Thm~7.4.3]{A} (stated without proof in \cite[Thm~4.2.1]{B})
says that this identifies ${\mathrm{Ho}}(\mathcal P_{(A,{\bar\partial},h)})$ with twisted perfect complexes.
\begin{verbatim}
def MainStatementTwisted : Prop :=
  ∀ (E : Type) [NormedAddCommGroup E] [NormedSpace ℂ E] [FiniteDimensional ℂ E]
    (X : Type) [TopologicalSpace X] [T2Space X] [CompactSpace X]
    [ChartedSpace E X] [IsManifold 𝓘(ℂ, E) ∞ X]
    (h : DolbeaultClosed2 E X) (T : TwistData E X h),
      letI := HasDerivedCategory.standard (TwMod (dolbeaultPreData E X).toSheafData.O T.U T.ch)
      (twistedBlockFunctor h T).Full ∧ (twistedBlockFunctor h T).Faithful ∧
        (twistedBlockFunctor h T).essImage = TwPerf (dolbeaultPreData E X).toSheafData.O T.U T.ch
\end{verbatim}
proved in \texttt{MainTwisted.lean} as \texttt{main\_theorem\_twisted}, with the same three axioms.
Here \texttt{T} is a \emph{twist datum} for $h$ (\S\ref{sec:tw-defs}); one exists for every $h$
(\texttt{exists\_twistData}), so the statement is not vacuous. The check K8
(\texttt{K8\_twisted\_implies\_unconditional}) proves \texttt{MainStatementTwisted → MainStatement}:
at $h=0$ the curved algebra is the flat one by \texttt{rfl}.

%======================================================================
\section{The definitions behind the statements}
The script \texttt{scripts/statement\_defs.py} computes, from the compiled environment, every
project declaration that the statements depend on (types and bodies of definitions; proofs are
not followed since they cannot change the meaning). For \texttt{MainStatement} and
\texttt{MainStatementConditional} the result is 202 declarations in 25 files, frozen at gates 1--5;
with \texttt{MainStatementTwisted} it is 276 declarations in 44 files, all frozen at human review gates,
where they were checked in excerpts (the twisted ones at gate~12). Instances on reached types that are not term dependencies (for
example the triangulated structure on \texttt{HoP}) are listed for information and are not frozen.
This section describes the definitions.

\subsection{Graded-commutative curved dgas (A1)}
A dga is encoded as a $k$-algebra $A$ with a Mathlib \texttt{GradedAlgebra} structure
$\mathcal A\colon\mathbb Z\to\mathrm{Submodule}\,k\,A$ (internal grading) and a class
\texttt{GCDga 𝒜} bundling: ${\mathcal{A}}^q=0$ for $q<0$ and $q>N$, graded commutativity
$ab=(-1)^{pq}ba$, a degree-one derivation $d$ with graded Leibniz rule and $d^2=0$, and a
curvature $c\in{\mathcal{A}}^2$ with $dc=0$ ($d^2=[c,-]$ is automatic, since $c$ is central by graded
commutativity). For Block's theorem $c=0$;
the twisted theorem uses $c=h$. ${\mathcal{A}}^0$ is a commutative ring, and $A$ is an ${\mathcal{A}}^0$-algebra.

\subsection{Cohesive modules and their homotopy category (A3--A5)}
We use \emph{left} modules (Block uses right modules; the translation is the sign convention
$c\leftrightarrow -c$, and a sanity item A13 proves the two definitions dg-equivalent).
\begin{verbatim}
structure QuasiCohesive where
  Q : ℤ → Type                         -- generators Q^p, each an 𝒜⁰-module
  X : Tot 𝒜 Q →ₗ[k] Tot 𝒜 Q          -- 𝕏 on the total module A ⊗_{A⁰} ⨁ Q^p
  X_mem   : ...                        -- 𝕏 has total degree 1
  leibniz : ∀ {p} {a}, a ∈ 𝒜 p → ∀ x,
    X (a • x) = GCDga.d a • x + p.negOnePow • a • X x
  flat    : ∀ x, X (X x) = GCDga.c • x

structure CohesiveModule extends QuasiCohesive 𝒜 where
  bounded    : ∃ a b, ∀ p, (p < a ∨ b < p) → Subsingleton (Q p)
  finite     : ∀ p, Module.Finite (𝒜 0) (Q p)
  projective : ∀ p, Module.Projective (𝒜 0) (Q p)
\end{verbatim}
A cohesive module is thus a bounded graded finitely generated projective ${\mathcal{A}}^0$-module $Q$ with a
$\mathbb Z$-graded superconnection ${\mathbb{X}}$ on $A\otimes_{{\mathcal{A}}^0}Q$ with ${\mathbb{X}}^2=c$ \cite[Def~2.3.2]{B}.

\emph{Hom complexes} (A4): ${\mathrm{Hom}}^n(M_1,M_2)$ consists of $k$-linear maps $\varphi$ of degree $n$
with $\varphi(a\cdot m)=(-1)^{n|a|}a\cdot\varphi(m)$, and
$\delta\varphi={\mathbb{X}}_2\varphi-(-1)^n\varphi{\mathbb{X}}_1$. The sign convention agrees with Mathlib's
\texttt{CochainComplex.HomComplex.δ} (proved). \emph{The homotopy category} \texttt{HoP 𝒜} (A5)
has cohesive modules as objects and $H^0({\mathrm{Hom}}^\bullet(M_1,M_2))$ as morphisms (closed degree-0
maps modulo coboundaries).

\emph{Base change} (A9): for a dga map $f\colon A\to B$, $f^*M$ has generators
$B^0\otimes_{A^0}Q^p$ and ${\mathbb{X}}_B(b\otimes m)=db\otimes m+(-1)^{|b|}b\cdot{\mathbb{X}} m$; it is used to
restrict cohesive modules to open sets. \emph{Modules from complexes} (A10): for $\rho\colon R\to
A^0$ with $d\rho=0$ and a bounded complex $W$ of finitely generated projective $R$-modules,
$A\otimes_R W$ with ${\mathbb{X}}=d\otimes1+1\otimes d_W$ (Koszul sign) is cohesive
(\texttt{CohesiveModule.fromComplex}).

\subsection{\texorpdfstring{Sheaves, perfect complexes, ${D_{\mathrm{perf}}}$ (S1, S9)}{Sheaves, perfect complexes, D\_perf (S1, S9)}}\label{sec:perf}
Sheaves of modules are Mathlib's \texttt{SheafOfModules} over a sheaf of rings
$R$ on \texttt{Opens X}; \texttt{Mod R} abbreviates the category. \texttt{ringOf} turns a sheaf of
commutative rings into a sheaf of rings. Restriction to an open $U$ uses Mathlib's over-site
\texttt{Over U}.
\begin{verbatim}
def IsStrictlyPerfect (P : CochainComplex (Mod R) ℤ) : Prop :=
  (∃ a b : ℤ, P.IsStrictlyGE a ∧ P.IsStrictlyLE b) ∧ ∀ n, IsFinFreeSummand (P.X n)

def IsPerfect (K : CochainComplex (Mod R) ℤ) : Prop :=
  ∀ x : X, ∃ U : Opens X, x ∈ U ∧
    ∃ (L P : CochainComplex (Mod (R.over U)) ℤ) (f : L ⟶ K|U) (g : L ⟶ P),
      QuasiIso f ∧ QuasiIso g ∧ IsStrictlyPerfect P

def Perf (R : RingSheaf X) [HasDerivedCategory.{w} (Mod R)] :
    ObjectProperty (DerivedCategory (Mod R)) :=
  fun D ↦ ∃ K, IsPerfect K ∧ Nonempty (D ≅ DerivedCategory.Q.obj K)
\end{verbatim}
A perfect complex is locally connected by a roof of quasi-isomorphisms to a bounded complex of
direct summands of finite free modules. This is equivalent to the usual
``$Q(K|_U)\cong Q(P)$ in $D(R|_U)$'' but avoids a derived category for every open set.
A sheaf of rings $S$ is \emph{fine} (\texttt{IsFine}) if every finite open cover admits a partition
of unity $\sum\rho_i=1$ by global sections with $\rho_i$ supported in $U_i$ (in a stalk-free
form). On a compact space finite covers suffice.

\subsection{The Dolbeault algebra and holomorphic functions (F1--F12)}
\begin{itemize}
\item \emph{Fibres} (F2): $\Lambda^{0,q}(E;F)$ is the space of continuous $\mathbb C$-alternating
$q$-linear maps on the conjugate space $\bar E$ (\texttt{Conj E}, a type synonym with the conjugate
action). A $\mathbb C$-linear map on $\bar E$ is an antilinear map on $E$, so these are exactly
the $(0,q)$-covectors.
\item \emph{${\bar\partial}$ on open subsets of $E$} (F3): $D''\omega(x)=\tfrac12(D\omega(x)+i\,D\omega(x)\circ J)$,
the antilinear part of the real derivative, and ${\bar\partial}\omega(x)$ is its alternation (Mathlib's
\texttt{alternatizeUncurryFin}). ${\bar\partial}{\bar\partial}=0$ (F4) by copying Mathlib's proof of $d\circ d=0$
with $D''$ for $D$.
\item \emph{Product} (F5): the total fibre $\Lambda^{0,\bullet}=\prod_{q\le n}\Lambda^{0,q}$ gets
its algebra structure transported from Mathlib's exterior algebra of $\mathrm{Dual}_{\mathbb C}\bar
E$ (associativity and graded commutativity for free); Leibniz for ${\bar\partial}$ (F6).
\item \emph{Holomorphic pullback} (F7) commutes with ${\bar\partial}$; this makes the chart-wise definition
independent of charts.
\item \emph{Sheaves on $X$} (F9--F10): a section of ${\mathcal{A}}_X$ over $U$ is a function into the total
fibre, read in charts, whose chart expressions are smooth (a Mathlib \texttt{LocalPredicate}).
${\mathcal{A}}_X(U)$ is a bounded graded-commutative dga with ${\bar\partial}$ (\texttt{dolbeaultGrade E U}); the
Dolbeault algebra is \texttt{dolbeault E X := dolbeaultGrade E ⊤}. The degree-$q$ component sheaves
are subtypes of the total section algebra, so that ${\mathcal{A}}^0$ is literally the dga's degree-0 part.
\item \emph{Holomorphic functions} (F12): \texttt{holSheaf E X := smoothSheafCommRing 𝓘(ℂ, E) 𝓘(ℂ, ℂ) X ℂ},
complex-$C^\infty$ functions, which are the holomorphic ones; $\rho\colon{\mathcal{O}}_X\to{\mathcal{A}}^0$ is the
inclusion.
\end{itemize}

\subsection{The hypotheses H1 and H2}\label{sec:hyp}
They are \texttt{Prop}-valued definitions in \texttt{Hypotheses.lean} (frozen at gate~3). They occur
in \texttt{MainStatementConditional} and, as fields, in the abstract resolving data below; they do
not occur in \texttt{MainStatement}.
\begin{verbatim}
def DolbeaultGrothendieck : Prop :=
  ∀ (U : Opens X) (q : ℕ) (x : X), x ∈ U →
    ∀ α ∈ dolbeaultGrade E U ((q : ℤ) + 1), GCDga.d α = 0 →
      ∃ (U' : Opens X) (h : U' ≤ U), x ∈ U' ∧
        ∃ η ∈ dolbeaultGrade E U' (q : ℤ), GCDga.d η = (dolbeaultRes E h) α

def LocalGaugeLemma : Prop :=
  ∀ (M : CohesiveModule (dolbeault E X)) (x : X),
    ∃ (U' : Opens X) (_ : x ∈ U')
      (W : CochainComplex (ModuleCat (O_X(U'))) ℤ) (hW : IsBoundedFGProj _ W),
      Nonempty (HoP.mk (M.baseChange (dolbeaultRes E (le_top : U' ≤ ⊤))) ≅
                HoP.mk (CohesiveModule.fromComplex (ρ := dolbeaultRhoU E U') W ...))
\end{verbatim}
H1 together with F12 says exactly that $0\to{\mathcal{O}}_X\to{\mathcal{A}}^{0,0}\xrightarrow{{\bar\partial}}{\mathcal{A}}^{0,1}\to\cdots$ is
a resolution. H2 says: near each point, the restriction $M|_{U'}$ is isomorphic in
${\mathrm{Ho}}(\mathcal P_{{\mathcal{A}}(U')})$ to ${\mathcal{A}}(U')\otimes_{{\mathcal{O}}(U')}W$. It is stated for sections over
small opens, only for restrictions of global modules, which is all the proof uses (\cite{A}'s
Assumption~6 is a sheaf version on every open).

\subsection{\texorpdfstring{The functor $J$ (C0, C1)}{The functor J (C0, C1)}}\label{sec:J}
The comparison is proved for abstract \emph{resolving data}, then applied to the Dolbeault data.
\texttt{PreResolvingData X} bundles, for a topological space $X$: a presheaf of bounded
graded-commutative dgas $U\mapsto{\mathcal{A}}(U)$ that is a sheaf (gluing, locality of degree), with flat
curvature; a sheaf of commutative rings ${\mathcal{O}}$ and ring maps $\rho_U\colon{\mathcal{O}}(U)\to{\mathcal{A}}^0(U)$,
natural, with $d\circ\rho=0$; ${\mathcal{A}}^0$ fine; each ${\mathcal{A}}^q$ locally free over ${\mathcal{A}}^0$ (local frames).
\texttt{ResolvingData} adds the resolution condition ($\rho$ injective, exactness in degree $0$ and
in degrees $q+1$) and the local lifting field \texttt{LL} (the abstract form of H2).

For pre-resolving data $P$ and a cohesive module $M$ over $A={\mathcal{A}}(X)$, the complex
$\alpha(M)$ (Lean: \texttt{P.sheafOf M}) has
\[ \alpha(M)^n(V) = \bigl({\mathcal{A}}(V)\otimes_{{\mathcal{A}}^0(V)}\textstyle\bigoplus_p {\mathcal{A}}^0(V)\otimes_{{\mathcal{A}}^0(X)}Q^p\bigr)^n
\;\cong\; \textstyle\bigoplus_{p+q=n}{\mathcal{A}}^q(V)\otimes_{{\mathcal{A}}^0(X)}Q^p, \]
the degree-$n$ part of the total module of the base change $M|_V$, with the differential ${\mathbb{X}}$ of
$M|_V$; ${\mathcal{O}}$ acts through $\rho_V$; restriction maps are the base-change maps. (The first form
reuses the base change A9 and its functoriality; it is canonically isomorphic to Block's
formula.) It is a sheaf (C1.a: $Q^p$ is finitely generated projective and ${\mathcal{A}}^q$ is a summand of
a finite free ${\mathcal{A}}^0$-module). A closed degree-0 map gives a chain map, a coboundary a
null-homotopic one (C1.b), so
\[ J := ({\mathrm{Ho}}(\mathcal P_A)\to K(\mathrm{Mod}\,{\mathcal{O}})\to D(\mathrm{Mod}\,{\mathcal{O}})),\qquad
\texttt{P.J := P.JK ⋙ DerivedCategory.Qh}. \]
Finally (\texttt{BlockFunctorDef.lean}):
\begin{verbatim}
def dolbeaultPreData : PreResolvingData X   -- the Dolbeault data; every field a theorem
def blockFunctor [h : HasDerivedCategory.{w} (Mod (ringOf (holSheaf E X)))] :
    HoP (dolbeault E X) ⥤ DerivedCategory (Mod (ringOf (holSheaf E X))) :=
  @PreResolvingData.J X _ (dolbeaultPreData E X) h
\end{verbatim}
Every field of \texttt{dolbeaultPreData} is a theorem of layer F: gluing and locality of forms,
$\rho$ natural with ${\bar\partial}\rho=0$, ${\mathcal{A}}^0$ fine (smooth partitions of unity, needs compact
Hausdorff), ${\mathcal{A}}^{0,q}$ locally free of rank $\binom nq$. No analytic hypothesis is used:
\texttt{blockFunctor} is defined without H1, H2.

\subsection{Curved Dolbeault algebras, twist data, twisted sheaves (Z3)}\label{sec:tw-defs}
\begin{itemize}
\item \texttt{DolbeaultClosed2 E X}: global forms $h$ of degree $2$ with ${\bar\partial} h=0$;
\texttt{curvedDolbeault h} is $({\mathcal{A}}_X(X),{\bar\partial},h)$, the Dolbeault algebra with curvature replaced
by $h$ (\texttt{GCDga.withC}, on the same carrier).
\item \texttt{TwistData E X h}: a finite open cover $U=(U_\alpha)$, primitives
$h^0_\alpha\in{\mathcal{A}}^{0,1}(U_\alpha)$ with ${\bar\partial} h^0_\alpha=h|_{U_\alpha}$, units
$u_{\alpha\beta}\in{\mathcal{A}}^0(U_\alpha\cap U_\beta)^\times$ with
${\bar\partial} u_{\alpha\beta}=(h^0_\beta-h^0_\alpha)u_{\alpha\beta}$, $u_{\alpha\alpha}=1$,
$u_{\beta\alpha}=u_{\alpha\beta}^{-1}$, and a normalized Čech 2-cocycle ${\check h}$ of units of
${\mathcal{O}}_X$ with $\rho({\check h}_{\alpha\beta\gamma})=u_{\beta\gamma}u_{\alpha\beta}u_{\alpha\gamma}^{-1}$.
\item \texttt{TwMod O U ȟ}: ${\check h}$-twisted sheaves of modules, i.e.\ sheaves $M_\alpha$ on the
over-sites \texttt{Over U\textsubscript{α}} with gluing maps $\varphi_{\alpha\beta}$ on sections over every
$W\le U_\alpha\cap U_\beta$, $\varphi_{\alpha\alpha}=\mathrm{id}$ and
$\varphi_{\beta\gamma}\circ\varphi_{\alpha\beta}={\check h}_{\alpha\beta\gamma}\varphi_{\alpha\gamma}$. It is an
abelian category (its \texttt{Abelian} instance is frozen, since the derived category depends on it)
with enough injectives.
\item \texttt{TwPerf}: objects of $D(\mathrm{TwMod})$ isomorphic to $Q(K)$ with every component
$K_\alpha$ perfect on $U_\alpha$.
\item \texttt{twistedBlockFunctor h T}: $M$ goes to the twisted complex with components
$\alpha(M)^n|_{U_\alpha}$, differential ${\mathbb{X}}-h^0_\alpha$ and gluing $u_{\alpha\beta}\cdot$.
\end{itemize}

%======================================================================

% Generated by scripts/blueprint_gen.py; do not edit.
\chapter{Cohesive modules over a dga}\label{chap:A}
\textbf{Setting for all of §2.} k a commutative ring. A a ℤ-graded k-algebra with A\textsuperscript{q} = 0 for q < 0 and q > N, graded-commutative (ab = (−1)\textsuperscript{pq} ba for a ∈ A\textsuperscript{p}, b ∈ A\textsuperscript{q}), with a degree-one k-linear derivation d satisfying the graded Leibniz rule and d² = 0. Write A\textsuperscript{0} for the (commutative) degree-0 part. Curvature c ∈ A\textsuperscript{2} with dc = 0 is carried in the definitions because it costs nothing; \textbf{this project only ever uses c = 0} (c ≠ 0 is Z3).

\begin{definition}[A1 — Bounded graded-commutative dga]
\label{A1}
\lean{Cohesive.GCDga}
\leanok
Let k be a commutative ring. A \emph{bounded graded-commutative curved dga} over k is a ℤ-graded k-algebra A = ⨁\textsubscript{q} A\textsuperscript{q} together with

\begin{itemize}
\item a top degree N with A\textsuperscript{q} = 0 for q < 0 and for q > N;
\item graded commutativity: ab = (−1)\textsuperscript{pq} ba for a ∈ A\textsuperscript{p}, b ∈ A\textsuperscript{q};
\item a k-linear differential d of degree +1 with d(ab) = da·b + (−1)\textsuperscript{p} a·db for a ∈ A\textsuperscript{p}, and d ∘ d = 0;
\item a curvature c ∈ A\textsuperscript{2} with dc = 0.
\end{itemize}
Deviation: Block allows non-commutative A with d² = [c, −]; for graded-commutative A this is d² = 0. Boundedness is added (it holds for the Dolbeault algebra).
\par\smallskip{\footnotesize Source: blueprint.md §2.}
\end{definition}
\paragraph{Blueprint notes.}
Definition of the setting above. Encoding: internal ℤ-grading, see conventions §5. Lean: \texttt{GradedAlgebra} ✓; no graded-commutativity class in Mathlib (✗ as a class), so state it as a field. \texttt{class GCDga} in \texttt{Algebra/Dga.lean}, with A\textsuperscript{0} a commutative ring and A an A\textsuperscript{0}-algebra built from graded commutativity. \warn{} ℤ-indexed (zero in negative degrees) rather than ℕ-indexed.

\begin{theorem}[A2 — Connections]
\label{A2}
\lean{Cohesive.Connection, Cohesive.Connection.exists_connection, Cohesive.Connection.diff}
\leanok
Let A be a bounded graded-commutative curved dga over a commutative ring k, and P an A\textsuperscript{0}-module. A \emph{connection} on P is an additive map ∇ : P → A\textsuperscript{1} ⊗\textsubscript{A\textasciicircum{}0} P with ∇(fe) = df ⊗ e + f·∇e for f ∈ A\textsuperscript{0} (it is then automatically k-linear).

\begin{itemize}
\item Every finitely generated projective A\textsuperscript{0}-module has a connection.
\item The difference ∇₁ − ∇₂ of two connections on P is A\textsuperscript{0}-linear.
\end{itemize}
Deviation: additivity is assumed and k-linearity derived; left modules.
\par\smallskip{\footnotesize Source: blueprint.md §2.}
\end{theorem}
\begin{proof}
\leanok
\uses{A1}
See the Lean proof and the cited items.
\paragraph{Blueprint notes.}
A connection on an A\textsuperscript{0}-module P is a k-linear ∇ : P → A\textsuperscript{1} ⊗\textsubscript{A\textasciicircum{}0} P with ∇(fe) = df ⊗ e + f∇e. Extension to A ⊗\textsubscript{A\textasciicircum{}0} P by ∇(a ⊗ e) = da ⊗ e + (−1)\textsuperscript{|a|} a·∇e.

\begin{itemize}
\item \textbf{A2.a} Every finitely generated projective A\textsuperscript{0}-module has a connection. Proof: P ⊕ P' ≅ (A\textsuperscript{0})\textsuperscript{m}; take ∇ = (1 ⊗ π) ∘ d\textsuperscript{⊕m} ∘ ι. [B, proof of Lemma 3.2.8].
\item \textbf{A2.b} The difference of two connections is A\textsuperscript{0}-linear.
\end{itemize}
\end{proof}

\begin{theorem}[A3 — Cohesive and quasi-cohesive modules]
\label{A3}
\lean{Cohesive.QuasiCohesive, Cohesive.CohesiveModule, Cohesive.QuasiCohesive.ofGenerators, Cohesive.QuasiCohesive.X_eq_of_generators, Cohesive.QuasiCohesive.comp_one_leibniz, Cohesive.QuasiCohesive.X0}
\leanok
Let A be a bounded graded-commutative curved dga over k with curvature c. A \emph{quasi-cohesive module} M = (Q, 𝕏) is a family of A\textsuperscript{0}-modules Q\textsuperscript{p} (p ∈ ℤ) together with a k-linear operator 𝕏 of total degree 1 on the total module Tot Q = A ⊗\textsubscript{A\textasciicircum{}0} ⨁\textsubscript{p} Q\textsuperscript{p} (A acting on the left factor) such that 𝕏(a·x) = da·x + (−1)\textsuperscript{|a|} a·𝕏x for homogeneous a, and 𝕏 ∘ 𝕏 = c·(−). It is \emph{cohesive} if Q\textsuperscript{p} = 0 outside a bounded range of p and every Q\textsuperscript{p} is finitely generated projective over A\textsuperscript{0}; P\textsubscript{A} denotes the cohesive modules.

\begin{itemize}
\item 𝕏 is determined by its values on the generators 1 ⊗ e, and conversely an additive map θ on the generators of total degree 1 with θ(fe) = df ⊗ e + f·θ(e) that is flat on generators extends to a quasi-cohesive module with 𝕏(a ⊗ e) = da ⊗ e + (−1)\textsuperscript{|a|} a·θ(e).
\item Writing 𝕏\textsuperscript{j} for the form-degree-j part of 𝕏 on generators: 𝕏\textsuperscript{1}(fe) = df ⊗ e + f·𝕏\textsuperscript{1}(e) (a connection, no sign); 𝕏\textsuperscript{0} is an A\textsuperscript{0}-linear map of degree +1 on ⨁\textsubscript{p} Q\textsuperscript{p} with 𝕏\textsuperscript{0} ∘ 𝕏\textsuperscript{0} = 0.
\end{itemize}
Deviation: left modules; the curvature is c = −c\textsubscript{Block}.
\par\smallskip{\footnotesize Source: blueprint.md §2.}
\end{theorem}
\begin{proof}
\leanok
\uses{A1}
See the Lean proof and the cited items.
\paragraph{Blueprint notes.}
A \emph{quasi-cohesive module} is a ℤ-graded A\textsuperscript{0}-module Q together with a k-linear operator 𝕏 of total degree 1 on A ⊗\textsubscript{A\textasciicircum{}0} Q (left A-module, A acting on the left factor) such that

\begin{itemize}
\item Leibniz: 𝕏(a·m) = da·m + (−1)\textsuperscript{|a|} a·𝕏m, and
\item flatness: 𝕏 ∘ 𝕏 = c·(−) (left multiplication; c = 0 here).
\end{itemize}
It is \emph{cohesive} if moreover Q\textsuperscript{p} ≠ 0 for only finitely many p and each Q\textsuperscript{p} is finitely generated projective. P\textsubscript{A} denotes the cohesive ones. [B, Def 2.3.2] \warn{} left modules, see conventions §3 for the translation to Block's right modules.

\begin{itemize}
\item \textbf{A3.a Components.} Restriction to Q determines 𝕏, with components 𝕏\textsuperscript{j} : Q\textsuperscript{p} → A\textsuperscript{j} ⊗ Q\textsuperscript{p+1−j}; 𝕏\textsuperscript{0} is an A\textsuperscript{0}-linear differential, 𝕏\textsuperscript{1} is a connection on each Q\textsuperscript{p} (no sign in the left convention), 𝕏\textsuperscript{j} (j ≥ 2) are A\textsuperscript{0}-linear. [B, after Def 2.3.1], [A, eq. (2)].
\end{itemize}
A3.b in degree 0 (\texttt{X0\_X0}) is part of this item; the other degrees are A3.b′ (not formalized, not needed).
\end{proof}

\begin{theorem}[A3.b′ — Flatness in all degrees]
\label{A3.bprime}
\uses{A3}
Let M = (Q, 𝕏) be a quasi-cohesive module over a bounded graded-commutative curved dga A with curvature c ∈ A\textsuperscript{2}, and 𝕏\textsuperscript{j} : Q\textsuperscript{p} → A\textsuperscript{j} ⊗\textsubscript{A\textasciicircum{}0} Q\textsuperscript{p+1−j} the components of 𝕏 on generators, extended to Tot Q by the Leibniz rule. Then flatness 𝕏 ∘ 𝕏 = c·(−) is equivalent to the equations Σ\textsubscript{i+j=n} 𝕏\textsuperscript{i} ∘ 𝕏\textsuperscript{j} = δ\textsubscript{n,2}·c·(−) for all n ≥ 0, one in each form degree n (in degree 0: 𝕏\textsuperscript{0} ∘ 𝕏\textsuperscript{0} = 0). Not formalized; not used by any other item.
\par\smallskip{\footnotesize Source: blueprint.md §2.}
\end{theorem}
\begin{proof}
See the Lean proof and the cited items.
\paragraph{Blueprint notes.}
\textbf{A3.b} Flatness unfolds degree by degree into the equations of [A, eq. (2)].
\end{proof}

\begin{definition}[A4 — Hom complexes]
\label{A4}
\lean{Cohesive.Hom, Cohesive.δ, Cohesive.homEquivGen}
\leanok
\uses{A3}
Let A be a bounded graded-commutative curved dga over k and M₁ = (Q₁, 𝕏₁), M₂ = (Q₂, 𝕏₂) quasi-cohesive modules. Hom\textsuperscript{n}(M₁, M₂) is the k-module of k-linear maps φ : Tot Q₁ → Tot Q₂ raising total degree by n with φ(a·x) = (−1)\textsuperscript{n|a|} a·φ(x) for homogeneous a, with differential δφ = 𝕏₂ ∘ φ − (−1)\textsuperscript{n} φ ∘ 𝕏₁.

\begin{itemize}
\item δ maps Hom\textsuperscript{n} to Hom\textsuperscript{n+1} and δ ∘ δ = 0 (also with curvature, since |c| = 2).
\item Composition maps Hom\textsuperscript{m} × Hom\textsuperscript{n} to Hom\textsuperscript{m+n}, and δ(ψ ∘ φ) = δψ ∘ φ + (−1)\textsuperscript{m} ψ ∘ δφ for ψ of degree m; the identity lies in Hom\textsuperscript{0} and is closed.
\item δ agrees with Mathlib's \texttt{CochainComplex.HomComplex.δ}: δφ = 𝕏₂ ∘ φ + (−1)\textsuperscript{n+1} φ ∘ 𝕏₁.
\item Restriction to generators is a k-linear isomorphism from Hom\textsuperscript{n}(M₁, M₂) to the A\textsuperscript{0}-linear maps ψ : ⨁\textsubscript{p} Q₁\textsuperscript{p} → Tot Q₂ with ψ(Q₁\textsuperscript{p}) in total degree p + n; the inverse is ψ̂(a ⊗ e) = (−1)\textsuperscript{n|a|} a·ψ(e).
\end{itemize}
Deviation: M₁ may be any quasi-cohesive module (the sources take M₁ cohesive).
\par\smallskip{\footnotesize Source: blueprint.md §2.}
\end{definition}
\paragraph{Blueprint notes.}
For M₁ cohesive and M₂ quasi-cohesive: Hom\textsuperscript{n}(M₁, M₂) = k-linear maps φ of degree n with φ(a·m) = (−1)\textsuperscript{n|a|} a·φ(m); differential δφ = 𝕏₂∘φ − (−1)\textsuperscript{n} φ∘𝕏₁.

\begin{itemize}
\item \textbf{A4.a} δ∘δ = 0, including with curvature (|c| even). [B, Lemma 2.3.4]
\item \textbf{A4.b} Composition is associative, unital and satisfies the Leibniz rule for δ.
\item \textbf{A4.c} Restriction to generators gives Hom\textsuperscript{n}(M₁, M₂) ≅ ⊕\textsubscript{p} Hom\textsubscript{A\textasciicircum{}0}(Q₁\textsuperscript{p}, (A ⊗ Q₂)\textsuperscript{p+n}).
\item \textbf{A4.d} The sign convention agrees with Mathlib's \texttt{CochainComplex.HomComplex.δ} ✓ (δz = z ≫ d + (−1)\textsuperscript{n+1} d ≫ z), conventions §2.
\end{itemize}

\begin{definition}[A5 — Homotopy category Ho]
\label{A5}
\lean{Cohesive.HoP}
\leanok
\uses{A4}
Let A be a bounded graded-commutative curved dga over k. The homotopy category Ho(P\textsubscript{A}) has the cohesive modules as objects and, as morphisms M → N, the degree-0 cohomology H\textsuperscript{0} = Z\textsuperscript{0}/B\textsuperscript{0} of the Hom complex, where Z\textsuperscript{0} = \{φ ∈ Hom\textsuperscript{0}(M, N) : δφ = 0\} and B\textsuperscript{0} = δ(Hom\textsuperscript{−1}(M, N)); composition is induced by composition of maps. Ho(P\textsubscript{A}) is a preadditive, k-linear category.
\par\smallskip{\footnotesize Source: blueprint.md §2.}
\end{definition}
\paragraph{Blueprint notes.}
Objects: cohesive modules. Morphisms: H\textsuperscript{0} of the Hom complex. A category (from A4.b), k-linear. Homotopy equivalences = isomorphisms in Ho(P\textsubscript{A}). [B, §2.4]

\begin{theorem}[A6 — Form-degree filtration]
\label{A6}
\lean{Cohesive.filTot, Cohesive.filHom, Cohesive.memFilHom_iff, Cohesive.δ_mem_filHom}
\leanok
Let A be a bounded graded-commutative curved dga over k with top degree N, and M₁, M₂ quasi-cohesive modules. F\textsuperscript{b} Tot Q is the k-span of the a ⊗ e with a ∈ A\textsuperscript{q}, q ≥ b, and F\textsuperscript{b} Hom\textsuperscript{n}(M₁, M₂) consists of the φ ∈ Hom\textsuperscript{n} with φ(F\textsuperscript{c} Tot Q₁) ⊆ F\textsuperscript{c+b} Tot Q₂ for all c.

\begin{itemize}
\item φ ∈ Hom\textsuperscript{n} lies in F\textsuperscript{b} iff φ(1 ⊗ e) ∈ F\textsuperscript{b} Tot Q₂ for every generator e, i.e. the components of φ into A\textsuperscript{b′} ⊗ Q₂ vanish for b′ < b.
\item δ maps F\textsuperscript{b} Hom\textsuperscript{n} into F\textsuperscript{b} Hom\textsuperscript{n+1}.
\item The filtration is finite: F\textsuperscript{0} Hom\textsuperscript{n} = Hom\textsuperscript{n} and F\textsuperscript{b} Hom\textsuperscript{n} = 0 for b > N.
\end{itemize}
\par\smallskip{\footnotesize Source: blueprint.md §2.}
\end{theorem}
\begin{proof}
\leanok
\uses{A4}
See the Lean proof and the cited items.
\paragraph{Blueprint notes.}
F\textsuperscript{b} Hom\textsuperscript{n} = maps whose components into A\textsuperscript{b'} ⊗ − vanish for b' < b. δ preserves it, the filtration is finite (0 ≤ b ≤ N), and gr\textsuperscript{b} ≅ Hom\textsubscript{A\textasciicircum{}0}(Q₁, A\textsuperscript{b} ⊗ Q₂) with differential induced by 𝕏₁\textsuperscript{0}, 𝕏₂\textsuperscript{0}. [B, Prop 2.5.1], [A, Def 3.3.2]
\end{proof}

\begin{theorem}[A7 — Finite filtration comparison]
\label{A7}
\lean{Cohesive.quasiIso_of_finite_filtration}
\leanok
Let C be an abelian category and N ∈ ℕ. Suppose given, for b = 0, …, N, maps f\textsubscript{b} : F\textsuperscript{b} K → F\textsuperscript{b} L of ℤ-indexed cochain complexes in C, and for b < N maps of short exact sequences 0 → F\textsuperscript{b+1} → F\textsuperscript{b} → gr\textsuperscript{b} → 0 (of the K-side into the L-side) whose first two components are f\textsubscript{b+1} and f\textsubscript{b}. If every induced map gr\textsuperscript{b} K → gr\textsuperscript{b} L (b < N) is a quasi-isomorphism and f\textsubscript{N} is a quasi-isomorphism (e.g. F\textsuperscript{N} = 0), then f\textsubscript{0} is a quasi-isomorphism.

Deviation: replaces the spectral sequences of the sources by the homology long exact sequence.
\par\smallskip{\footnotesize Source: blueprint.md §2.}
\end{theorem}
\begin{proof}
\leanok
See the Lean proof and the cited items.
\paragraph{Blueprint notes.}
A map of finitely filtered complexes inducing quasi-isomorphisms on all gr\textsuperscript{b} is a quasi-isomorphism. Proof: induction on the filtration length with the long exact sequence. Replaces the spectral sequences in [B] and [A]. Lean: homology long exact sequence ✓ (\texttt{HomologicalComplex.HomologySequence}).
\end{proof}

\begin{theorem}[A8 — Homotopy criterion]
\label{A8}
\uses{A7, F15, Z2.g}
Let A be a bounded graded-commutative dga over k and φ : M₁ → M₂ a closed degree-0 map of cohesive modules. Then φ is an isomorphism in Ho(P\textsubscript{A}) iff its (0,0)-component φ\textsuperscript{0} : (Q₁, 𝕏₁\textsuperscript{0}) → (Q₂, 𝕏₂\textsuperscript{0}) is a homotopy equivalence of complexes of A\textsuperscript{0}-modules (equivalently, since these are bounded complexes of projectives, a quasi-isomorphism). Not formalized; not used by any other item.
\par\smallskip{\footnotesize Source: blueprint.md §2.}
\end{theorem}
\begin{proof}
See the Lean proof and the cited items.
\paragraph{Blueprint notes.}
Not on the critical path: no step of C or F uses it. Schedule it after F15, or with Z2. A closed degree-0 φ : M₁ → M₂ between cohesive modules is an isomorphism in Ho(P\textsubscript{A}) iff φ\textsuperscript{0} : (Q₁, 𝕏₁\textsuperscript{0}) → (Q₂, 𝕏₂\textsuperscript{0}) is a homotopy equivalence of complexes of A\textsuperscript{0}-modules. [A, Lemma 3.3.5], [B, Prop 2.5.2]

\begin{itemize}
\item Proof (⇐): for every cohesive M, φ\textsubscript{\emph{} : Hom(M, M₁) → Hom(M, M₂) induces on gr\textsuperscript{b} the map Hom\textsubscript{A\textasciicircum{}0}(Q, A\textsuperscript{b} ⊗ φ\textsuperscript{0}), a homotopy equivalence (any additive functor preserves them); by A7 φ\textsubscript{}} is a quasi-isomorphism; Yoneda as in [A, Lemma 3.3.5].
\item Proof (⇒): take the (0,0)-components of the homotopies.
\item \textbf{A8.a} For bounded complexes of projective A\textsuperscript{0}-modules, homotopy equivalence ⇔ quasi-isomorphism ([B]'s form). Lean: \texttt{CochainComplex.isKProjective\_of\_projective} ✓.
\end{itemize}
\end{proof}

\begin{theorem}[A9 — Base change along dga maps]
\label{A9}
\lean{Cohesive.DgaHom, Cohesive.CohesiveModule.baseChange, Cohesive.bcHoFunctor, Cohesive.bcCompNatIso, Cohesive.bcHomL, Cohesive.hom_eq_of_bcF, Cohesive.DgaHom.comp, Cohesive.bcCompTot}
\leanok
Let A, B be bounded graded-commutative curved dgas over k. A \emph{dga map} f : A → B is a k-algebra map preserving degrees, the differential and the curvature (f(c\textsubscript{A}) = c\textsubscript{B}). For a quasi-cohesive A-module M = (Q, 𝕏), the base change f\textsuperscript{*M} is the quasi-cohesive B-module with generators B\textsuperscript{0} ⊗\textsubscript{A\textasciicircum{}0} Q\textsuperscript{p} and 𝕏\textsubscript{B}(b ⊗ m) = db ⊗ m + (−1)\textsuperscript{|b|} b·𝕏m.

\begin{itemize}
\item If M is cohesive, so is f\textsuperscript{*M}.
\item φ ∈ Hom\textsuperscript{n}(M₁, M₂) gives f\textsuperscript{\emph{φ} ∈ Hom\textsuperscript{n}(f\textsuperscript{*M₁}, f\textsuperscript{*M₂}), (f\textsuperscript{*φ})(b ⊗ m) = (−1)\textsuperscript{n|b|} b·φ(m), with f\textsuperscript{}}(δφ) = δ(f\textsuperscript{*φ}); this gives a functor Ho(P\textsubscript{A}) ⥤ Ho(P\textsubscript{B}).
\item Maps of degree n out of f\textsuperscript{*M₁} are determined by their values on the image of Tot Q₁ → Tot(f\textsuperscript{*Q₁}), a ⊗ e ↦ f(a) ⊗ (1 ⊗ e).
\item For a second dga map g : B → C, γ ⊗ q ↦ γ ⊗ (1 ⊗ q) is a closed degree-0 isomorphism (g ∘ f)\textsuperscript{\emph{M} ≅ g\textsuperscript{}}(f\textsuperscript{\emph{M}), natural on Hom complexes in every degree; it gives a natural isomorphism of functors (g ∘ f)\textsuperscript{}} ≅ g\textsuperscript{\emph{} ∘ f\textsuperscript{}} on Ho(P\textsubscript{A}).
\end{itemize}
Deviation: left modules; Block's ω is 0.
\par\smallskip{\footnotesize Source: blueprint.md §2.}
\end{theorem}
\begin{proof}
\leanok
\uses{A5}
See the Lean proof and the cited items.
\paragraph{Blueprint notes.}
For a dga map f : A → B, M = (Q, 𝕏) ↦ f\textsuperscript{*M} with generators B\textsuperscript{0} ⊗\textsubscript{A\textasciicircum{}0} Q\textsuperscript{p} (Block's E\textsuperscript{•} ⊗\textsubscript{𝒜\_1} 𝒜\textsubscript{2}, [B, §2.7]) and 𝕏\textsubscript{B}(b ⊗ m) = db ⊗ m + (−1)\textsuperscript{|b|} b·𝕏m; its total module B ⊗\textsubscript{B\textasciicircum{}0} ⨁\textsubscript{p} (B\textsuperscript{0} ⊗\textsubscript{A\textasciicircum{}0} Q\textsuperscript{p}) is canonically B ⊗\textsubscript{A} M. Cohesive goes to cohesive (B\textsuperscript{0} ⊗\textsubscript{A\textasciicircum{}0} Q\textsuperscript{p} is finitely generated projective); a dg functor; preserves homotopy equivalences. Used for restriction of cohesive modules to open subsets. [B, §2.7 with ω = 0]

\begin{itemize}
\item \textbf{A9.a} Base change along a composite B → C is naturally isomorphic to the composite of base changes, compatibly with Hom complexes (used when restricting from X to U′ to V ⊆ U′).
\end{itemize}
\end{proof}

\begin{theorem}[A10 — Cohesive modules from complexes]
\label{A10}
\lean{Cohesive.QuasiCohesive.fromComplex, Cohesive.IsBoundedFGProj, Cohesive.fcResP_δ, Cohesive.homRδ, Cohesive.homRδ_mem, Cohesive.fc_hom_ext}
\leanok
Let A be a bounded graded-commutative dga over k with curvature c = 0, ρ : R → A\textsuperscript{0} a map of commutative rings with d ∘ ρ = 0, and (W, d\textsubscript{W}) a cochain complex of R-modules. Then A ⊗\textsubscript{R} W, with generators A\textsuperscript{0} ⊗\textsubscript{R} W\textsuperscript{p} and 𝕏(a ⊗ w) = da ⊗ w + (−1)\textsuperscript{|a|} a ⊗ d\textsubscript{W} w, is a quasi-cohesive module; it is cohesive if W is bounded with finitely generated projective terms. For every quasi-cohesive N, restriction along w ↦ 1 ⊗ w is an additive bijection from Hom\textsuperscript{n}(A ⊗\textsubscript{R} W, N) to the families of ρ-semilinear maps f\textsubscript{p} : W\textsuperscript{p} → Tot N landing in total degree p + n; it carries δ to (δf)\textsubscript{p} = 𝕏\textsubscript{N} ∘ f\textsubscript{p} − (−1)\textsuperscript{n} f\textsubscript{p+1} ∘ d\textsubscript{W}.

Deviation: requires c = 0 (otherwise 𝕏 ∘ 𝕏 = 0 ≠ c·(−)).
\par\smallskip{\footnotesize Source: blueprint.md §2.}
\end{theorem}
\begin{proof}
\leanok
\uses{A9}
See the Lean proof and the cited items.
\paragraph{Blueprint notes.}
Let ρ : R → A\textsuperscript{0} be a map of commutative rings with d∘ρ = 0, and (W, d\textsubscript{W}) a bounded complex of finitely generated projective R-modules. Then A ⊗\textsubscript{R} W with 𝕏 = d ⊗ 1 + 1 ⊗ d\textsubscript{W} is cohesive. (v1.10) With the Koszul sign of conventions §2: 𝕏(a ⊗ w) = da ⊗ w + (−1)\textsuperscript{|a|} a ⊗ d\textsubscript{W} w; on generators 1 ⊗ w no sign appears, and this is what \texttt{fromComplex} implements.

\begin{itemize}
\item \textbf{A10.a Adjunction.} Hom\textsuperscript{•\_}\{P\textsubscript{A}\}(A ⊗\textsubscript{R} W, M) ≅ Hom\textsuperscript{•\_R}(W, M) as complexes, via precomposition with w ↦ 1 ⊗ w. Used in C4.
\end{itemize}
\end{proof}

\begin{theorem}[A11 — Lifting theorem]
\label{A11}
\lean{Cohesive.lifting, Cohesive.CohesiveModule.ofGenData, Cohesive.lift_core, Cohesive.liftStage_one, Cohesive.liftStage_succ, Cohesive.ConeLift}
\leanok
\uses{A2}
Let A be a bounded graded-commutative curved dga over k (curvature c arbitrary) such that every A\textsuperscript{b} is flat over A\textsuperscript{0}. Let N = (Q, 𝕏) be a quasi-cohesive module, (P, ∂) a bounded cochain complex of finitely generated projective A\textsuperscript{0}-modules, and e\textsuperscript{0} : (P, ∂) → (Q, 𝕏\textsuperscript{0}) a quasi-isomorphism. Then there are a cohesive module (A ⊗\textsubscript{A\textasciicircum{}0} P, 𝔼) with generators P\textsuperscript{p} and 𝔼\textsuperscript{0} = ∂, and a closed degree-0 map e : (A ⊗ P, 𝔼) → N whose (0,0)-component is e\textsuperscript{0} (ε(e(1 ⊗ x)) = e\textsuperscript{0}(x), ε the form-degree-0 part). The proof constructs 𝔼 and e form degree by form degree; at each step it only needs that the cone of e\textsuperscript{0} ∘ − : gr\textsuperscript{b} Hom(P, P) → gr\textsuperscript{b} Hom(P, Q) has no cohomology in degree 1.
\par\smallskip{\footnotesize Source: blueprint.md §2.}
\end{theorem}
\begin{proof}
\leanok
\uses{A6}
See the Lean proof and the cited items.
\paragraph{Blueprint notes.}
Assume each A\textsuperscript{b} is flat over A\textsuperscript{0}. Let N = (Q, 𝕏) be quasi-cohesive, (P, ∂) a bounded complex of finitely generated projective A\textsuperscript{0}-modules, and e\textsuperscript{0} : (P, ∂) → (Q, 𝕏\textsuperscript{0}) a quasi-isomorphism. Then there are a cohesive structure 𝔼 on A ⊗ P with 𝔼\textsuperscript{0} = ∂ and a closed degree-0 e : (A ⊗ P, 𝔼) → N with e\textsuperscript{0} as (0,0)-component. [B, Thm 3.2.7(2) and Lemma 3.2.8], [A, Prop 3.3.4, case 2]

\begin{itemize}
\item Proof: choose connections on each P\textsuperscript{p} (A2.a). Construct 𝔼\textsuperscript{b} and e\textsuperscript{b} by induction on b. At each step the defect pair (P\textsubscript{b}, Q\textsubscript{b}) of [A, (10)–(12)] is a cycle in the cone of the map of A11.a, which is acyclic; a preimage gives (𝔼\textsuperscript{b+1}, e\textsuperscript{b+1}). Stop at b = N.
\item \warn{} Re-derive every sign of [A, (10)–(12)] in the left convention before coding (Antweiler composes with a Koszul sign (−1)\textsuperscript{a'b}, conventions §4).
\item \textbf{A11.a Internal hypothesis.} The proof only uses that, for every b, the map Hom\textsubscript{A\textasciicircum{}0}(P, A\textsuperscript{b} ⊗ P) → Hom\textsubscript{A\textasciicircum{}0}(P, A\textsuperscript{b} ⊗ Q) induced by e\textsuperscript{0} is a quasi-isomorphism ([A]'s form). This follows from the assumptions: A\textsuperscript{b} ⊗ − preserves quasi-isomorphisms, and so does Hom out of a bounded complex of projectives.
\end{itemize}
\end{proof}

\begin{theorem}[A12 — Total quasi-isomorphism]
\label{A12}
\lean{Cohesive.totMap_quasiIso, Cohesive.lifting_quasiIso, Cohesive.QuasiCohesive.totComplex, Cohesive.totMap, Cohesive.Z0.comp0}
\leanok
\uses{A11, A7}
Let A be a bounded graded-commutative dga over k with curvature c = 0 such that every A\textsuperscript{b} is flat over A\textsuperscript{0}. Let φ : M₁ → M₂ be a closed degree-0 map of quasi-cohesive modules whose (0,0)-component φ\textsuperscript{0} : (Q₁, 𝕏₁\textsuperscript{0}) → (Q₂, 𝕏₂\textsuperscript{0}) is a quasi-isomorphism of complexes of A\textsuperscript{0}-modules. Then φ is a quasi-isomorphism of total complexes (Tot Q₁, 𝕏₁) → (Tot Q₂, 𝕏₂) (complexes of k-modules graded by total degree). In particular, in the lifting theorem (A11) with c = 0, the map e is a quasi-isomorphism of total complexes.

Deviation: stated for any such φ (the sources state it for the map e of A11).
\par\smallskip{\footnotesize Source: blueprint.md §2.}
\end{theorem}
\begin{proof}
\leanok
See the Lean proof and the cited items.
\paragraph{Blueprint notes.}
In A11, e is a quasi-isomorphism of total complexes. Proof: A7 with the form-degree filtration. [A, proof of Prop 7.2.18]
\end{proof}

\begin{definition}[A13 — Bridge to Block's definition]
\label{A13}
\lean{Cohesive.TotR.rAct, Cohesive.BlockCohesive, Cohesive.BlockHom, Cohesive.blockδ, Cohesive.blockSwap, Cohesive.BlockQuasiCohesive.toLeft}
\leanok
\uses{A4}
Let A be a bounded graded-commutative curved dga over k with curvature c, and c\textsubscript{B} ∈ A with c = −c\textsubscript{B}. A quasi-cohesive module in Block's convention is a family of A\textsuperscript{0}-modules Q\textsuperscript{p} with a k-linear operator 𝔼 of total degree 1 on the right module (⨁\textsubscript{p} Q\textsuperscript{p}) ⊗\textsubscript{A\textasciicircum{}0} A satisfying 𝔼(eω) = 𝔼(e ⊗ 1)·ω + (−1)\textsuperscript{|e|} e ⊗ dω for e ∈ Q\textsuperscript{p}, and 𝔼 ∘ 𝔼 + (−)·c\textsubscript{B} = 0; its morphisms of degree n are the right A-linear maps of degree n, with dφ = 𝔼₂ ∘ φ − (−1)\textsuperscript{n} φ ∘ 𝔼₁. Transport along σ : e ⊗ a ↦ (−1)\textsuperscript{|e||a|} a ⊗ e gives

\begin{itemize}
\item mutually inverse bijections between Block's (cohesive, quasi-cohesive) modules and those of A3;
\item for each n a k-linear isomorphism from Block's Hom\textsuperscript{n} to Hom\textsuperscript{n} of A4, φ ↦ σ ∘ φ ∘ σ\textsuperscript{−1}, compatible with the differentials, composition and identities.
\end{itemize}
\par\smallskip{\footnotesize Source: blueprint.md §2.}
\end{definition}
\paragraph{Blueprint notes.}
The left-module definition is dg-equivalent to [B, Def 2.3.2] (right modules) for graded-commutative A, with curvature c ↔ −c. Not on the critical path; proves the conventions are faithful.

% Generated by scripts/blueprint_gen.py; do not edit.
\chapter{Sheaves on a compact space}\label{chap:S}
\textbf{Setting for all of §3.} X a compact Hausdorff space, \texttt{X : Type} (conventions §7). Sheaves on \texttt{Opens X} with the standard topology. R a sheaf of commutative rings, modules are Mathlib \texttt{SheafOfModules} ✓ (abelian ✓ \texttt{SheafOfModules.instAbelian}, Grothendieck abelian ✓, hence enough injectives ✓ \texttt{IsGrothendieckAbelian.enoughInjectives}).

\begin{theorem}[S1 — Fine sheaves of rings]
\label{S1}
\lean{Cohesive.SupportedIn, Cohesive.IsFine, Cohesive.existsUnique_extend, Cohesive.exists_hom_extend}
\leanok
Let X be a topological space and R a sheaf of rings on X. A global section ρ ∈ R(X) is \emph{supported in} an open U if there is an open V with U ∪ V = X and ρ|\textsubscript{V} = 0 (equivalently supp ρ ⊆ U). R is \emph{fine} if for every finite open cover (U\textsubscript{i}) of X there are ρ\textsubscript{i} ∈ R(X) with Σ\textsubscript{i} ρ\textsubscript{i} = 1 and each ρ\textsubscript{i} supported in U\textsubscript{i}.

\begin{itemize}
\item \textbf{Extension by zero.} If ρ is supported in U, M is a sheaf of R-modules and t ∈ M(U), there is a unique s ∈ M(X) with s|\textsubscript{U} = ρ|\textsubscript{U} · t and s|\textsubscript{V} = 0 on every open V where ρ|\textsubscript{V} = 0.
\item \textbf{For morphisms.} If R is commutative and ρ is supported in U, every morphism φ : M|\textsubscript{U} → N|\textsubscript{U} gives a morphism ψ : M → N with ψ = ρ·φ over opens V ⊆ U and ψ = 0 over opens where ρ vanishes.
\end{itemize}
Deviation: only finite covers are required (the usual notion on compact spaces).
\par\smallskip{\footnotesize Source: blueprint.md §3.}
\end{theorem}
\begin{proof}
\leanok
See the Lean proof and the cited items.
\paragraph{Blueprint notes.}
S is \emph{fine} if for every finite open cover (U\textsubscript{i}) of X there are ρ\textsubscript{i} ∈ S(X) with Σρ\textsubscript{i} = 1 and supp ρ\textsubscript{i} ⊆ U\textsubscript{i} (support = closed set of points where the germ is nonzero). ✗ in Mathlib.

\begin{itemize}
\item \textbf{S1.a Extension by zero of sections.} If supp ρ ⊆ U and t ∈ M(U) for an S-module M, then ρt extends uniquely to a global section vanishing off supp ρ.
\item Finite covers suffice because X is compact; \warn{} [A] uses softness and arbitrary covers.
\end{itemize}
\end{proof}

\begin{theorem}[S2 — Γ is exact on modules over a fine sheaf]
\label{S2}
\lean{Cohesive.exists_local_preimage_of_epi, Cohesive.surjective_globalSections_of_epi}
\leanok
\uses{S1}
Let R be a sheaf of rings on a topological space X and f : M → N an epimorphism of sheaves of R-modules.

\begin{itemize}
\item f is locally surjective: for every s ∈ N(U) and x ∈ U there are an open x ∈ V ⊆ U and t ∈ M(V) with f(t) = s|\textsubscript{V}.
\item If X is compact and R is fine, then Γ(X, f) : M(X) → N(X) is surjective.
\end{itemize}
\par\smallskip{\footnotesize Source: blueprint.md §3.}
\end{theorem}
\begin{proof}
\leanok
See the Lean proof and the cited items.
\paragraph{Blueprint notes.}
For S fine, M ↠ N an epimorphism of S-modules, Γ(X, M) → Γ(X, N) is surjective. Proof: lift locally on a finite cover, glue with a partition of unity via S1.a. Lean: epi ⇔ locally surjective (?) (\texttt{Sheaf.isLocallySurjective\_iff\_epi'} ✓ for sheaves of types/abelian groups; check the module version).
\end{proof}

\begin{theorem}[S3 — Acyclic complexes of S-modules]
\label{S3}
\lean{Cohesive.acyclic_globalSections, Cohesive.exactAt_globalSections, Cohesive.quasiIso_globalSections}
\leanok
\uses{S2}
Let X be compact and R a fine sheaf of rings on X. For ℤ-indexed cochain complexes of sheaves of R-modules (with R-linear differentials):

\begin{itemize}
\item if K is acyclic, the complex of global sections Γ(X, K) is acyclic;
\item degreewise: if K is exact at n, then Γ(X, K) is exact at n;
\item if f : K → L is a quasi-isomorphism, so is Γ(X, f).
\end{itemize}
No boundedness is assumed.
\par\smallskip{\footnotesize Source: blueprint.md §3.}
\end{theorem}
\begin{proof}
\leanok
See the Lean proof and the cited items.
\paragraph{Blueprint notes.}
For S fine, an acyclic complex of S-modules has acyclic complex of global sections. Proof: split into short exact sequences of cycles; apply S2. No sheaf cohomology needed. Degreewise form: exactness of K at n gives exactness of Γ(K) at n (used in S10). Replaces [A, Lemma 7.2.8] for complexes \emph{of} S-modules (S-linear differentials); for complexes whose terms are S-modules but whose differentials are not S-linear (C4), use S3.b.
\end{proof}

\begin{theorem}[S3.b — Soft sheaves]
\label{S3.b}
\lean{Cohesive.IsSoft, Cohesive.IsSoft.of_action, Cohesive.isSoft_of_isFine, Cohesive.surjective_globalSections_of_isSoft, Cohesive.IsSoft.of_shortExact, Cohesive.acyclic_globalSections_of_isSoft}
\leanok
Let R be a sheaf of rings on a topological space X. A sheaf M of R-modules is \emph{soft} if for every closed K ⊆ X, open U ⊇ K and s ∈ M(U) there are t ∈ M(X) and an open V with K ⊆ V ⊆ U and t|\textsubscript{V} = s|\textsubscript{V}.

\begin{itemize}
\item (i) Let S be a fine sheaf of rings on X acting on M through an additive map S(X) → End(M) with 1 acting as the identity and ρ acting as 0 over every open where ρ vanishes. Then M is soft. In particular every module over a fine commutative sheaf of rings is soft.
\item Let X be compact Hausdorff and 0 → A → B → C → 0 a short exact sequence of R-modules with A soft. (ii) Γ(X, B) → Γ(X, C) is surjective. (iii) If B is also soft, so is C.
\item (iv) On compact Hausdorff X, a bounded-below acyclic complex of R-modules with soft terms has acyclic complex of global sections.
\end{itemize}
Deviation: softness is defined through open neighbourhoods of closed sets rather than stalks or sections over closed sets; on compact Hausdorff spaces this is the usual notion.
\par\smallskip{\footnotesize Source: blueprint.md §3.}
\end{theorem}
\begin{proof}
\leanok
\uses{S2}
See the Lean proof and the cited items.
\paragraph{Blueprint notes.}
X compact Hausdorff. Stalk-free softness: a sheaf F of abelian groups is \emph{soft} if every section over an open neighbourhood of a closed K agrees near K with a global section.

\begin{itemize}
\item \textbf{S3.b(i)} A sheaf of abelian groups on which the global sections of a fine sheaf of rings S act (by sheaf endomorphisms, additively, unitally, with ρ· = 0 on opens where ρ vanishes) is soft. In particular every S-module is soft.
\item \textbf{S3.b(ii)} If 0 → A → B → C → 0 is exact and A is soft, Γ(B) → Γ(C) is surjective. Proof: finite cover with local lifts, shrink to closed sets (normality), glue inductively, correcting by global sections of A (softness).
\item \textbf{S3.b(iii)} If A and B are soft, so is C.
\item \textbf{S3.b(iv)} A bounded-below acyclic complex of sheaves (of R-modules) with soft terms has acyclic complex of global sections: the cycle sheaves are soft by induction from the bottom ((iii)), and (ii) lifts global cycles.
\end{itemize}
Replaces [A, Lemma 7.2.8] (Γ-acyclic terms, Stacks 015E) in C4 without sheaf cohomology.
\end{proof}

\begin{theorem}[S4 — Projective modules over a fine sheaf]
\label{S4}
\lean{Cohesive.projective_free, Cohesive.projective_of_isFinFreeSummand, Cohesive.isFinFreeSummand_of_locally, Cohesive.finite_globalSections, Cohesive.projective_globalSections, Cohesive.bijective_globalSections_map}
\leanok
\uses{S2}
Let R be a sheaf of rings on a topological space X. A sheaf of R-modules is a \emph{finite free summand} if it is a retract of R\textsuperscript{m} for some m ∈ ℕ.

\begin{itemize}
\item If X is compact and R is fine, every free R-module R\textsuperscript{(I)} (any index type I) is projective; hence every finite free summand is projective.
\item If X is compact and R is fine and commutative, a module that is locally a finite free summand (each point has an open U with M|\textsubscript{U} a finite free summand) is globally one.
\item For any R and any finite free summand M: Γ(X, M) is a finitely generated projective R(X)-module, and Γ : Hom\textsubscript{R}(M, N) → Hom\textsubscript{R(X)}(Γ(X, M), Γ(X, N)) is bijective for every R-module N.
\end{itemize}
\par\smallskip{\footnotesize Source: blueprint.md §3.}
\end{theorem}
\begin{proof}
\leanok
\uses{S8}
See the Lean proof and the cited items.
\paragraph{Blueprint notes.}
For S fine:

\begin{itemize}
\item \textbf{S4.a} Finite free S-modules are projective (Hom(S, M) = M(X) and S2).
\item \textbf{S4.b} A locally finitely generated, locally projective S-module is a direct summand of a finite free module (finite cover, partition of unity). [A, Lemma 7.2.15, Cor 7.2.13]
\item \textbf{S4.c} Γ is an equivalence between direct summands of finite free S-modules and finitely generated projective S(X)-modules; inverse P ↦ P̃, P̃(U) = S(U) ⊗\textsubscript{S(X)} P. Moreover Hom\textsubscript{S}(P̃, N) = Hom\textsubscript{S(X)}(P, N(X)) for every S-module N. [A, remark after Def 7.2.5, and proof of Lemma 7.2.7]
\item \textbf{S4.d} For such P̃, Γ(P̃ ⊗\textsubscript{S} N) = P ⊗\textsubscript{S(X)} Γ(N). (v1.9: proved in the stronger form "the sectionwise tensor product P̃ ⊗\textsuperscript{pre} N is already a sheaf", on any site and any object, for S a sheaf of commutative rings; needs no fineness.)
\item Lean: \texttt{SheafOfModules.free} ✓, \texttt{IsLocallyFree} ✓.
\end{itemize}
\end{proof}

\begin{theorem}[S5 — Restriction and extension by zero]
\label{S5}
\lean{Cohesive.restrFreeIso, Cohesive.IsFinFreeSummand.restr, Cohesive.extZero, Cohesive.extZero_preservesFiniteLimits, Cohesive.isIso_extZeroAdj_unit, Cohesive.extZero_locally_zero}
\leanok
Let R be a sheaf of rings on a topological space X and U ⊆ X open. Restriction M ↦ M|\textsubscript{U} sends sheaves of R-modules to sheaves of R|\textsubscript{U}-modules on U; it is exact.

\begin{itemize}
\item R\textsuperscript{(I)}|\textsubscript{U} ≅ (R|\textsubscript{U})\textsuperscript{(I)}, and the restriction of a finite free summand is a finite free summand.
\item Extension by zero j\_! is the left adjoint of restriction. The unit M → (j\_! M)|\textsubscript{U} is an isomorphism; every section of j\_! M vanishes near each point outside U; and j\_! is exact.
\item Consequently restriction preserves injective objects.
\end{itemize}
\par\smallskip{\footnotesize Source: blueprint.md §3.}
\end{theorem}
\begin{proof}
\leanok
\uses{S1}
See the Lean proof and the cited items.
\paragraph{Blueprint notes.}
\begin{itemize}
\item \textbf{S5.a} Restriction M ↦ M|U to an open U, exact. Lean: \texttt{overFunctor.\{0\}} ✓ through the project abbreviation \texttt{restr} (conventions §7, decided in P3); exact because it has adjoints on both sides.
\item \textbf{S5.b} Extension by zero j\_! for an open j : U → X: left adjoint to restriction, exact, stalks (j\_!M)\textsubscript{x} = M\textsubscript{x} for x ∈ U and 0 otherwise. ✗ as such in Mathlib, but \texttt{SheafOfModules.pullback} ✓ (left adjoint of \texttt{pushforward} along a continuous functor) applied to the inclusion of opens may already be j\_! (?); exactness still needs a stalk argument.
\item \textbf{S5.c} Restriction preserves injective objects (it has the exact left adjoint j\_!).
\end{itemize}
\end{proof}

\begin{theorem}[S6 — Locality]
\label{S6}
\lean{Cohesive.shortComplex_exact_iff, Cohesive.acyclic_of_locally, Cohesive.quasiIso_of_locally, Cohesive.quasiIso_of_sectionwise}
\leanok
Let R be a sheaf of rings on a topological space X.

\begin{itemize}
\item A short complex M₁ → M₂ → M₃ of R-modules is exact iff every section of M₂ over U killed by the second map comes, near every point of U, from a section of M₁.
\item A complex K of R-modules is acyclic if every point has an open neighbourhood U with K|\textsubscript{U} acyclic; a map f of complexes is a quasi-isomorphism if every point has an open neighbourhood U with f|\textsubscript{U} a quasi-isomorphism.
\item On any site, a map of complexes of sheaves of modules that is a quasi-isomorphism on sections over every object is a quasi-isomorphism.
\end{itemize}
\par\smallskip{\footnotesize Source: blueprint.md §3.}
\end{theorem}
\begin{proof}
\leanok
\uses{S2, S5}
See the Lean proof and the cited items.
\paragraph{Blueprint notes.}
A map of complexes of R-modules is a quasi-isomorphism iff its restrictions to the members of some open cover are; a complex is acyclic iff locally acyclic. Also: a sectionwise homotopy equivalence of presheaves of complexes induces a quasi-isomorphism of the associated sheaf complexes.
\end{proof}

\begin{definition}[S7 — Hom sheaves]
\label{S7}
\lean{Cohesive.homSheaf, Cohesive.homPre, Cohesive.homPost, Cohesive.homSheafComplex, Cohesive.homComplexSectionsEquiv, Cohesive.homSheafOverIso}
\leanok
\uses{S3.b, S5, S6, S9}
Let R be a sheaf of commutative rings on a site C (in particular on X or on the over-site of an open U). For R-modules M, N the \emph{Hom sheaf} 𝓗om(M, N) is the sheaf of R-modules with 𝓗om(M, N)(V) = Hom(M|\textsubscript{V}, N|\textsubscript{V}), R acting through the target, functorial in both variables (pre- and post-composition). For ℤ-indexed complexes K, L, the \emph{Hom complex} 𝓗om(K, L) is the complex of sheaves of R-modules with 𝓗om(K, L)\textsuperscript{n} = ∏\textsubscript{p+n=q} 𝓗om(K\textsuperscript{p}, L\textsuperscript{q}) and differential δz = z ∘ d + (−1)\textsuperscript{n+1} d ∘ z. Its sections over V are the degree-n cochains K|\textsubscript{V} → L|\textsubscript{V}, and 𝓗om(M, N)|\textsubscript{V} ≅ 𝓗om(M|\textsubscript{V}, N|\textsubscript{V}).
\par\smallskip{\footnotesize Source: blueprint.md §3.}
\end{definition}
\paragraph{Blueprint notes.}
For complexes M, N of R-modules, 𝓗om(M, N) is the complex of sheaves of R-modules with 𝓗om(M, N)(V) = Hom\textsuperscript{•}(M|V, N|V). ✗ for \texttt{SheafOfModules} (Mathlib has \texttt{presheafHom} ✓ for type-valued Homs only).

\begin{itemize}
\item \textbf{S7.a′} (v1.9, replaces S7.a) \emph{Termwise source action.} Let a ring T act on each term M\textsuperscript{i} of a complex of R-modules by R-linear sheaf endomorphisms (additively, unitally; not necessarily commuting with the differential). Precomposition gives an action of T on each term 𝓗om(M, N)\textsuperscript{n} by R-linear sheaf endomorphisms, additive and unital, and vanishing on V whenever the action on every M\textsuperscript{i} vanishes on V. It is not a chain map. With T = S(X) for a fine S this is the input of S3.b(i) (\texttt{IsSoft.of\_action}), used in C4 step 2. (The old S7.a, "a complex of S-modules via the source", is false as stated when the differentials of M are not S-linear, which is exactly the case in C3/C4.)
\item \textbf{S7.b} If W is strictly perfect (S9), 𝓗om(W, −) preserves quasi-isomorphisms.
\item \textbf{S7.c} If I is a bounded-below complex of injectives, then 𝓗om(C, I) is acyclic for every acyclic C (sections over V use S5.c and \texttt{CochainComplex.isKInjective\_of\_injective} ✓).
\end{itemize}

\begin{definition}[S8 — Tensor products and K-flat resolutions]
\label{S8}
\lean{Cohesive.tensorBifunctor, Cohesive.tensorC, Cohesive.tensorCMap₂, Cohesive.IsFlat, Cohesive.IsKFlat, Cohesive.restrTensorIso}
\leanok
\uses{S5, S9}
Let R be a sheaf of commutative rings on a site.

\begin{itemize}
\item The \emph{tensor product} M ⊗ N of sheaves of R-modules is the sheafification of the presheaf V ↦ M(V) ⊗\textsubscript{R(V)} N(V).
\item The \emph{tensor product of complexes} K ⊗ L is the total complex of the double complex K\textsuperscript{p} ⊗ L\textsuperscript{q} (Koszul sign on the differential of the second factor); it is functorial in both variables.
\item M is \emph{flat} if M ⊗ − preserves monomorphisms; a complex K is \emph{K-flat} if K ⊗ L is acyclic for every acyclic complex L.
\item On a space X, tensor products commute with restriction to an open U: (M ⊗ N)|\textsubscript{U} ≅ M|\textsubscript{U} ⊗ N|\textsubscript{U}.
\end{itemize}
\par\smallskip{\footnotesize Source: blueprint.md §3.}
\end{definition}
\paragraph{Blueprint notes.}
This package replaces Malgrange's flatness theorem (feasibility report, G6).

\begin{itemize}
\item \textbf{S8.a} Tensor product and base change of sheaves of modules along R → S. Lean: \texttt{PresheafOfModules} monoidal ✓ plus sheafification ✓; no monoidal structure on \texttt{SheafOfModules} ✗.
\item \textbf{S8.b} Flat R-modules; j\_! R|\textsubscript{U} is flat (stalks, S5.b).
\item \textbf{S8.c} Every R-module is a quotient of a direct sum of modules j\_! R|\textsubscript{U}.
\item \textbf{S8.d} Every bounded-above complex has a quasi-isomorphism from a bounded-above complex of flat modules; bounded-above complexes of flats are K-flat. [Stacks, chapter \emph{Cohomology of Sheaves}, section on K-flat complexes; look up the exact tags in phase 2].
\item \textbf{S8.e} If K is K-flat and acyclic then K ⊗ L is acyclic for every L (resolve L by S8.d).
\item \textbf{S8.f} Restriction to opens preserves K-flatness (use j\_!).
\item \textbf{S8.g} If K, K′ are K-flat and connected by a roof of quasi-isomorphisms, so are S ⊗\textsubscript{R} K and S ⊗\textsubscript{R} K′ (resolve the middle term by S8.d). (v1.9) Lean form: for a single quasi-isomorphism K → K′ between K-flat complexes. For a roof K ← L → K′ with K′ bounded (e.g. strictly perfect) and L arbitrary, as in \texttt{IsPerfect}: first replace L by τ\textsubscript{≤b}L (\texttt{CochainComplex.truncLE}; L → K′ a quasi-isomorphism with K′ in degrees ≤ b makes τ\textsubscript{≤b}L → L one), then resolve τ\textsubscript{≤b}L by S8.d and apply S8.g to both legs.
\end{itemize}

\begin{theorem}[S9 — Perfect complexes]
\label{S9}
\lean{Cohesive.IsStrictlyPerfect, Cohesive.IsPerfect, Cohesive.perf_Q_obj_iff, Cohesive.IsPerfectIn, Cohesive.IsPerfect.exists_isPerfectIn, Cohesive.IsPerfectIn.exactAt}
\leanok
Let R be a sheaf of rings on a topological space X; complexes are ℤ-indexed cochain complexes of sheaves of R-modules.

\begin{itemize}
\item K is \emph{strictly perfect} if K\textsuperscript{n} = 0 outside some interval [a, b] and every K\textsuperscript{n} is a finite free summand (on any site).
\item K is \emph{perfect} if every point has an open neighbourhood U, a strictly perfect complex P of R|\textsubscript{U}-modules and quasi-isomorphisms K|\textsubscript{U} ← L → P.
\item D\textsubscript{perf}(X, R) is the full subcategory of the derived category D(R) on the objects isomorphic to Q(K) for some perfect K; Q(K) lies in D\textsubscript{perf}(X, R) iff K is perfect.
\item If X is compact, every perfect K has local models as above all concentrated in one interval [a, b], and then the homology sheaves of K vanish outside [a, b].
\end{itemize}
Deviation: perfectness is a roof of quasi-isomorphisms on each U instead of an isomorphism in D(R|\textsubscript{U}); the two are equivalent.
\par\smallskip{\footnotesize Source: blueprint.md §3.}
\end{theorem}
\begin{proof}
\leanok
\uses{S5}
See the Lean proof and the cited items.
\paragraph{Blueprint notes.}
\begin{itemize}
\item \textbf{S9.a} A complex of R-modules is \emph{strictly perfect} if it is bounded and each term is a direct summand of a finite free module. [A, Def 7.2.11]
\item \textbf{S9.b} A complex K is \emph{perfect} if every point has an open neighbourhood U, a strictly perfect P on U and quasi-isomorphisms K|U ← L → P (\texttt{QuasiIso} ✓). Equivalent to Q(K|U) ≅ Q(P) in D(R|U), but avoids a derived category for every open set. Lean: with \texttt{overFunctor} (chosen in P3), "strictly perfect on U" lives over \texttt{R.over U} on the over-site \texttt{Over U}, not over a ring on \texttt{Opens U}; S9.d and S10 must use the same API.
\item \textbf{S9.c} D\textsubscript{perf}(X, R) is the full subcategory of \texttt{DerivedCategory (SheafOfModules R)} ✓ on objects isomorphic to Q(K) with K perfect (perfectness is invariant under quasi-isomorphism).
\item \textbf{S9.d} On compact X every perfect complex is globally bounded (uniform degree range and number of generators), by a finite cover.
\end{itemize}
\end{proof}

\begin{theorem}[S10 — Globalization over a fine sheaf]
\label{S10}
\lean{Cohesive.exists_strictlyPerfect_quasiIso, Cohesive.nullHomotopic_locally, Cohesive.lift_locally, Cohesive.split_locally, Cohesive.isFinFreeSummand_locally_of_acyclic, Cohesive.homology_finGen_locally, Cohesive.exists_epi_free_of_locally}
\leanok
\uses{S1, S11, S2, S4, S9}
Let X be compact and R a fine sheaf of commutative rings on X. Every perfect complex K of R-modules receives a quasi-isomorphism P → K from a strictly perfect complex P.

Supporting lemmas, for a sheaf of rings R on any site:

\begin{itemize}
\item a chain map from a strictly perfect complex to an acyclic complex is locally null-homotopic;
\item for E strictly perfect and a quasi-isomorphism s : L → K, every chain map E → K lifts locally through s up to homotopy;
\item an epimorphism onto a finite free summand splits locally; in a bounded acyclic complex in degrees [m, n] whose terms other than the (m+1)-st are finite free summands, that term is locally a finite free summand;
\item if P is strictly perfect and exact in all degrees ≥ k, then H\textsuperscript{k−1}(P) is locally a quotient of a finite free module.
\end{itemize}
And on X with R fine: a module that is locally a quotient of a finite free module is a quotient of a finite free module.
\par\smallskip{\footnotesize Source: blueprint.md §3.}
\end{theorem}
\begin{proof}
\leanok
\uses{S7}
See the Lean proof and the cited items.
\paragraph{Blueprint notes.}
For S fine, every globally bounded perfect complex of S-modules is quasi-isomorphic to a strictly perfect complex of S-modules. [A, Lemma 7.2.16, with Lemmas 7.2.9, 7.2.10]

\begin{itemize}
\item Sub-lemmas (v1.8), on any space Y and sheaf of rings R on Y (applied on opens U through Mathlib's \texttt{Opens.sheafOfModulesEquivOver}):
\begin{itemize}
\item \textbf{S10.a Local null-homotopy.} A chain map from a strictly perfect complex to an acyclic complex is locally null-homotopic (finitely many sections, degree by degree, by local surjectivity of epis, S2 local part).
\item \textbf{S10.b Local lifting.} For E strictly perfect and a quasi-isomorphism s : L → K, every chain map E → K lifts locally through s up to homotopy (S10.a on the cone of s).
\item \textbf{S10.c Local splitting.} An epimorphism onto a direct summand of a finite free module splits locally. Hence, for a bounded acyclic complex C in degrees [m, n] whose terms are summands of finite free modules except possibly C\textsuperscript{m+1}: splitting from the top, the cycle sheaves Z\textsuperscript{i} (i ≥ m + 2) are locally such summands, and C\textsuperscript{m+1} ≅ C\textsuperscript{m} ⊕ Z\textsuperscript{m+2} locally, so C\textsuperscript{m+1} is locally a summand of a finite free module too.
\item \textbf{S10.d} The top homology sheaf of a strictly perfect complex is locally finitely generated (S10.c on the part above it).
\item \textbf{S10.e Global generation.} For R fine on compact X, a locally finitely generated R-module is a quotient of a finite free module (finite cover, S1.a).
\end{itemize}
\item Proof outline (v1.8): by S9.d, K has local strictly perfect models in [a, b] and homology in [a, b]. (1) Downward induction from b to a − 1 builds finite free S-modules Fⁿ and F → K, iso on homology above the current degree, as in [TT, Lemma 1.9.5]: the next homology sheaf of the cone is locally the top homology of the cone of a local lift F|U → P\textsubscript{U} (S10.b), a strictly perfect complex, so locally finitely generated (S10.d), hence a quotient of a finite free module (S10.e); lift to global cycles (S4.a). (2) B = coker(F\textsuperscript{a−1} → F\textsuperscript{a}) and P = (B → F\textsuperscript{a+1} → ⋯ → F\textsuperscript{b}), with P → τ\textsubscript{≥a}K ≃ K a quasi-isomorphism. Locally the model P\textsubscript{U} lifts to a quasi-isomorphism P\textsubscript{U} → P|U (S10.b); its cone is bounded and acyclic in degrees [a − 1, b] with terms summands of finite free modules except cone\textsuperscript{a} = P\textsubscript{U}\textsuperscript{a+1} ⊕ B; S10.c makes B locally a summand of a finite free module, and S4.b globally. So P is strictly perfect, hence K-projective (S4.a, S11), and P → τ\textsubscript{≥a}K lifts to P → K (S11).
\item \warn{} The proof of [A, Lemma 7.2.16] applies [A, Lemma 7.2.10] on small opens U, which needs free S|U-modules to be projective (true for soft S on paracompact U, not given by S1 on compact X); replaced by local lifting and local splitting (S10.a–c) and S4.b (v1.8).
\item \warn{} Compactness replaces [A, Lemma 7.2.12] (finite covering dimension).
\end{itemize}
\end{proof}

\begin{theorem}[S11 — Derived-category facts from Mathlib]
\label{S11}
\lean{Cohesive.exists_injective_resolution, Cohesive.exists_lift_of_isKProjective}
\leanok
Let C be an abelian category.

\begin{itemize}
\item If C has enough injectives, every complex K with K\textsuperscript{i} = 0 for i < n has a quasi-isomorphism K → I to a complex I of injective objects with I\textsuperscript{i} = 0 for i < n (so I is K-injective).
\item If P is K-projective and s : L → K is a quasi-isomorphism, every chain map g : P → K lifts through s up to homotopy: there is f : P → L with s ∘ f homotopic to g.
\end{itemize}
Morphisms in D(C) into a K-injective complex, or out of a K-projective one, are homotopy classes of chain maps; bounded complexes of projectives are K-projective.
\par\smallskip{\footnotesize Source: blueprint.md §3.}
\end{theorem}
\begin{proof}
\leanok
See the Lean proof and the cited items.
\paragraph{Blueprint notes.}
\begin{itemize}
\item Bounded-below injective resolutions: \texttt{exists\_quasiIso\_injective} ✓ (\texttt{DerivabilityStructureInjectives}).
\item Hom in D into a K-injective complex equals homotopy classes: \texttt{CochainComplex.IsKInjective} ✓, \texttt{Qh\_map\_bijective\_of\_isKInjective} ✓.
\item Bounded complexes of projectives are K-projective: \texttt{isKProjective\_of\_projective} ✓.
\item Universes: \texttt{HasDerivedCategory.standard} ✓ as a local instance (conventions §7).
\end{itemize}
\end{proof}

% Generated by scripts/blueprint_gen.py; do not edit.
\chapter{The comparison theorem}\label{chap:C}

\begin{definition}[C0 — Resolving data]
\label{C0}
\lean{Cohesive.DgaPresheaf, Cohesive.DgaPresheaf.pre0, Cohesive.PreResolvingData, Cohesive.ResolvingData, Cohesive.PreResolvingData.piecePre, Cohesive.PreResolvingData.pieceSheaf}
\leanok
\uses{A1, A10, A9, S4, S6, S8}
Let X be a topological space. \emph{Pre-resolving data} on X consist of:

\begin{itemize}
\item for every open U a bounded graded-commutative ℂ-dga 𝒜(U) (A1), with restriction maps of dgas for V ⊆ U that are functorial; sections glue uniquely, being of degree q is a local property, and 𝒜\textsuperscript{0} : U ↦ 𝒜\textsuperscript{0}(U) is a sheaf of commutative rings;
\item curvature c = 0 in every 𝒜(U);
\item a sheaf of commutative rings O and ring maps ρ\textsubscript{U} : O(U) → 𝒜\textsuperscript{0}(U), natural in U, with d ∘ ρ\textsubscript{U} = 0;
\item 𝒜\textsuperscript{0} is fine (S1), and every 𝒜\textsuperscript{q} is locally free of finite rank over 𝒜\textsuperscript{0}: each point has an open U and finitely many e\textsubscript{i} ∈ 𝒜\textsuperscript{q}(U) such that, over every V ⊆ U, each β ∈ 𝒜\textsuperscript{q}(V) is uniquely Σ\textsubscript{i} f\textsubscript{i}·e\textsubscript{i}|\textsubscript{V} with f\textsubscript{i} ∈ 𝒜\textsuperscript{0}(V).
\end{itemize}
From these, each 𝒜\textsuperscript{q} is a sheaf of 𝒜\textsuperscript{0}-modules. \emph{Resolving data} add:

\begin{itemize}
\item O → 𝒜\textsuperscript{0} → 𝒜\textsuperscript{1} → ⋯ is a resolution: each ρ\textsubscript{U} is injective, every closed section of 𝒜\textsuperscript{0}(U) lies in the image of ρ\textsubscript{U}, and in degrees q + 1 ≥ 1 every closed form is locally exact;
\item local lifting (LL): for every cohesive module M over 𝒜(X) and x ∈ X there are an open x ∈ U′, a bounded complex W of finitely generated projective O(U′)-modules and an isomorphism in Ho(P\textsubscript{𝒜(U′)}) between M|\textsubscript{U′} and 𝒜(U′) ⊗\textsubscript{O(U′)} W.
\end{itemize}
Deviation: LL is required only for restrictions of global cohesive modules, and K-flatness of 𝒜 over O is not assumed. No compactness or Hausdorff condition is part of the data.
\par\smallskip{\footnotesize Source: blueprint.md §4.}
\end{definition}
\paragraph{Blueprint notes.}
The abstract input, Antweiler's Assumptions 7.2 at h = 0 in the form used here. It is split in two structures: **C0.a \texttt{PreResolvingData}\textbf{ has every field below except the resolution condition and LL; }C0.b \texttt{ResolvingData}** extends it by those two. C1 (the functor J) uses only C0.a, so J can be defined for X without H1 and H2 (see §1).

\begin{itemize}
\item X compact Hausdorff, \texttt{X : Type}; R a sheaf of commutative ℂ-algebras on X.
\item For each open U a bounded graded-commutative dga 𝒜(U) (A1, degrees ≤ N), restriction maps that are dga maps and functorial; each 𝒜\textsuperscript{q} a sheaf of ℂ-modules.
\item 𝒜\textsuperscript{0} a fine sheaf of commutative rings (S1); each 𝒜\textsuperscript{q} a locally free 𝒜\textsuperscript{0}-module of finite rank.
\item (v1.9, audit F1) \textbf{One carrier for A⁰.} The sheaves are not separate data: they are derived from the presheaf of dgas U ↦ (A\textsubscript{U}, 𝒜\textsubscript{U}). 𝒜\textsuperscript{0} is the \texttt{Sheaf \_ CommRingCat} with 𝒜\textsuperscript{0}(U) = \texttt{CommRingCat.of ↥(𝒜\_U 0)} and the restriction maps of the dgas; it is used as a \texttt{RingSheaf} through \texttt{ringOf} (S8). 𝒜\textsuperscript{q}(U) = ↥(𝒜\textsubscript{U} q) as an 𝒜\textsuperscript{0}(U)-module (\texttt{GCDga.piece}). The sheaf conditions are fields. So Γ(X, 𝒜\textsuperscript{0}) and A⁰ = ↥(𝒜\textsubscript{X} 0) are the same type with the same ring structure (prototype P7: \texttt{rfl} at instance transparency), modules over Γ(X, 𝒜\textsuperscript{0}) are \texttt{𝒜 0}-modules (A11 applies to P = Γ(P̃) without transport), and ρ\textsubscript{U} : R(U) →+* ↥(𝒜\textsubscript{U} 0) is the ring map A10 takes.
\item ρ : R → 𝒜\textsuperscript{0} a map of sheaves of rings with d∘ρ = 0 such that 0 → R → 𝒜\textsuperscript{0} → 𝒜\textsuperscript{1} → ⋯ is exact (a resolution).
\item (v1.10) \textbf{Flat curvature: c = 0} in every 𝒜(U) (a field of C0.a). Needed for α(M) to be a complex (𝕏² = c·(−)) and for \texttt{fromComplex} in LL (A10 takes \texttt{hc : c = 0}). It holds for the Dolbeault dga by definition (F10), so F14 fills it; it is not an analytic hypothesis. Curved data (c ≠ 0) is the twisted case Z3 (LATER).
\item \textbf{LL (local lifting).} For every cohesive module M over 𝒜(X) and x ∈ X, there are an open x ∈ U′, a bounded complex W of finitely generated projective R(U′)-modules and a homotopy equivalence in P\textsubscript{𝒜(U′)} between M|U′ (A9) and 𝒜(U′) ⊗\textsubscript{R(U′)} W (A10).
\item \warn{} [A] Assumption 6 is a sheaf version for modules on every open set. C2 and C4 only ever apply LL to restrictions of global modules, so this weaker, sectionwise form suffices; it is also the form of the gauge step in Block's proof of [B, Lemma 4.1.5].
\item \warn{} The first half of [A] Assumption 5 (𝒜 K-flat over R) is \textbf{not} assumed: C5 uses S8 instead. (Its second half, each 𝒜\textsuperscript{q} flat over 𝒜\textsuperscript{0}, follows from local freeness above.)
\end{itemize}

\begin{definition}[C1 — The functor J]
\label{C1}
\lean{Cohesive.PreResolvingData.sheafOf, Cohesive.PreResolvingData.JK, Cohesive.PreResolvingData.J, Cohesive.PreResolvingData.sheafOfOver}
\leanok
\uses{A9, C0, S4}
Let pre-resolving data on X be given (C0), A = 𝒜(X) and R = O. For a cohesive module M = (Q, 𝕏) over A, α(M) is the complex of sheaves of O-modules with α(M)\textsuperscript{n}(V) the degree-n part of the base change M|\textsubscript{V} (A9), i.e. of 𝒜(V) ⊗\textsubscript{𝒜\textasciicircum{}0(V)} ⨁\textsubscript{p} (𝒜\textsuperscript{0}(V) ⊗\textsubscript{A\textasciicircum{}0} Q\textsuperscript{p}) ≅ ⨁\textsubscript{p+q=n} 𝒜\textsuperscript{q}(V) ⊗\textsubscript{A\textasciicircum{}0} Q\textsuperscript{p}, O acting through ρ, with differential 𝕏 of M|\textsubscript{V} (d² = 0 because c = 0).

\begin{itemize}
\item Closed degree-0 maps give chain maps and coboundaries give null-homotopic maps, so M ↦ α(M) defines J\textsubscript{K} : Ho(P\textsubscript{A}) ⥤ K(O) and J : Ho(P\textsubscript{A}) ⥤ D(O), the composite with the localization.
\item For an open U and a cohesive module M over 𝒜(U), α\textsubscript{U}(M) is the same construction on the over-site of U, and α(M)|\textsubscript{U} ≅ α\textsubscript{U}(M|\textsubscript{U}).
\end{itemize}
J lands in the whole derived category D(O); that it lands in D\textsubscript{perf} is C2.
\par\smallskip{\footnotesize Source: blueprint.md §4.}
\end{definition}
\paragraph{Blueprint notes.}
Input: pre-resolving data (C0.a) only; J : Ho(P\textsubscript{A}) ⥤ D(R), into the whole derived category. For a cohesive module M = (Q, 𝔼) over A = 𝒜(X), α(M) is the complex of sheaves of R-modules α(M)\textsuperscript{n}(V) = ⊕\textsubscript{p+q=n} 𝒜\textsuperscript{q}(V) ⊗\textsubscript{𝒜\textasciicircum{}0(X)} Q\textsuperscript{p} with the differential of the base change M|V (A9).

\begin{itemize}
\item \textbf{C1.a} It is a sheaf (Q\textsuperscript{p} is finitely generated projective; S4.c). R acts through ρ, and the differential is R-linear because d∘ρ = 0.
\item \textbf{C1.b} Closed degree-0 maps go to chain maps, homotopies to homotopies, so J : Ho(P\textsubscript{A}) ⥤ K(R) ⥤ D(R). [B, §4.1, the sheaves ℰ\textsubscript{X}]
\item \textbf{C1.c} α(M)|U ≅ α\textsubscript{U}(M|U), naturally (sheaf on U from the cohesive module over 𝒜(U)).
\end{itemize}

\begin{theorem}[C2 — J lands in perfect complexes]
\label{C2}
\lean{Cohesive.ResolvingData.quasiIso_of_gamma_eq, Cohesive.ResolvingData.exists_quasiIso_sheafOfOver_fromComplex, Cohesive.ResolvingData.isPerfect_sheafOf, Cohesive.ResolvingData.perf_J_obj}
\leanok
\uses{C1, C7, S4, S9}
Let resolving data be given on a compact Hausdorff space X.

\begin{itemize}
\item For an open U and a bounded complex W of finitely generated projective O(U)-modules there are a strictly perfect complex L of O|\textsubscript{U}-modules and a quasi-isomorphism L → α\textsubscript{U}(𝒜(U) ⊗\textsubscript{O(U)} W) (A10).
\item For every cohesive module M over 𝒜(X), α(M) is a perfect complex; hence J(M) ∈ D\textsubscript{perf}(X, O).
\end{itemize}
Deviation: X is assumed compact Hausdorff, which the proof uses through the explicit model C7.
\par\smallskip{\footnotesize Source: blueprint.md §4.}
\end{theorem}
\begin{proof}
\leanok
See the Lean proof and the cited items.
\paragraph{Blueprint notes.}
For every cohesive M, α(M) is perfect (S9.b). [A, Lemma 7.2.17] Proof: by LL and C1.c, α(M)|U′ is homotopy equivalent to 𝒜|U′ ⊗\textsubscript{R} W̃, W̃(V) = R(V) ⊗ W strictly perfect; W̃ → 𝒜 ⊗\textsubscript{R} W̃ is a quasi-isomorphism because ρ is a resolution and each W̃\textsuperscript{p} is a direct summand of a finite free module.

\begin{itemize}
\item (v1.11) \textbf{Proof via C7} (the identification above was not itemized). X compact Hausdorff. e : Tot(𝒜(U′) ⊗ W) → Tot(N\textsubscript{W̃}(U′)), a ⊗ (α ⊗ w) ↦ aα ⊗ η(1 ⊗ w̃) with w̃ ∈ Γ(W̃) the section corresponding to w. gr\textsuperscript{0} Ξ\textsubscript{e} is an isomorphism (both sides summands of finite free 𝒜⁰|U′-modules with the same global sections, S4.c on the over-site), so Ξ\textsubscript{e} is a quasi-isomorphism (C7.f). u ≫ Ξ\textsubscript{e} = (W̃ → T(W̃)) (both determined by Γ, S4.c), a quasi-isomorphism by C7.d; so u is one (2-out-of-3).
\item (v1.9, audit F5) \textbf{The roof.} \texttt{IsPerfect} wants α(M)|U′ ← L → P. Take L := W̃, P := W̃, g := id and f := α(h) ∘ u : W̃ → 𝒜|U′ ⊗\textsubscript{R} W̃ → α(M)|U′, where h : 𝒜(U′) ⊗ W → M|U′ is the direction of the LL homotopy equivalence pointing into M|U′ (a homotopy equivalence, so α(h) is a quasi-isomorphism).
\item (v1.9) \textbf{W̃ without fineness.} R|U′ is not fine. W̃ is built on the over-site of U′ from S4.c, whose parts used here (\texttt{exists\_globalSections\_iso}, the Hom bijection) need neither fineness nor compactness; transfer to the over-site through \texttt{TopologicalSpace.Opens.sheafOfModulesEquivOver}, or restate S4.c on any site. The differentials of W lift uniquely by the Hom bijection.
\end{itemize}
\end{proof}

\begin{theorem}[C3 — Hom sheaves of cohesive modules]
\label{C3}
\lean{Cohesive.PreResolvingData.homAlong, Cohesive.PreResolvingData.homTop, Cohesive.PreResolvingData.isSoft_homSh}
\leanok
\uses{A4, S3.b, S4, S7}
Let pre-resolving data be given on X and M₁, M₂ cohesive modules over 𝒜(U₀). For V ⊆ U₀ let ℋ\textsuperscript{n}(M₁, M₂)(V) be the ℂ-linear maps χ : Tot(M₁) → Tot(M₂|\textsubscript{V}) of degree n with χ(a·x) = (−1)\textsuperscript{n|a|} a|\textsubscript{V}·χ(x); by base change these are the degree-n maps M₁|\textsubscript{V} → M₂|\textsubscript{V}.

\begin{itemize}
\item For U₀ = X these form a complex ℋ(M₁, M₂) of sheaves of O-modules with δχ = 𝕏₂ ∘ χ − (−1)\textsuperscript{n} χ ∘ 𝕏₁, and ℋ(M₁, M₂)(X) ≅ Hom\textsuperscript{•\_}\{P\textsubscript{A}\}(M₁, M₂) in each degree.
\item If X is compact Hausdorff, every term ℋ\textsuperscript{n}(M₁, M₂) is soft (S3.b): 𝒜\textsuperscript{0}(X) acts by multiplication. The differential is not 𝒜\textsuperscript{0}-linear.
\end{itemize}
\par\smallskip{\footnotesize Source: blueprint.md §4.}
\end{theorem}
\begin{proof}
\leanok
\uses{Z3.m}
See the Lean proof and the cited items.
\paragraph{Blueprint notes.}
For cohesive M₁, M₂ over A, ℋ(M₁, M₂)(V) = Hom\textsuperscript{•\_}\{P\textsubscript{𝒜(V)}\}(M₁|V, M₂|V). It is a complex of sheaves of R-modules whose terms are 𝒜\textsuperscript{0}-modules (A4.c and S4.c); \warn{} the differential is not 𝒜\textsuperscript{0}-linear (δ(fφ) = df·φ + f·δφ). ℋ(M₁, M₂)(X) = Hom\textsuperscript{•\_}\{P\textsubscript{A}\}(M₁, M₂).

\begin{itemize}
\item \textbf{C3.a} (v1.9, audit F3) \emph{The actions needed for softness.} (i) 𝒜\textsuperscript{0}(X) acts on each term α(M₁)\textsuperscript{i} by R-linear sheaf endomorphisms (multiplication; R-linear because R → 𝒜\textsuperscript{0} is a map of commutative rings), additively, unitally, vanishing on V where f|V = 0; by S7.a′ it acts on each term of 𝓗om\textsubscript{R}(α(M₁), N) for every N. (ii) 𝒜\textsuperscript{0}(X) acts on each term of ℋ(M₁, M₂) (by restriction and multiplication in 𝒜(V)), with the same three properties. Both actions are termwise, not chain maps; these are the inputs of S3.b(i) in C4 step 2.
\end{itemize}
\end{proof}

\begin{theorem}[C4 — Full faithfulness]
\label{C4}
\lean{Cohesive.ResolvingData.bijective_J_map, Cohesive.ResolvingData.quasiIso_Φ, Cohesive.ResolvingData.bijective_J_map_comp, Cohesive.PreResolvingData.Ψ_top, Cohesive.ResolvingData.quasiIso_Ψ_comp, Cohesive.PreResolvingData.exists_isStrictlyGE_ℋ, Cohesive.ResolvingData.bijective_J_map}
\leanok
\uses{C2, C3, S11, S3.b, S7}
Let resolving data be given on a compact Hausdorff space X. Then for all M, N in Ho(P\textsubscript{A}), J : Hom(M, N) → Hom\textsubscript{D(O)}(J M, J N) is bijective. Main steps:

\begin{itemize}
\item Ψ : ℋ(M₁, M₂) → 𝓗om(α(M₁), α(M₂)), χ ↦ α(χ), is a map of complexes, and ℋ(M₁, M₂) is bounded below; at X it sends a closed degree-0 map φ to α(φ).
\item For a quasi-isomorphism r : α(M₂) → I to a bounded-below complex of injectives, r\textsubscript{*} ∘ Ψ : ℋ(M₁, M₂) → 𝓗om(α(M₁), I) is a quasi-isomorphism (locally via LL, then S6).
\item Hence f ↦ J(f) ∘ Q(r) is a bijection Hom\textsubscript{Ho(P\_A)}(M, N) → Hom\textsubscript{D(O)}(α(M), I).
\end{itemize}
\par\smallskip{\footnotesize Source: blueprint.md §4.}
\end{theorem}
\begin{proof}
\leanok
See the Lean proof and the cited items.
\paragraph{Blueprint notes.}
J is fully faithful. [A, Prop 7.2.19]

\begin{enumerate}
\item Choose a bounded-below injective resolution r : α(M₂) → I (S11).
\item Ψ : ℋ(M₁, M₂) → 𝓗om\textsubscript{R}(α(M₁), I), φ ↦ r ∘ α(φ), is a map of complexes of sheaves of R-modules; 𝒜\textsuperscript{0}(X) acts on the terms of both sides (C3.a, S7.a′), hence they are soft (S3.b(i)), but the differentials are not 𝒜\textsuperscript{0}-linear (v1.7). The terms of cone(Ψ) are biproducts of terms of both sides, soft by the closure lemmas of S3.b.
\item Ψ is a quasi-isomorphism of sheaf complexes. Locally (S6), at x take U′ and the homotopy equivalence h : 𝒜(U′) ⊗ W → M₁|U′ from LL. In the square
\end{enumerate}
\begin{verbatim}
   ℋ(M₁,M₂)|U′ ───Ψ───> 𝓗om_R(α(M₁), I)|U′
       │ h^*                  │ α(h)^*
   ℋ(𝒜⊗W, M₂) ───Ψ───> 𝓗om_R(𝒜 ⊗_R W̃, I)
       │ ≅ (A10.a)            │ u^*,  u : W̃ → 𝒜 ⊗_R W̃
   𝓗om_R(W̃, α(M₂)) ──r_*──> 𝓗om_R(W̃, I)
\end{verbatim}
both upper vertical maps are sectionwise homotopy equivalences, u\textsuperscript{\emph{} is a quasi-isomorphism by S7.c (cone(u) is acyclic), and r\textsubscript{}} is one by S7.b. (v1.11) Lean: the square lives on the over-site of U′, so it needs ℋ\textsubscript{U′} (C3's section data with base U′) with ℋ(M₁, M₂)|U′ ≅ ℋ\textsubscript{U′}(M₁|U′, M₂|U′), and Ψ\textsubscript{U′} compatible with Ψ under this iso, \texttt{restrSheafOfIso} and \texttt{homComplexOverIso}; u is C2's map.

\begin{enumerate}
\item Γ(X, Ψ) is a quasi-isomorphism: S3.b(iv) applied to cone(Ψ), bounded below (α(M₁) and ℋ bounded, I bounded below) with soft terms (v1.7; S3 does not apply, see C3).
\item Take H\textsuperscript{0}: Hom\textsubscript{Ho(P\_A)}(M₁, M₂) ≅ H\textsuperscript{0} Hom(α(M₁), I) ≅ Hom\textsubscript{D(R)}(α(M₁), α(M₂)) (S11), and check that the composite is J on morphisms.
\end{enumerate}
\end{proof}

\begin{theorem}[C5 — Essential surjectivity]
\label{C5}
\lean{Cohesive.ResolvingData.exists_J_obj_iso, Cohesive.PreResolvingData.ρSheaf, Cohesive.ResolvingData.exists_flat_model, Cohesive.ResolvingData.isPerfect_sheafBC, Cohesive.ResolvingData.exists_sheafOfMap_quasiIso, Cohesive.ResolvingData.quasiIso_tensorCMap_ρhat, Cohesive.ResolvingData.tensorC_single_unit_iso}
\leanok
\uses{A10, A11, A7, A9, C7, S10, S11, S3, S4, S8, S9}
Let resolving data be given on a compact Hausdorff space X. Every object of D\textsubscript{perf}(X, O) is isomorphic to J(M) for some cohesive module M over 𝒜(X). Main steps:

\begin{itemize}
\item a perfect K receives a quasi-isomorphism from a bounded-above complex K′ of flat O-modules;
\item for such K′ (perfect), the base change 𝒜\textsuperscript{0} ⊗\textsubscript{O} K′ along ρ : O → 𝒜\textsuperscript{0} is perfect;
\item given a strictly perfect complex P̃ of 𝒜\textsuperscript{0}-modules and a quasi-isomorphism P̃ → 𝒜\textsuperscript{0} ⊗\textsubscript{O} K′, there are a cohesive M and a quasi-isomorphism α(M) → 𝒜\textsuperscript{•} ⊗\textsubscript{O} K′ (𝒜\textsuperscript{•} the first factor);
\item for bounded-above, termwise flat K′ (no compactness needed), ρ̂ ⊗ K′ : O[0] ⊗ K′ → 𝒜\textsuperscript{•} ⊗ K′ is a quasi-isomorphism, and O[0] ⊗ K′ ≅ K′.
\end{itemize}
Deviation: K-flat resolutions replace Malgrange's flatness theorem.
\par\smallskip{\footnotesize Source: blueprint.md §4.}
\end{theorem}
\begin{proof}
\leanok
See the Lean proof and the cited items.
\paragraph{Blueprint notes.}
Every object of D\textsubscript{perf}(X, R) is isomorphic to some J(M). [A, Prop 7.2.18], \warn{} K-flat variant.

\begin{itemize}
\item \textbf{C5.0 Tensor input} (the only place where flatness enters). For every perfect K there is a K′ connected to K by quasi-isomorphisms such that (i) K′ → 𝒜 ⊗\textsubscript{R} K′ is a quasi-isomorphism and (ii) 𝒜\textsuperscript{0} ⊗\textsubscript{R} K′ is a globally bounded perfect complex of 𝒜\textsuperscript{0}-modules. Steps 1, 2 and 6 below prove C5.0 with S8; steps 3–5 and 7 use only C5.0. Keeping this interface lets C5 be written before S8 is finished, and keeps open the alternative proof by Malgrange flatness (K′ = K).
\item Let K be perfect. It is globally bounded (S9.d); replace it by a bounded complex (truncations \texttt{CochainComplex.truncGE} ✓, \texttt{truncLE}) and then by a K-flat resolution K′ → K (S8.d).
\item 𝒜\textsuperscript{0} ⊗\textsubscript{R} K′ is a globally bounded perfect complex of 𝒜\textsuperscript{0}-modules (S8.f, S8.g). (v1.9, audit F4) The middle term of each local roof K′|U ← L → P\textsubscript{U} is unbounded; truncate it to τ\textsubscript{≤b}L and resolve by S8.d before applying S8.g, as in S8.g's v1.9 note.
\item By S10 there is a strictly perfect complex P̃ of 𝒜\textsuperscript{0}-modules quasi-isomorphic to it; P̃ is K-projective (S4.a, S11), so the quasi-isomorphism is an honest map f\textsuperscript{0} : P̃ → 𝒜\textsuperscript{0} ⊗\textsubscript{R} K′.
\item Pass to global sections: P = Γ(P̃) is a bounded complex of finitely generated projective 𝒜\textsuperscript{0}(X)-modules (S4.c); Γ(f\textsuperscript{0}) is a quasi-isomorphism (S3); N := the quasi-cohesive module of C5.a. Each 𝒜\textsuperscript{q}(X) is projective, hence flat, over 𝒜\textsuperscript{0}(X) (S4.b, S4.c), so A11 gives a cohesive M on 𝒜(X) ⊗ P and a closed e : M → N extending Γ(f\textsuperscript{0}).
\item Localize e to a map of complexes of sheaves α(M) → 𝒜 ⊗\textsubscript{R} K′ (C5.b); it is a quasi-isomorphism (C5.c).
\end{itemize}
(v1.11) \textbf{C5.a–c are proved via C7} with U = X and K = K′, which replaces the sheaf-tensor arguments below (they relied on a base-change associativity that S8.a does not contain, and on an 𝒜(V)-action on sections of a sheafified tensor product): N := N\textsubscript{K′}(X) (C7.c); A11 with e⁰ = Γ(f⁰) (a quasi-isomorphism by S3: 𝒜⁰ is fine, no boundedness needed); C5.b := Ξ\textsubscript{e} (C7.e); C5.c := C7.f, with gr\textsuperscript{0} Ξ\textsubscript{e} ≅ f⁰ via Γ (S4.c); step 7 goes through T(K′) and C7.d. The text below is kept as the informal picture.

Sub-items (v1.9, audit F2):

\begin{itemize}
\item \textbf{C5.a Quasi-cohesive structure on Γ(𝒜 ⊗\textsubscript{R} K′).} Q\textsuperscript{p} := Γ(𝒜\textsuperscript{0} ⊗\textsubscript{R} K′\textsuperscript{p}) (an 𝒜\textsuperscript{0}(X)-module). This is \emph{not} A10 (K′ is not finitely generated projective, and Γ(𝒜 ⊗\textsubscript{R} K′) ≠ 𝒜(X) ⊗\textsubscript{R(X)} Γ(K′)). Instead: 𝒜\textsuperscript{q} ⊗\textsubscript{R} K′\textsuperscript{p} ≅ 𝒜\textsuperscript{q} ⊗\textsubscript{𝒜\textasciicircum{}0} (𝒜\textsuperscript{0} ⊗\textsubscript{R} K′\textsuperscript{p}) (base-change associativity, S8.a), and Γ of the right side is 𝒜\textsuperscript{q}(X) ⊗\textsubscript{𝒜\textasciicircum{}0(X)} Q\textsuperscript{p} by S4.d (𝒜\textsuperscript{q} is a summand of a finite free 𝒜\textsuperscript{0}-module, F11 and S4.b). 𝕏 is Γ of the sheaf operator 𝕏(b ⊗ k) = db ⊗ k + (−1)\textsuperscript{|b|} b ⊗ d\textsubscript{K′}k (well defined because d∘ρ = 0), transported along these isomorphisms; Leibniz and 𝕏² = 0 hold sectionwise, hence after transport. (v1.10, sign fix) The Koszul sign is essential: without it Leibniz fails whenever 𝒜 ≠ 𝒜\textsuperscript{0} (𝕏(a·(b⊗k)) = da·(b⊗k) + (−1)\textsuperscript{|a|} a·𝕏(b⊗k) needs (−1)\textsuperscript{|a|+|b|} on ab ⊗ dk). This operator is the differential of \texttt{tensorC R 𝒜\textasciicircum{}• K′} with 𝒜\textsuperscript{•} as the \textbf{first} factor (Mathlib puts the sign on the second factor), as fixed in conventions §2. Its (0,0)-component is Γ(𝒜\textsuperscript{0} ⊗\textsubscript{R} d\textsubscript{K′}), so A11's target complex (Q, 𝕏\textsuperscript{0}) is Γ(𝒜\textsuperscript{0} ⊗\textsubscript{R} K′).
\item \textbf{C5.b The sheaf map.} On V, e\textsubscript{V} : 𝒜(V) ⊗\textsubscript{𝒜\textasciicircum{}0(X)} P → Γ(V, 𝒜 ⊗\textsubscript{R} K′), a ⊗ p ↦ a·res\textsubscript{V}(e(1 ⊗ p)). (v1.10) No sign: e is closed of degree 0, and a degree-0 map satisfies φ(a·m) = a·φ(m) (conventions §2, Hom\textsuperscript{n} with n = 0). For a map of degree n the generator extension carries (−1)\textsuperscript{n|a|}. N|V is \emph{not} the base change of N, so A9 does not supply the map; α(M)(V) is A9's base change of M, and e\textsubscript{V} is defined on generators. It is a chain map because res\textsubscript{V} is a dga map and e is closed; it is compatible with restriction, so it defines α(M) → 𝒜 ⊗\textsubscript{R} K′ (S4.c Hom property on each α(M)\textsuperscript{n}, which is a finite sum of summands of finite free 𝒜\textsuperscript{0}-modules).
\item \textbf{C5.c Quasi-isomorphism.} Form-degree filtration on both sides as complexes of sheaves of abelian groups; A7 applies as stated (\texttt{quasiIso\_of\_finite\_filtration} is proved in any abelian category). On gr\textsuperscript{b} the map is 𝒜\textsuperscript{b} ⊗\textsubscript{𝒜\textasciicircum{}0} f\textsuperscript{0} (up to the shift sign of A12), a quasi-isomorphism because 𝒜\textsuperscript{b} is a summand of a finite free 𝒜\textsuperscript{0}-module (tensoring with it preserves quasi-isomorphisms; flat by S8.b, or directly).
\item K′ → 𝒜 ⊗\textsubscript{R} K′ is a quasi-isomorphism (𝒜\textsuperscript{•} first, as in C5.a): its cone is cone(ρ) ⊗\textsubscript{R} K′, acyclic because ρ is a resolution and K′ is K-flat.
\item Hence α(M) ≅ 𝒜 ⊗\textsubscript{R} K′ ≅ K′ ≅ K in D(R).
\end{itemize}
\end{proof}

\begin{theorem}[C6 — Comparison theorem for resolving data]
\label{C6}
\lean{Cohesive.ResolvingData.comparison}
\leanok
\uses{C2, C4, C5}
Let resolving data be given on a compact Hausdorff space X, with A = 𝒜(X). Then the functor J : Ho(P\textsubscript{A}) ⥤ D(O) is full and faithful, and its essential image is exactly D\textsubscript{perf}(X, O).

This is the abstract comparison theorem, for any resolving data (C0). Its application to the Dolbeault data is F15 (\texttt{blockPerf}), which is the proof of M.
\par\smallskip{\footnotesize Source: blueprint.md §4.}
\end{theorem}
\begin{proof}
\leanok
See the Lean proof and the cited items.
\paragraph{Blueprint notes.}
For resolving data, J : Ho(P\textsubscript{𝒜(X)}) ⥤ D(R) is full and faithful (C4) with essential image exactly D\textsubscript{perf}(X, R) (C2, C5); equivalently its corestriction to D\textsubscript{perf} is an equivalence. [A, Thm 7.2.20 at h = 0], [CHL, Thm 7.10] (for twisted modules) and [CHL, Cor 7.13] (the same for cohesive modules, via [CHL, Cor 3.4]) (full faithfulness: [CHL, Prop 7.4]).
\end{proof}

\begin{theorem}[C7 — The explicit model T(K) of 𝒜 ⊗\textsubscript{R} K]
\label{C7}
\lean{Cohesive.PreResolvingData.exists_TC_iso_tensorC, Cohesive.PreResolvingData.quasiIso_Ξ, Cohesive.PreResolvingData.Conn, Cohesive.PreResolvingData.Conn.N, Cohesive.PreResolvingData.TSh, Cohesive.PreResolvingData.Conn.TC, Cohesive.PreResolvingData.Ξ}
\leanok
\uses{A3, A6, A7, A9, S4, S8}
Let pre-resolving data be given on a compact space X, U ⊆ X open, and F a complex of 𝒜\textsuperscript{0}|\textsubscript{U}-modules on the over-site of U. A \emph{flat connection} on F is, for every V, generator data θ\textsubscript{V} on the 𝒜\textsuperscript{0}(V)-modules F\textsuperscript{p}(V) of total degree 1, flat (𝕏\textsubscript{V²} = c = 0), with form-degree-0 part d\textsubscript{F}, and compatible with restriction. Given one:

\begin{itemize}
\item N(V) = (⨁\textsubscript{p} F\textsuperscript{p}(V), 𝕏\textsubscript{V}) is a quasi-cohesive module over 𝒜(V);
\item T(F)\textsuperscript{n}(V) = (𝒜(V) ⊗\textsubscript{𝒜\textasciicircum{}0(V)} ⨁\textsubscript{p} F\textsuperscript{p}(V))\textsuperscript{n} is a sheaf of O|\textsubscript{U}-modules, and with 𝕏 a complex T;
\item a closed degree-0 map e : M → N(U) from a cohesive module over 𝒜(U) induces a map of complexes Ξ\textsubscript{e} : α\textsubscript{U}(M) → T.
\item For K a complex of O|\textsubscript{U}-modules and F = 𝒜\textsuperscript{0} ⊗\textsubscript{O} K, with a connection satisfying ∇(1 ⊗ k) = 0 (which exists), T ≅ 𝒜\textsuperscript{•}|\textsubscript{U} ⊗ K, and under this and K ≅ O|\textsubscript{U}[0] ⊗ K the map K → T, k ↦ 1 ⊗ (1 ⊗ k), is ρ̂|\textsubscript{U} ⊗ K.
\end{itemize}
\par\smallskip{\footnotesize Source: blueprint.md §4.}
\end{theorem}
\begin{proof}
\leanok
\uses{A11, C1}
See the Lean proof and the cited items.
\paragraph{Blueprint notes.}
Used by C2, C4 step 3 and C5.a–c in place of the sheafified tensor product, whose sections are not computable. Setting: resolving data, X compact Hausdorff, an open U, the over-site of U. For a complex K of R|U-modules let F := 𝒜⁰ ⊗\textsubscript{R} K (S8 \texttt{sheafBC} along ρ), a complex of 𝒜⁰|U-modules, and T(K)\textsuperscript{n}(V) := degree-n part of 𝒜(V) ⊗\textsubscript{𝒜⁰(V)} ⊕\textsubscript{p} F\textsuperscript{p}(V) = ⊕\textsubscript{q} 𝒜\textsuperscript{q}(V) ⊗\textsubscript{𝒜⁰(V)} F\textsuperscript{n−q}(V), with plain tensor products of sections (A3's \texttt{Tot} over 𝒜(V), generators F\textsuperscript{p}(V)).

\begin{itemize}
\item \textbf{C7.a Sheaf.} T(K)\textsuperscript{n} is a sheaf of R|U-modules: 𝒜\textsuperscript{q}|U is a summand of a finite free 𝒜⁰|U-module (S4.b on X, restricted), so T(K)\textsuperscript{n} is sectionwise a retract of a finite power of ⊕\textsubscript{q} F\textsuperscript{n−q}, compatibly with restriction.
\item \textbf{C7.b Connection.} ∇ : F\textsuperscript{p} → 𝒜¹ ⊗\textsubscript{𝒜⁰} F\textsuperscript{p} (target a sheaf by C7.a) is the sheafification of α ⊗ k ↦ dα ⊗ η(1 ⊗ k) on the presheaf tensor 𝒜⁰(V) ⊗\textsubscript{R(V)} K\textsuperscript{p}(V), η the unit of sheafification (R-linear, well defined since d∘ρ = 0). ∇(αf) = dα ⊗ f + α∇f, ∇ commutes with d\textsubscript{F}, and ∇² = 0 (in 𝒜² ⊗ F): checked where f comes from the presheaf, then everywhere (η is locally surjective; the target is separated).
\item \textbf{C7.c Differential.} 𝕏\textsubscript{V}(a ⊗ f) = da ⊗ f + (−1)\textsuperscript{|a|}(a·∇f + a ⊗ d\textsubscript{F} f) (Koszul sign of v1.10; ∇ extended to 𝒜 ⊗ F by Leibniz). Leibniz holds by construction, 𝕏² = 0 from C7.b and c = 0, and 𝕏 commutes with restriction. Hence (F\textsuperscript{•}(V), 𝕏\textsubscript{V}) is a quasi-cohesive module N\textsubscript{K}(V) over 𝒜(V) (A3; generators not finitely generated), the restrictions are closed degree-0 maps semilinear along res, and T(K) is a complex of sheaves of R|U-modules.
\item \textbf{C7.d Comparison with tensorC.} The presheaf map 𝒜\textsuperscript{q}(V) ⊗\textsubscript{R(V)} K\textsuperscript{p}(V) → T(K)(V), a ⊗ k ↦ a ⊗ η(1 ⊗ k), is locally bijective (a retract of a finite power of η), so it induces an isomorphism tensorC(𝒜\textsuperscript{•}|U, K) ≅ T(K) of complexes (differentials compared on a ⊗ k). With C5 step 6 on any site: for K bounded above with flat terms, the map K → T(K), k ↦ 1 ⊗ η(1 ⊗ k), is a quasi-isomorphism.
\item \textbf{C7.e Comparison map.} For M cohesive over 𝒜(U) and e : Tot(M) → Tot(N\textsubscript{K}(U)) closed of degree 0, Ξ\textsubscript{e} : α\textsubscript{U}(M) → T(K) is on V the A9 extension (\texttt{extAlong}) of res\textsubscript{V} ∘ e; it is a map of complexes of sheaves (closedness, C7.c).
\item \textbf{C7.f Quasi-isomorphism criterion.} Filter α\textsubscript{U}(M) and T(K) by form degree (A6 \texttt{filTot}, sectionwise; subsheaves by C7.a and \texttt{degree\_local}), with short exact sequences sectionwise (A12 \texttt{filSC\_shortExact}) and A7 in sheaves. gr\textsuperscript{b} Ξ\textsubscript{e} is a retract of a finite power of gr\textsuperscript{0} Ξ\textsubscript{e} (the splitting of 𝒜\textsuperscript{b} over 𝒜⁰), so Ξ\textsubscript{e} is a quasi-isomorphism as soon as gr\textsuperscript{0} Ξ\textsubscript{e} : (V ↦ 𝒜⁰(V) ⊗\textsubscript{𝒜⁰(U)} Q\textsuperscript{•}) → F is one.
\end{itemize}
\end{proof}

% Generated by scripts/blueprint_gen.py; do not edit.
\chapter{The Dolbeault sheaf and the main theorem}\label{chap:F}
\textbf{Setting.} E a finite-dimensional complex normed space, n = dim\textsubscript{ℂ} E, \texttt{E : Type}. F a finite-dimensional complex normed space of values (F = ℂ for the main theorem; general F is for Z3 and Z8). M a complex manifold modelled on E (\texttt{ChartedSpace E M}, \texttt{IsManifold 𝓘(ℂ, E) ∞ M}), not necessarily compact. Only F = ℂ is needed on the critical path; F2–F4 and F7 are stated for general F because it is free there. Conventions §6 fixes all normalizations.

\begin{theorem}[M — Main theorem conditional]
\label{M}
\lean{Cohesive.MainStatementConditional, main_theorem_conditional}
\leanok
\uses{C1, C2, F10, F12, H1, H2, S9}
Let E be a finite-dimensional complex normed space and X a compact Hausdorff complex manifold modelled on E. Let A = 𝒜\textsubscript{X}(X) be the Dolbeault algebra of X and O\textsubscript{X} its sheaf of holomorphic functions. Assume that X satisfies the Dolbeault–Grothendieck lemma (H1) and the local gauge lemma in its weak form (H2). Then the functor J : Ho(P\textsubscript{A}) ⥤ D(O\textsubscript{X}), M ↦ α(M), into the derived category of sheaves of O\textsubscript{X}-modules is full and faithful, and its essential image is exactly D\textsubscript{perf}(X, O\textsubscript{X}). Equivalently, J corestricts to an equivalence Ho(P\textsubscript{A}) ≃ D\textsubscript{perf}(X, O\textsubscript{X}). Both hypotheses are theorems (H1, H2), so this form is a corollary of Z5.j.

Deviation: D\textsubscript{perf}(X, O\textsubscript{X}) instead of Block's D\textsuperscript{b\_}\{coh\}(X); the equivalence is stated as full, faithful and essential image D\textsubscript{perf}.
\par\smallskip{\footnotesize Source: blueprint.md §1.}
\end{theorem}
\begin{proof}
\leanok
\uses{F15}
See the Lean proof and the cited items.
\paragraph{Blueprint notes.}
For a finite-dimensional complex normed space E and a compact Hausdorff complex manifold X modelled on E, let A = 𝒜\textsubscript{X}(X) be the Dolbeault algebra (F10) and O\textsubscript{X} the sheaf of holomorphic functions (F12). Assume the two analytic hypotheses H1 (Dolbeault–Grothendieck) and H2 (local gauge lemma, weak form) on X. Then the functor

J : Ho(P\textsubscript{A}) ⥤ D\textsubscript{perf}(X, O\textsubscript{X}),  M ↦ α(M)

of C1 is an equivalence of categories.

Lean shape (names indicative, fixed in M0):

\begin{verbatim}
theorem blockPerf
    (E : Type) [NormedAddCommGroup E] [NormedSpace ℂ E] [FiniteDimensional ℂ E]
    (X : Type) [TopologicalSpace X] [T2Space X] [CompactSpace X]
    [ChartedSpace E X] [IsManifold 𝓘(ℂ, E) ∞ X]
    (h₁ : DolbeaultGrothendieck E X) (h₂ : LocalGaugeLemma E X) :
    (blockFunctor E X).Full ∧ (blockFunctor E X).Faithful ∧
      (blockFunctor E X).essImage = Perf (ringOf (holSheaf E X))
\end{verbatim}
Here \texttt{blockFunctor E X : HoP (dolbeault E X) ⥤ DerivedCategory (Mod (ringOf (holSheaf E X)))} is J with values in the whole derived category, and \texttt{Perf (ringOf (holSheaf E X))} is the \texttt{ObjectProperty} of S9.c (v1.9: \texttt{Perf} takes the ring sheaf; v1.10: \texttt{Perf} and \texttt{Mod} take a \texttt{RingSheaf X}, while \texttt{holSheaf E X} is a sheaf of \texttt{CommRingCat}, hence \texttt{ringOf}; \texttt{E} is explicit everywhere because the model space of \texttt{ChartedSpace} is not an \texttt{outParam}). \texttt{blockFunctor} and \texttt{Perf} use the same derived-category instance, \texttt{HasDerivedCategory.standard}, fixed once in \texttt{Statements.lean}. \warn{} **Why not \texttt{IsEquivalence} onto D\textsubscript{perf}.** That J lands in D\textsubscript{perf} is C2, which uses LL = H2, so a functor \texttt{Ho(P\_A) ⥤ D\_perf(X)} cannot be defined before H2 is assumed. Fully faithful with essential image exactly D\textsubscript{perf} is equivalent to "the corestriction is an equivalence"; that corollary is proved in \texttt{Checks.lean} (K6).

\warn{} \textbf{Smoothness exponents.} In Mathlib \texttt{⊤ : ℕ∞ω} is ω (analytic) and \texttt{∞} is C\textsuperscript{∞}. Complex manifolds use \texttt{∞} (complex C\textsuperscript{∞} transition maps, i.e. holomorphic); all forms and partitions of unity are real \texttt{∞}. Never write \texttt{⊤} for smoothness in this project.

\textbf{Relation to the sources.} This is [B, Thm 4.1.3] with D\textsuperscript{b\_coh} replaced by D\textsubscript{perf}, and [A, Thm 7.4.3] at h = 0 for compact X. The step D\textsubscript{perf}(X) = D\textsuperscript{b\_coh}(X) is Z1 (LATER).

\textbf{What is deliberately not here} (all in §6): D\textsuperscript{b\_coh} (Z1), triangulated structure (Z2), twisted sheaves and curvature (Z3), proofs of H1 and H2 (Z4, Z5), the dg-level quasi-equivalence (Z6), open manifolds (Z7), comparison with intrinsic forms (Z8).
\end{proof}

\begin{theorem}[F1 — Realification]
\label{F1}
\lean{Cohesive.isManifold_real, Cohesive.contMDiffRing_real_complex}
\leanok
Let E be a complex normed space and M a complex C\textsuperscript{∞} manifold modelled on E (charts with complex C\textsuperscript{∞}, i.e. holomorphic, transition maps). Then M, with the same charts, is a real C\textsuperscript{∞} manifold modelled on E (E regarded as a real normed space). Moreover ℂ, as a real manifold, is a smooth ring: multiplication ℂ × ℂ → ℂ is real C\textsuperscript{∞}. Hence the real-smooth ℂ-valued functions on M form a sheaf of commutative rings C\textsuperscript{∞\_M}(ℂ).
\par\smallskip{\footnotesize Source: blueprint.md §5.}
\end{theorem}
\begin{proof}
\leanok
See the Lean proof and the cited items.
\paragraph{Blueprint notes.}
\texttt{IsManifold 𝓘(ℝ, E) ∞ M} from \texttt{IsManifold 𝓘(ℂ, E) ∞ M}, same charts. Proof: transition maps are \texttt{ContDiffOn ℂ ∞}, hence \texttt{ContDiffOn ℝ ∞} (\texttt{ContDiffOn.restrict\_scalars} ✓). \texttt{NormedSpace ℝ E} comes from the instance \texttt{NormedSpace.complexToReal} ✓. The realification lemma itself is ✗ in Mathlib. Provide it as a theorem plus local instance; avoid a global instance (conventions §7).
\end{proof}

\begin{theorem}[F2 — The fibre Λ\textsuperscript{0,q}]
\label{F2}
\lean{Cohesive.Conj.map, Cohesive.form0Equiv, Cohesive.extToFibre, Cohesive.extToFibre_bijective, Cohesive.extEquivFibre, Cohesive.dzbarBasis}
\leanok
Let E and F be complex normed spaces. The conjugate space Ē is E with the same norm and the scalar action c · v := c̄ v. For q ∈ ℕ, the fibre Λ\textsuperscript{0,q}(E; F) is the normed space of continuous ℂ-multilinear alternating maps Ē\textsuperscript{q} → F (equivalently: real alternating maps on E that are antilinear in each slot).

\begin{itemize}
\item A continuous ℂ-linear map T : E → E′ is also a continuous ℂ-linear map Ē → Ē′.
\item Λ\textsuperscript{0,0}(E; F) ≅ F, isometrically: a 0-form at a point is a value.
\item If E is finite-dimensional, Φ\textsubscript{q} : ⋀\textsuperscript{q\_ℂ} (Ē)\textsuperscript{*} → Λ\textsuperscript{0,q}(E; ℂ), φ\textsubscript{1} ∧ ⋯ ∧ φ\textsubscript{q} ↦ (v ↦ det(φ\textsubscript{j}(v\textsubscript{i}))), is a linear isomorphism.
\item For a basis b of E indexed by a finite linearly ordered set ι, the forms dz̄\textsubscript{I} = Φ\textsubscript{q}(dz̄\textsubscript{i\_1} ∧ ⋯ ∧ dz̄\textsubscript{i\_q}) (dz̄\textsubscript{i} the dual basis of b in Ē), I running over the q-element subsets of ι, form a basis of Λ\textsuperscript{0,q}(E; ℂ). So dim Λ\textsuperscript{0,q}(E; ℂ) = C(n, q), n = dim\textsubscript{ℂ} E, and Λ\textsuperscript{0,q} = 0 for q > n.
\end{itemize}
\par\smallskip{\footnotesize Source: blueprint.md §5.}
\end{theorem}
\begin{proof}
\leanok
See the Lean proof and the cited items.
\paragraph{Blueprint notes.}
Ē = \texttt{Conj E}: E with the conjugate ℂ-action, a type synonym declared with \texttt{def} (not \texttt{abbrev}), normed like E. ✗ in Mathlib. A ℂ-linear map on Ē is an antilinear map on E, and a ℂ-linear map of E (for example DG(x) for holomorphic G) is also ℂ-linear on Ē.

Λ\textsuperscript{0,q}(E; F) := \texttt{Conj E [⋀\textasciicircum{}Fin q]→L[ℂ] F}, a ℂ-normed space (\texttt{instNormedAddCommGroup}, \texttt{instNormedSpace} ✓). \warn{} This replaces the earlier subtype of ℝ-alternating maps antilinear in each slot: the same maps, without membership side conditions, and Mathlib's exterior-power API applies to Ē directly. P2 confirmed \texttt{Conj E} works without instance diamonds; no fallback.

\begin{itemize}
\item \textbf{F2.a} Λ\textsuperscript{0,0}(E; F) ≅ F (\texttt{constOfIsEmpty} ✓).
\item \textbf{F2.b Exterior-algebra bridge.} Φ\textsubscript{q} : ⋀\textsuperscript{q\_ℂ}(Dual\textsubscript{ℂ} Ē) ≃ Λ\textsuperscript{0,q}(E; ℂ), φ₁ ∧ ⋯ ∧ φ\textsubscript{q} ↦ (v ↦ det(φ\textsubscript{i}(v\textsubscript{j}))), from \texttt{exteriorPower.pairingDual} ✓ (\texttt{pairingDual\_ιMulti\_ιMulti} ✓) and \texttt{exteriorPower.alternatingMapLinearEquiv} ✓, plus continuity in finite dimension. Bijective via the basis lemmas \texttt{ιMultiDual\_apply\_diag}, \texttt{ιMultiDual\_apply\_nondiag} ✓ and \texttt{Module.Basis.exteriorPower} ✓; packaging it as an equivalence is ✗.
\item \textbf{F2.c} Λ\textsuperscript{0,q} = 0 for q > n; dim = C(n, q)·dim F; a basis of E gives the basis dz̄\textsubscript{I}.
\end{itemize}
\end{proof}

\begin{theorem}[F3 — The ∂̄-operator on open subsets of E]
\label{F3}
\lean{Cohesive.antiPart, Cohesive.dbarPart, Cohesive.dbar, Cohesive.dbar_add, Cohesive.dbar_const_smul, Cohesive.dbar_congr}
\leanok
Let E, F be complex normed spaces and q ∈ ℕ. The antilinear part of a real-linear map L : E → G is L″(v) = ½(L v + i·L(i v)), a ℂ-linear map Ē → G. For α : E → Λ\textsuperscript{0,q}(E; F) and x ∈ E let D″α(x) be the antilinear part of the real derivative Dα(x), and ∂̄α(x)(v\textsubscript{0}, …, v\textsubscript{q}) = Σ\textsubscript{i} (−1)\textsuperscript{i} D″α(x)(v\textsubscript{i})(v\textsubscript{0}, …, v̂\textsubscript{i}, …, v\textsubscript{q}) ∈ Λ\textsuperscript{0,q+1}(E; F) (Mathlib's normalization of the exterior derivative). Properties:

\begin{itemize}
\item ∂̄(α + β)(x) = ∂̄α(x) + ∂̄β(x) and ∂̄(cα)(x) = c ∂̄α(x) (c ∈ ℂ) if α, β are real-differentiable at x;
\item ∂̄α(x) depends only on the germ of α at x.
\end{itemize}
Deviation: ∂̄ is defined for every function on E (with value 0 where α is not differentiable); by locality this gives ∂̄ on open subsets of E.
\par\smallskip{\footnotesize Source: blueprint.md §5.}
\end{theorem}
\begin{proof}
\leanok
\uses{F5}
See the Lean proof and the cited items.
\paragraph{Blueprint notes.}
For ω : E → Λ\textsuperscript{0,q}(E; F), smooth near x: D″ω(x) = ½ (Dω(x) + i·Dω(x) ∘ J), the antilinear part of the real derivative (J = multiplication by i), regarded as an element of \texttt{Conj E →L[ℂ] Λ\textasciicircum{}\{0,q\}(E; F)}; and ∂̄ω(x) = \texttt{alternatizeUncurryFin} (D″ω(x)) ✓, over ℂ on Ē.

\begin{itemize}
\item \textbf{F3.a} D″ω(x) is ℂ-linear on Ē (the defining computation).
\item \textbf{F3.b} ℂ-linear in ω, local (depends on the germ), smooth in, smooth out.
\item \textbf{F3.c} (sanity) ∂̄ω is the (0, q+1)-part of \texttt{extDeriv} ✓ ω.
\end{itemize}
\end{proof}

\begin{theorem}[F4 — ∂̄∂̄ = 0]
\label{F4}
\lean{Cohesive.dbar_dbar}
\leanok
Let E, F be complex normed spaces and α : E → Λ\textsuperscript{0,q}(E; F). If α is real C² at x, then ∂̄(∂̄α)(x) = 0.
\par\smallskip{\footnotesize Source: blueprint.md §5.}
\end{theorem}
\begin{proof}
\leanok
\uses{F3}
See the Lean proof and the cited items.
\paragraph{Blueprint notes.}
For ω C² near x, ∂̄(∂̄ω)(x) = 0. Proof: copy \texttt{extDerivWithin\_extDerivWithin\_apply} ✓ with D″ in place of D, over ℂ on Ē. The ℂ-bilinear map (u, v) ↦ D″D″ω(x)(u, v) on Ē is symmetric (from \texttt{ContDiffAt.isSymmSndFDerivAt} ✓ over ℝ); conclude with \texttt{alternatizeUncurryFin\_alternatizeUncurryFinCLM\_comp\_of\_symmetric} ✓.
\end{proof}

\begin{theorem}[F5 — Product on the total fibre]
\label{F5}
\lean{Cohesive.TotFibre.ofDeg, Cohesive.TotFibre.deg, Cohesive.TotFibre.mulCLM, Cohesive.TotFibre.extEquivTotFibre_mul, Cohesive.TotFibre.ofDeg_mul_ofDeg, Cohesive.TotFibre.ofDeg_mul_comm}
\leanok
Let E be a finite-dimensional complex normed space, n = dim\textsubscript{ℂ} E. The total fibre is the normed space Λ\textsuperscript{0,•}(E) = ∏\textsubscript{q=0}\textsuperscript{n} Λ\textsuperscript{0,q}(E; ℂ). It is a ℂ-algebra, with the product transported from the exterior algebra ⋀\textsubscript{ℂ} (Ē)\textsuperscript{*} along Φ = (Φ\textsubscript{q})\textsubscript{q} of F2, so that Φ is an isomorphism of ℂ-algebras. Viewing a q-form as an element of Λ\textsuperscript{0,•}(E) in degree q (zero if q > n):

\begin{itemize}
\item the product of a p-form and a q-form is a (p+q)-form;
\item graded commutativity: α ∧ β = (−1)\textsuperscript{pq} β ∧ α for a p-form α and a q-form β;
\item Laplace formula: for a 1-form φ and a q-form β, (φ ∧ β)(v\textsubscript{0}, …, v\textsubscript{q}) = Σ\textsubscript{i} (−1)\textsuperscript{i} φ(v\textsubscript{i}) β(v\textsubscript{0}, …, v̂\textsubscript{i}, …, v\textsubscript{q});
\item the product is a continuous bilinear map.
\end{itemize}
\par\smallskip{\footnotesize Source: blueprint.md §5.}
\end{theorem}
\begin{proof}
\leanok
\uses{F2}
See the Lean proof and the cited items.
\paragraph{Blueprint notes.}
The total fibre is Λ\textsuperscript{0,•}(E) := Π\textsubscript{q ≤ n} Λ\textsuperscript{0,q}(E; ℂ), a finite product of normed spaces. Its multiplication is transported along Φ = (Φ\textsubscript{q})\textsubscript{q} from \texttt{ExteriorAlgebra ℂ (Dual\_ℂ Ē)}, which is graded (\texttt{ExteriorAlgebra.gradedAlgebra} ✓) and graded-commutative on generators (\texttt{ιMulti\_mul\_ιMulti\_anticomm} ✓). Associativity, unit and graded commutativity come for free; continuity by finite dimension. \warn{} This replaces a hand-built shuffle product, whose associativity is not in Mathlib (only \texttt{AlternatingMap.domCoprod} ✓ and its alternatization lemmas are).

\begin{itemize}
\item \textbf{F5.a Laplace formula.} For a 1-form φ and a q-form ω, φ ∧ ω = \texttt{alternatizeUncurryFin} (v ↦ φ(v) • ω), i.e. (φ ∧ ω)(v₀, …, v\textsubscript{q}) = Σ\textsubscript{i} (−1)\textsuperscript{i} φ(v\textsubscript{i}) ω(…, v̂\textsubscript{i}, …). From expansion of the determinant along the first row (\texttt{Matrix.det\_succ\_row\_zero} ✓). This is the only link between the product and the analysis.
\item \textbf{F5.b} F-valued forms are a module over Λ\textsuperscript{0,•}(E) via Λ\textsuperscript{0,q}(E; F) ≅ Λ\textsuperscript{0,q}(E; ℂ) ⊗ F (later only).
\end{itemize}
\end{proof}

\begin{theorem}[F6 — Leibniz rule]
\label{F6}
\lean{Cohesive.wedge, Cohesive.dbar_wedge}
\leanok
\uses{F2, F5}
Let E be a finite-dimensional complex normed space. For a p-form α and a q-form β at a point, α ∧ β ∈ Λ\textsuperscript{0,p+q}(E; ℂ) is the degree-(p+q) part of their product in Λ\textsuperscript{0,•}(E). If α : E → Λ\textsuperscript{0,p}(E; ℂ) and β : E → Λ\textsuperscript{0,q}(E; ℂ) are real-differentiable at x, then in Λ\textsuperscript{0,•}(E) ∂̄(α ∧ β)(x) = ∂̄α(x) ∧ β(x) + (−1)\textsuperscript{p} α(x) ∧ ∂̄β(x).
\par\smallskip{\footnotesize Source: blueprint.md §5.}
\end{theorem}
\begin{proof}
\leanok
\uses{F3}
See the Lean proof and the cited items.
\paragraph{Blueprint notes.}
∂̄(α ∧ β) = ∂̄α ∧ β + (−1)\textsuperscript{p} α ∧ ∂̄β. Proof: (i) ∂̄(f·η) = ∂̄f ∧ η + f·∂̄η for a function f, by the product rule and F5.a; (ii) constant forms are ∂̄-closed; (iii) locally every form is Σ\textsubscript{I} f\textsubscript{I} dz̄\textsubscript{I} (F2.c); expand and use graded commutativity. No general alternation identity is needed.
\end{proof}

\begin{theorem}[F7 — Holomorphic pullback]
\label{F7}
\lean{Cohesive.TotFibre.pull, Cohesive.pullback, Cohesive.dbar_pullback, Cohesive.pullback_smul, Cohesive.TotFibre.pull_mul, Cohesive.TotFibre.pull_one}
\leanok
Let E, F be complex normed spaces. For G : E → E and α : E → Λ\textsuperscript{0,q}(E; F) the pullback is G\textsuperscript{*α}(x) = α(G x) ∘ DG(x), with the complex derivative DG(x) acting on Ē (junk where G is not complex-differentiable).

\begin{itemize}
\item If α is real-differentiable at G x and G is complex C² at x, then ∂̄(G\textsuperscript{\emph{α})(x) = G\textsuperscript{}}(∂̄α)(x).
\item G\textsuperscript{*}(f α)(x) = f(G x) · G\textsuperscript{*α}(x) for f : E → ℂ.
\item (G ∘ H)\textsuperscript{\emph{} α (x) = H\textsuperscript{}}(G\textsuperscript{*α})(x) if H is complex-differentiable at x and G at H x.
\item If E is finite-dimensional, the pullback T\textsuperscript{*} along a continuous ℂ-linear T : E → E, applied degreewise on Λ\textsuperscript{0,•}(E), is a unital ring homomorphism.
\end{itemize}
Deviation: G is a map E → E rather than a holomorphic map between open subsets; since all hypotheses are local at x, this covers the latter.
\par\smallskip{\footnotesize Source: blueprint.md §5.}
\end{theorem}
\begin{proof}
\leanok
\uses{F3}
See the Lean proof and the cited items.
\paragraph{Blueprint notes.}
For G : V → V′ holomorphic between open subsets of E, G\textsuperscript{*ω}(x) = ω(G x) ∘ DG(x) (\texttt{compContinuousLinearMap} ✓), with DG(x) regarded as a ℂ-linear map of Ē.

\begin{itemize}
\item \textbf{F7.a} G\textsuperscript{\emph{} ∂̄ = ∂̄ G\textsuperscript{}}. Proof: copy \texttt{extDerivWithin\_pullback} ✓ with D″. The second-derivative term is antilinear in the new slot, so D″ keeps it, and alternation kills it (\texttt{alternatizeUncurryFin\_fderivCompContinuousLinearMap\_eq\_zero} ✓).
\item \textbf{F7.b} On fibres G\textsuperscript{*} is the algebra map induced by the dual of DG(x) on the exterior algebra (\texttt{ExteriorAlgebra.map}), so it respects products; it respects multiplication by functions.
\item \textbf{F7.c} (G ∘ H)\textsuperscript{\emph{} = H\textsuperscript{}} G\textsuperscript{\emph{}, id\textsuperscript{}} = id.
\end{itemize}
\end{proof}

\begin{theorem}[F8 — Cauchy–Riemann]
\label{F8}
\lean{Cohesive.dbar_eq_zero_iff_contDiffOn}
\leanok
Let E, F be complex normed spaces, V ⊆ E open and f : E → F real C\textsuperscript{∞} on V. Then ∂̄f = 0 at every point of V (f regarded as a 0-form) if and only if f is complex C\textsuperscript{∞} (holomorphic) on V.
\par\smallskip{\footnotesize Source: blueprint.md §5.}
\end{theorem}
\begin{proof}
\leanok
\uses{F3}
See the Lean proof and the cited items.
\paragraph{Blueprint notes.}
For f : E → F real C\textsuperscript{∞} on an open V: ∂̄f = 0 on V ⇔ f is \texttt{ContDiffOn ℂ ∞} on V. Proof (⇒): ∂̄f = 0 means Df(x) is ℂ-linear (\texttt{differentiableAt\_iff\_restrictScalars} ✓). A symmetric ℝ-bilinear map that is ℂ-linear in one slot is ℂ-linear in both, so the derivative of fderiv ℂ f is again ℂ-linear; induct. \warn{} This avoids several-variable Cauchy or Osgood theory, which Mathlib does not have (✗ "Osgood", "Hartogs").
\end{proof}

\begin{theorem}[F9 — The sheaves 𝒜\textsuperscript{0,q}\textsubscript{M}(F) and 𝒜\textsuperscript{0,•}\textsubscript{M}]
\label{F9}
\lean{Cohesive.IsDolbeault, Cohesive.IsOfDegree.dbarSec, Cohesive.restrictForm, Cohesive.dolbeaultLocalPredicate, Cohesive.IsDolbeault.contDiffAt_chartExpr}
\leanok
\uses{F7}
Let E be a finite-dimensional complex normed space and X a complex manifold modelled on E. A form on an open U ⊆ X is a function α : U → Λ\textsuperscript{0,•}(E), the value α(y) being written in the chart at y. Its expression in a chart e is z ↦ T\textsuperscript{\emph{α}(y), y = e\textsuperscript{−1}(z), T the complex derivative at z of the transition map chartAt(y) ∘ e\textsuperscript{−1}. α is a *Dolbeault form} if, for every x ∈ U, its expression in φ = chartAt(x) is real C\textsuperscript{∞} at φ(x); its ∂̄ is (∂̄α)(x) = ∂̄(expression of α in φ)(φ(x)), degree by degree. A form has degree q if all its other components vanish.

\begin{itemize}
\item Restrictions of Dolbeault forms are Dolbeault forms, and being Dolbeault is local, so Dolbeault forms are a sheaf of sets on X.
\item The expression of a Dolbeault form in any atlas chart e is real C\textsuperscript{∞} at e(x) for x ∈ U ∩ e.source.
\item ∂̄ of a Dolbeault form is a Dolbeault form, and ∂̄ maps forms of degree q to degree q + 1.
\end{itemize}
Deviation: forms take values in the total fibre, and only ℂ-valued forms are treated; 𝒜\textsuperscript{0,q} is the degree-q part (equivalent, since the direct sum is finite).
\par\smallskip{\footnotesize Source: blueprint.md §5.}
\end{theorem}
\begin{proof}
\leanok
\uses{F4, F6}
See the Lean proof and the cited items.
\paragraph{Blueprint notes.}
Sections over an open U are functions ω : U → Λ\textsuperscript{0,q}(E; F) (or into the total fibre Λ\textsuperscript{0,•}(E)), ω(x) read in the chart \texttt{chartAt x}, satisfying the local predicate: for each x ∈ U, with φ = \texttt{chartAt x}, the chart expression z ↦ ω(y) ∘ D(\texttt{chartAt y} ∘ φ⁻¹)(z) at z = φ(y) is smooth near φ(x). Built with \texttt{TopCat.LocalPredicate} ✓ and \texttt{TopCat.subsheafToTypes} ✓, as Mathlib builds \texttt{smoothSheaf} ✓; the ring structure on 𝒜\textsuperscript{0,•} follows \texttt{smoothSheafCommRing} ✓. Making each 𝒜\textsuperscript{0,q} a \texttt{SheafOfModules} over 𝒜\textsuperscript{0} has no Mathlib precedent and is part of F9's size.

\begin{itemize}
\item \textbf{F9.a} The predicate is independent of the chart choices (F7). The transition factor is the coordinate change of \texttt{tangentBundleCore} ✓ (\texttt{tangentBundleCore\_coordChange} ✓), so its cocycle lemmas apply.
\item \textbf{F9.b} ∂̄ chartwise: (∂̄ω)(x) = (∂̄ω\textsubscript{φ})(φ(x)), φ = \texttt{chartAt x}; well defined by F7.a.
\item \textbf{F9.c} ∧ and the module structure over smooth functions are pointwise.
\item \textbf{F9.d} 𝒜\textsuperscript{0,0}\textsubscript{M}(ℂ) ≅ C\textsuperscript{∞\_M}(ℂ) := \texttt{smoothSheafCommRing 𝓘(ℝ, E) 𝓘(ℝ, ℂ) M ℂ} ✓.
\item \textbf{F9.e} (sanity) For an open V ⊆ E with its induced structure (\texttt{TopologicalSpace.Opens.instChartedSpace} ✓), sections over V are just smooth functions V → Λ\textsuperscript{0,q}(E; F).
\end{itemize}
\end{proof}

\begin{theorem}[F10 — The Dolbeault dga]
\label{F10}
\lean{Cohesive.dolbeault0, Cohesive.dolbeaultSubalgebra, Cohesive.dbarL, Cohesive.dolbeaultRes, Cohesive.dolbeault_glue, Cohesive.isOfDegree_of_locally}
\leanok
\uses{A1, F5, F9}
Let E be a finite-dimensional complex normed space, n = dim\textsubscript{ℂ} E, and X a complex manifold modelled on E. For an open U ⊆ X, 𝒜\textsubscript{X}(U) is the ℂ-algebra of Dolbeault forms on U (F9) with pointwise operations, graded by 𝒜\textsuperscript{0,q}\textsubscript{X}(U) = the forms of degree q (zero for q < 0 and q > n).

\begin{itemize}
\item 𝒜\textsubscript{X}(U) with d = ∂̄ is a bounded graded-commutative dga (A1) with top degree n and curvature 0.
\item Restriction along V ⊆ U is a map of dgas 𝒜\textsubscript{X}(U) → 𝒜\textsubscript{X}(V).
\item Sheaf conditions: Dolbeault forms on opens U\textsubscript{i} that agree on overlaps glue to a unique Dolbeault form on ⋃ U\textsubscript{i}; having degree q is a local property.
\item U ↦ 𝒜\textsuperscript{0,0}\textsubscript{X}(U) is a sheaf 𝒜\textsuperscript{0} of commutative rings on X.
\end{itemize}
The Dolbeault algebra of X is A = 𝒜\textsubscript{X}(X).
\par\smallskip{\footnotesize Source: blueprint.md §5.}
\end{theorem}
\begin{proof}
\leanok
\uses{A9}
See the Lean proof and the cited items.
\paragraph{Blueprint notes.}
𝒜\textsubscript{M}(U) := sections over U of the total-fibre sheaf 𝒜\textsuperscript{0,•}\textsubscript{M} (F9), with pointwise product (F5) and ∂̄ (F9.b); internally graded by the components q, and 𝒜\textsuperscript{0,q}\textsubscript{M} is the q-th component sheaf. A bounded graded-commutative dga (A1) for every open U, with restrictions dga maps. (v1.9) The component sheaves are \emph{defined} as the subtypes 𝒜\textsuperscript{0,q}\textsubscript{M}(U) := ↥(𝒜\textsubscript{M}(U) q) of the total section algebra, so that the 𝒜\textsuperscript{0} of C0 is the dga's own \texttt{𝒜 0} (C0, one carrier for A⁰). \textbf{The Dolbeault algebra of X is A = 𝒜\textsubscript{X}(X).} \warn{} The total fibre removes the direct-sum construction of conventions §5: the ring axioms hold pointwise.
\end{proof}

\begin{theorem}[F11 — Fineness and local freeness]
\label{F11}
\lean{Cohesive.isFine_dolbeault0, Cohesive.dolbeault_locallyFree}
\leanok
\uses{F2}
Let E be a finite-dimensional complex normed space, n = dim\textsubscript{ℂ} E, and X a complex manifold modelled on E.

\begin{itemize}
\item If X is compact Hausdorff, 𝒜\textsuperscript{0} is fine: for every finite open cover (U\textsubscript{i}) of X there are global sections ρ\textsubscript{i} ∈ 𝒜\textsuperscript{0,0}\textsubscript{X}(X) with Σ\textsubscript{i} ρ\textsubscript{i} = 1, each ρ\textsubscript{i} supported in U\textsubscript{i} (it vanishes on an open V\textsubscript{i} with U\textsubscript{i} ∪ V\textsubscript{i} = X).
\item 𝒜\textsuperscript{0,q} is locally free of rank C(n, q) over 𝒜\textsuperscript{0}: every x ∈ X has an open neighbourhood U with forms e\textsubscript{I} ∈ 𝒜\textsuperscript{0,q}\textsubscript{X}(U), I running over the q-element subsets of \{0, …, n−1\}, such that on every open V ⊆ U each β ∈ 𝒜\textsuperscript{0,q}\textsubscript{X}(V) is uniquely Σ\textsubscript{I} f\textsubscript{I} · e\textsubscript{I}|\textsubscript{V} with f\textsubscript{I} ∈ 𝒜\textsuperscript{0,0}\textsubscript{X}(V).
\end{itemize}
\par\smallskip{\footnotesize Source: blueprint.md §5.}
\end{theorem}
\begin{proof}
\leanok
\uses{F1, F12, S1}
See the Lean proof and the cited items.
\paragraph{Blueprint notes.}
For compact X: 𝒜\textsuperscript{0} is fine (\texttt{SmoothPartitionOfUnity.exists\_isSubordinate} ✓ for 𝓘(ℝ, E); needs T2 and σ-compact); 𝒜\textsuperscript{0,q} is locally free of rank C(n, q) over 𝒜\textsuperscript{0} (F2.c in charts); 𝒜\textsuperscript{0,q} = 0 for q > n. (v1.9) \texttt{IsFine} asks for the stalk-free \texttt{SupportedIn} (∃ V, U ⊔ V = ⊤ ∧ ρ|V = 0): take V := (tsupport ρ\textsubscript{i})ᶜ; the real-valued ρ\textsubscript{i} become sections of 𝒜\textsuperscript{0} through ℝ → ℂ.
\end{proof}

\begin{theorem}[F12 — Holomorphic functions]
\label{F12}
\lean{Cohesive.dolbeaultRhoU, Cohesive.dolbeaultRho, Cohesive.dolbeaultRhoU_injective, Cohesive.dbar_dolbeaultRhoU, Cohesive.exists_dolbeaultRhoU_of_dbar_eq_zero, Cohesive.holSheaf}
\leanok
\uses{F8}
Let E be a finite-dimensional complex normed space and X a complex manifold modelled on E. O\textsubscript{X} is the sheaf of commutative rings of complex C\textsuperscript{∞} (holomorphic) functions on X. ρ : O\textsubscript{X} → 𝒜\textsuperscript{0} is the morphism of sheaves of rings that sends a holomorphic function f on U to f as a form of degree 0. For every open U ⊆ X:

\begin{itemize}
\item ρ\textsubscript{U} is injective;
\item ∂̄ ρ\textsubscript{U}(f) = 0 for every f ∈ O\textsubscript{X}(U);
\item every α ∈ 𝒜\textsuperscript{0,0}\textsubscript{X}(U) with ∂̄α = 0 is ρ\textsubscript{U}(f) for some f ∈ O\textsubscript{X}(U).
\end{itemize}
So 0 → O\textsubscript{X}(U) → 𝒜\textsuperscript{0,0}\textsubscript{X}(U) → 𝒜\textsuperscript{0,1}\textsubscript{X}(U) is exact.
\par\smallskip{\footnotesize Source: blueprint.md §5.}
\end{theorem}
\begin{proof}
\leanok
\uses{F10}
See the Lean proof and the cited items.
\paragraph{Blueprint notes.}
O\textsubscript{M} := \texttt{smoothSheafCommRing 𝓘(ℂ, E) 𝓘(ℂ, ℂ) M ℂ} ✓ (complex C\textsuperscript{∞} = holomorphic). ρ : O\textsubscript{M} → C\textsuperscript{∞\_M}(ℂ) is the inclusion (F1 and restriction of scalars); it is injective with image ker(∂̄ : 𝒜\textsuperscript{0} → 𝒜\textsuperscript{1}) by F8. So 0 → O → 𝒜\textsuperscript{0} → 𝒜\textsuperscript{1} is exact.
\end{proof}

\begin{theorem}[H1 — Dolbeault–Grothendieck lemma]
\label{H1}
\lean{Cohesive.DolbeaultGrothendieck, Cohesive.dolbeaultGrothendieck}
\leanok
Let E be a finite-dimensional complex normed space and X a complex manifold modelled on E. X satisfies the \emph{Dolbeault–Grothendieck lemma} if for every open U ⊆ X, every q ∈ ℕ, every x ∈ U and every α ∈ 𝒜\textsuperscript{0,q+1}\textsubscript{X}(U) with ∂̄α = 0, there are an open U′ with x ∈ U′ ⊆ U and η ∈ 𝒜\textsuperscript{0,q}\textsubscript{X}(U′) with ∂̄η = α|\textsubscript{U′}. Theorem: this holds for every such X (no compactness or Hausdorff assumption). With F12 it says that ρ : O\textsubscript{X} → 𝒜\textsuperscript{0,•}\textsubscript{X} is a resolution.

Deviation: stated on the manifold X in sheaf form, not on polydiscs in ℂ\textsuperscript{n}.
\par\smallskip{\footnotesize Source: blueprint.md §5.}
\end{theorem}
\begin{proof}
\leanok
\uses{F10}
See the Lean proof and the cited items.
\paragraph{Blueprint notes.}
\texttt{DolbeaultGrothendieck E X : Prop}: for every open U ⊆ X, q ∈ ℕ, x ∈ U and ω ∈ Γ(U, 𝒜\textsuperscript{0,q+1}\textsubscript{X}) with ∂̄ω = 0, there are an open x ∈ U′ ⊆ U and η ∈ Γ(U′, 𝒜\textsuperscript{0,q}\textsubscript{X}) with ∂̄η = ω|U′. Together with F12 this says exactly that ρ : O\textsubscript{X} → 𝒜\textsuperscript{•\_X} is a resolution.

\begin{itemize}
\item True on every complex manifold (Dolbeault–Grothendieck lemma). Proof is Z4.
\item Trivial cases to prove as lemmas: q + 1 > n; n = 0.
\item \warn{} Stated on X, not on open subsets of E: this keeps chart transport (F13) off the critical path, and it is the standard sheaf-theoretic form of the lemma.
\end{itemize}
\end{proof}

\begin{theorem}[H2 — Local gauge lemma, weak form]
\label{H2}
\lean{Cohesive.LocalGaugeLemma, Cohesive.localGaugeLemma}
\leanok
\uses{A10, A9}
Let E be a finite-dimensional complex normed space and X a complex manifold modelled on E. X satisfies the \emph{local gauge lemma} (weak form) if for every cohesive module M over A = 𝒜\textsubscript{X}(X) and every x ∈ X there are an open U′ ∋ x, a bounded complex W of finitely generated projective O\textsubscript{X}(U′)-modules, and an isomorphism in Ho(P\textsubscript{𝒜\_X(U′)}) between the restriction M|\textsubscript{U′} (base change along 𝒜\textsubscript{X}(X) → 𝒜\textsubscript{X}(U′), A9) and 𝒜\textsubscript{X}(U′) ⊗\textsubscript{O\_X(U′)} W (A10, along ρ\textsubscript{U′}). Theorem: this holds for every such X.

Deviation: the weak, sectionwise form of [B, Lemma 4.1.5]: only restrictions of global cohesive modules, a homotopy equivalence instead of a gauge equivalence, and W need not be free.
\par\smallskip{\footnotesize Source: blueprint.md §5.}
\end{theorem}
\begin{proof}
\leanok
\uses{F12}
See the Lean proof and the cited items.
\paragraph{Blueprint notes.}
\texttt{LocalGaugeLemma E X : Prop}: for every cohesive module M over A = 𝒜\textsubscript{X}(X) and x ∈ X, there are an open x ∈ U′, a bounded complex W of finitely generated projective O\textsubscript{X}(U′)-modules, and a homotopy equivalence in P\textsubscript{𝒜\_X(U′)} between M|U′ (A9) and 𝒜\textsubscript{X}(U′) ⊗\textsubscript{O\_X(U′)} W (A10).

\begin{itemize}
\item True: the proof of [B, Lemma 4.1.5] (whose statement is full faithfulness of α) gives a gauge equivalence on a smaller polydisc with W free; this is its weak form, and [A, Assumption 6] at the level of sections. It is field LL of C0 verbatim. Proof is Z5.
\item Trivial case to prove: n = 0.
\item (v1.9, audit F9) Lean form, to be reviewed at gate 3: \texttt{LocalGaugeLemma E X} says: for every \texttt{M : CohesiveModule (dolbeault E X)} and x, there are an open U′ ∋ x, a complex \texttt{W : CochainComplex (ModuleCat.\{0\} (O(U′))) ℤ} with \texttt{IsBoundedFGProj}, and an isomorphism in \texttt{HoP} of the U′-dga between \texttt{M.baseChange (res\_\{U′\})} (A9) and \texttt{CohesiveModule.fromComplex ρ\_\{U′\} … W} (A10, with \texttt{hc : c = 0} and \texttt{d ∘ ρ\_\{U′\} = 0}, both facts about the Dolbeault dga). ρ\textsubscript{U′} : O(U′) →+* ↥(𝒜\textsubscript{U′} 0) (C0, one carrier).
\end{itemize}
\end{proof}

\begin{definition}[F14 — The fields of C0]
\label{F14}
\lean{Cohesive.dolbeaultPresheaf, Cohesive.dolbeaultPreData, Cohesive.dolbeaultData}
\leanok
\uses{C0, F10, F11, F12, H1, H2}
Let E be a finite-dimensional complex normed space and X a complex manifold modelled on E.

\begin{itemize}
\item U ↦ 𝒜\textsubscript{X}(U) with restriction of forms is a presheaf of dgas on X.
\item If X is compact Hausdorff, the Dolbeault data form pre-resolving data (C0.a): the dgas 𝒜\textsubscript{X}(U), which glue and have curvature 0; O = O\textsubscript{X} with ρ\textsubscript{U} : O\textsubscript{X}(U) → 𝒜\textsuperscript{0,0}\textsubscript{X}(U) and ∂̄ ∘ ρ = 0 (F12); 𝒜\textsuperscript{0} a fine sheaf and each 𝒜\textsuperscript{0,q} locally free over it (F11). No analytic hypothesis is used.
\item If moreover X satisfies H1 and H2, they form resolving data (C0.b): ρ is injective, O\textsubscript{X} → 𝒜\textsuperscript{0} → 𝒜\textsuperscript{1} is exact (F12), the Dolbeault complex is locally exact in degrees ≥ 1 (H1), and the local lifting field LL is H2.
\end{itemize}
\par\smallskip{\footnotesize Source: blueprint.md §5.}
\end{definition}
\paragraph{Blueprint notes.}
F10, F11 and F12 give the pre-resolving data C0.a for X, unconditionally. Given H1 and H2, F12 and H1 give the resolution field and H2 is the field LL, completing C0.b.

\begin{theorem}[F15 — Main theorem conditional: proof]
\label{F15}
\lean{Cohesive.blockFunctor, Cohesive.blockPerf}
\leanok
\uses{C0, C6, F14, H1, H2}
The proof of M, as the Lean theorem \texttt{blockPerf}. Let E be a finite-dimensional complex normed space and X a compact Hausdorff complex manifold modelled on E, A = 𝒜\textsubscript{X}(X). The Dolbeault data of X are pre-resolving data (F14) without any analytic hypothesis, so C1 gives the functor J : Ho(P\textsubscript{A}) ⥤ D(O\textsubscript{X}). If X satisfies H1 and H2, the Dolbeault data are resolving data (F14), and the comparison theorem C6 specialized to them says that J is full and faithful with essential image D\textsubscript{perf}(X, O\textsubscript{X}).

This is the only node that applies the abstract theorem C6 to the Dolbeault data; M is the statement itself (\texttt{MainStatementConditional}), and Z5.j is its form without H1 and H2.

Deviation: proved for any model of the derived category (any \texttt{HasDerivedCategory} instance), not only the standard one.
\par\smallskip{\footnotesize Source: blueprint.md §5.}
\end{theorem}
\begin{proof}
\leanok
See the Lean proof and the cited items.
\paragraph{Blueprint notes.}
\texttt{blockFunctor E X := J} for the C0.a data of F14 (no hypotheses). Assemble C0.b from H1, H2 (F14) and apply C6. This proves M.
\end{proof}

% Generated by scripts/blueprint_gen.py; do not edit.
\chapter{Proof of H1: the Dolbeault–Grothendieck lemma}\label{chap:Z4}
The plan proves the lemma first on open subsets of ℂⁿ by Grothendieck's induction, then transfers it to X through holomorphic charts. The analytic core reduces to one parametric Cauchy transform, built from TauCeti's Cauchy–Pompeiu file and Mathlib's parametric convolution lemmas.

\textbf{Target.} For every open U ⊆ X, q ∈ ℕ, x ∈ U and ω ∈ Γ(U, 𝒜\textsuperscript{0,q+1}) with ∂̄ω = 0, there are an open U′ with x ∈ U′ ⊆ U and η ∈ Γ(U′, 𝒜\textsuperscript{0,q}) with ∂̄η = ω|U′.

\textbf{Four layers.}

\begin{enumerate}
\item \emph{Parametric ∂̄-solver in one variable} (pure analysis, Banach-valued). Solves ∂g/∂z̄\textsubscript{k} = f near x, keeps g holomorphic in every variable in which f is, and commutes with continuous linear maps.
\item \emph{Flat Dolbeault complex on ℂⁿ}. Forms are maps into Mathlib's \texttt{ℂⁿ [⋀\textasciicircum{}Fin q]→L[ℝ] ℂ}; ∂̄ω = Σ\textsubscript{j} dz̄\textsubscript{j} ∧ ∂̄\textsubscript{j} ω. Contains the induction and the flat lemma.
\item \emph{Comparison with TauCeti}. Identifies TauCeti's \texttt{delBar} on open subsets of ℂⁿ with the flat ∂̄.
\item \emph{Transfer to X} through charts and \texttt{mpullback\_delBar}, then instantiates \texttt{DolbeaultGrothendieck E X}.
\end{enumerate}
\textbf{Two simplifications drive the design.} The statement is local, so each induction step may shrink the neighbourhood; Hörmander's exhaustion argument for a whole polydisc is not needed. The sup norm on \texttt{Fin n → ℂ} makes \texttt{Metric.ball x r} a polydisc, so no polydisc API is needed.

\textbf{Upstream route.} TauCeti's roadmap ComplexKahlerGeometry, stage K2, lists \texttt{dolbeault-grothendieck-lemma} as a target. As of 2026-10-29 K2 still needs local exactness and the parameter-dependent solver (\href{https://github.com/TauCetiProject/TauCeti/pull/22585}{\#22585}). Layers 1–3 map onto exactly these items; layer 4 is your project's glue.

The proof needs five algebraic identities, five calculus lemmas, one solver with four properties, and a one-line induction step. Indices here are 1-based.

\textbf{Notation.} e\textsubscript{j} is the j-th unit vector of ℂⁿ and dz̄\textsubscript{j}(v) = conj(v\textsubscript{j}). For f with values in a complex Banach space, ∂̄\textsubscript{j} f = ½(D\textsubscript{e\_j} f + i·D\textsubscript{i·e\_j} f), where D\textsubscript{v} is the real directional derivative. The contraction is ι\textsubscript{v} a = a(v, ·). For a real 1-form φ with complex values:

\$\$ (\textbackslash{}varphi\textbackslash{}wedge a)(v\textsubscript{0},\textbackslash{}dots,v\textsubscript{q})=\textbackslash{}sum\textsubscript{i=0}\textsuperscript{q}(-1)\textsuperscript{i}\textbackslash{},\textbackslash{}varphi(v\textsubscript{i})\textbackslash{},a(v\textsubscript{0},\textbackslash{}dots,\textbackslash{}widehat\{v\textsubscript{i}\},\textbackslash{}dots,v\textsubscript{q}) \$\$

The flat Dolbeault operator is

\$\$ \textbackslash{}bar\textbackslash{}partial\textbackslash{}omega=\textbackslash{}sum\textsubscript{j=1}\textsuperscript{n} d\textbackslash{}bar z\textsubscript{j}\textbackslash{}wedge\textbackslash{}bar\textbackslash{}partial\textsubscript{j}\textbackslash{}omega \$\$

\paragraph{A. Pointwise algebra}
\begin{itemize}
\item \textbf{A1 (contraction of a wedge).} ι\textsubscript{v}(φ ∧ a) = φ(v)·a − φ ∧ ι\textsubscript{v} a. Direct from the formula above.
\item \textbf{A2.} ι\textsubscript{v} ι\textsubscript{v} = 0 and ι\textsubscript{v} ι\textsubscript{w} = −ι\textsubscript{w} ι\textsubscript{v}. Immediate from alternation.
\item \textbf{A3.} φ ∧ φ ∧ a = 0 and φ ∧ ψ ∧ a = −ψ ∧ φ ∧ a.
\item \textbf{A4 (type).} If a is of type (0,q), so are ι\textsubscript{v} a and dz̄\textsubscript{j} ∧ a (types (0,q−1) and (0,q+1)). dz̄\textsubscript{j} is conjugate-linear, and the defining relation is linear in each slot.
\item \textbf{A5 (expansion).} If a has type (0,q+1), then a(v, w) = Σ\textsubscript{j} conj(v\textsubscript{j})·(ι\textsubscript{e\_j} a)(w). Write v = Σ\textsubscript{j} (Re v\textsubscript{j})e\textsubscript{j} + (Im v\textsubscript{j})(i e\textsubscript{j}) and use a(i e\textsubscript{j}, ·) = −i·a(e\textsubscript{j}, ·). Corollary: if ι\textsubscript{e\_j} a = 0 for every j, then a = 0.
\end{itemize}
\paragraph{B. Calculus on ℂⁿ}
Let U ⊆ ℂⁿ be open and ω smooth on U.

\begin{itemize}
\item \textbf{B1 (naturality).} For a continuous ℂ-linear L, ∂̄\textsubscript{j}(L ∘ ω) = L ∘ ∂̄\textsubscript{j} ω. In particular ∂̄\textsubscript{j} commutes with ι\textsubscript{v}, with dz̄\textsubscript{i} ∧ (·), and with every map defining type.
\item \textbf{B2.} ∂̄\textsubscript{i} ∂̄\textsubscript{j} ω = ∂̄\textsubscript{j} ∂̄\textsubscript{i} ω on U. Use symmetry of the second real derivative.
\item \textbf{B3 (∂̄² = 0).} Expand ∂̄∂̄ω = Σ\textsubscript{i,j} dz̄\textsubscript{i} ∧ dz̄\textsubscript{j} ∧ ∂̄\textsubscript{i}∂̄\textsubscript{j} ω with B1. The coefficient is symmetric in (i,j) by B2 and the prefix antisymmetric by A3, so the sum equals its own negative.
\item \textbf{B4 (contraction of ∂̄).} ι\textsubscript{e\_m} ∂̄ω = ∂̄\textsubscript{m} ω − ∂̄(ι\textsubscript{e\_m} ω). Apply A1 termwise with dz̄\textsubscript{j}(e\textsubscript{m}) = δ\textsubscript{jm}, then B1.
\item \textbf{B5 (closed ⇒ holomorphic in the tail).} Suppose ω involves only dz̄\textsubscript{1}, …, dz̄\textsubscript{k} on U and ∂̄ω = 0 on U. Then ∂̄\textsubscript{m} ω = 0 on U for every m > k. Indeed ι\textsubscript{e\_m} ω ≡ 0 on the open set U, so its ∂̄ vanishes there, and B4 gives ∂̄\textsubscript{m} ω = ι\textsubscript{e\_m} ∂̄ω = 0.
\item \textbf{B6 (slices).} ∂̄\textsubscript{k} F(z) is the ∂̄-derivative at z\textsubscript{k} of w ↦ F(z with z\textsubscript{k} replaced by w). Chain rule along \texttt{update z k}.
\end{itemize}
\paragraph{C. The parametric solver}
Fix x, k and R > 0 with the closed sup-ball (a polydisc) of radius 2R around x inside U. Choose a cutoff χ : ℂ → ℝ that equals 1 on B(x\textsubscript{k}, R) and has support in the closed disc B̄(x\textsubscript{k}, 3R/2). Write z[k ↦ ζ] for z with its k-th coordinate replaced by ζ, and set

\$\$ G(z)=\textbackslash{}frac\{1\}\{\textbackslash{}pi\}\textbackslash{}int\textsubscript{\textbackslash{}mathbb C}\textbackslash{}frac\{1\}\{t\}\textbackslash{},\textbackslash{}chi(z\textsubscript{k}-t)\textbackslash{},f\textbackslash{}bigl(z[k\textbackslash{}mapsto z\textsubscript{k}-t]\textbackslash{}bigr)\textbackslash{},dA(t) \$\$

This is the Mathlib convolution (K ⋆ g\textsubscript{z})(z\textsubscript{k}) with K(t) = t⁻¹ and g\textsubscript{z}(ζ) = χ(ζ)·f(z[k ↦ ζ]). For z in B(x,R) and ζ in the support of χ, the point z[k ↦ ζ] lies in B(x, 2R) ⊆ U. So g is jointly smooth on B(x,R) × ℂ and supported in a fixed compact set in ζ.

\begin{itemize}
\item \textbf{C1 (smoothness).} G is smooth on B(x,R), by \texttt{contDiffOn\_convolution\_right\_with\_param\_comp} with v(z) = z\textsubscript{k}.
\item \textbf{C2 (solves the k-th equation).} ∂̄\textsubscript{k} G = f on B(x,R). By B6, the slice through z is 2T h, where h(ζ) = χ(ζ)f(z[k ↦ ζ]) is smooth with compact support. TauCeti's Pompeiu converse gives ∂̄(2T h) = h, and h(z\textsubscript{k}) = f(z) because χ(z\textsubscript{k}) = 1.
\item \textbf{C3 (keeps holomorphy).} If j ≠ k and ∂̄\textsubscript{j} f = 0 on U, then ∂̄\textsubscript{j} G = 0 on B(x,R). Differentiate under the integral (\texttt{hasFDerivAt\_convolution\_right\_with\_param}); the integrand's ∂̄\textsubscript{j} derivative is χ·(∂̄\textsubscript{j} f)(z[k ↦ ζ]) = 0.
\item \textbf{C4 (naturality).} L(G\textsubscript{f}(z)) = G\textsubscript{L∘f}(z) for continuous ℂ-linear L, by \texttt{integral\_comp\_comm}. Hence if L ∘ f = 0 on U, then L ∘ G = 0 everywhere.
\end{itemize}
\paragraph{D. The induction}
\textbf{Claim P(k).} Let ω be a smooth (0,q+1)-form on an open U that is closed and involves only dz̄\textsubscript{1}, …, dz̄\textsubscript{k} on U. Then for each x ∈ U there are an open U′ with x ∈ U′ ⊆ U and a smooth (0,q)-form η on U′ with ∂̄η = ω on U′.

\textbf{Base, k = 0.} ω has degree q+1 ≥ 1 and all contractions vanish, so ω = 0 by A5. Take η = 0.

\textbf{Step, k ≥ 1.} Put α := ι\textsubscript{e\_k} ω.

\begin{enumerate}
\item α has type (0,q) (A4) and involves only dz̄\textsubscript{1}, …, dz̄\textsubscript{k−1}: ι\textsubscript{e\_m} ι\textsubscript{e\_k} ω = −ι\textsubscript{e\_k} ι\textsubscript{e\_m} ω is 0 for m > k, and is 0 for m = k by A2.
\item ∂̄\textsubscript{m} α = ι\textsubscript{e\_k} ∂̄\textsubscript{m} ω = 0 on U for m > k (B1, B5).
\item Let G be the solver of C applied to α. By C1–C4, G is smooth on B(x,R), has type (0,q), involves only dz̄\textsubscript{1}, …, dz̄\textsubscript{k−1} everywhere, satisfies ∂̄\textsubscript{k} G = α, and satisfies ∂̄\textsubscript{m} G = 0 for m > k.
\item Put ω′ := ω − ∂̄G on B(x,R). It is closed by B3 and has type (0,q+1) by A4 and B1.
\item ω′ involves only dz̄\textsubscript{1}, …, dz̄\textsubscript{k−1}. For m ≥ k, B4 gives ι\textsubscript{e\_m} ω′ = ι\textsubscript{e\_m} ω − ∂̄\textsubscript{m} G + ∂̄(ι\textsubscript{e\_m} G), and ι\textsubscript{e\_m} G ≡ 0. For m = k this is α − α = 0; for m > k it is 0 − 0.
\item P(k−1) applied to ω′ on B(x,R) gives η′ on U′ ∋ x. Then η := η′ + G satisfies ∂̄η = ω on U′.
\end{enumerate}
\textbf{Theorem.} Every (0,q+1)-form involves only dz̄\textsubscript{1}, …, dz̄\textsubscript{n}, so P(n) is the flat Dolbeault–Grothendieck lemma.

\begin{theorem}[Z4.a — Cauchy kernel]
\label{Z4.a}
\lean{Cohesive.locallyIntegrable_inv_complex}
\leanok
The Cauchy kernel t ↦ t\textsuperscript{−1} is locally integrable on ℂ with respect to Lebesgue measure.
\par\smallskip{\footnotesize Source: blueprint.md §11 v1.14.}
\end{theorem}
\begin{proof}
\leanok
See the Lean proof and the cited items.
\paragraph{Blueprint notes.}
(\texttt{Analysis/CauchyKernel.lean}): \texttt{LocallyIntegrable (fun t : ℂ ↦ t⁻¹)}, from Mathlib's \texttt{locallyIntegrable\_of\_norm\_le\_rpow} (α = 1 < 2 = dim\textsubscript{ℝ} ℂ).
\end{proof}

\begin{theorem}[Z4.b — One-variable ∂̄ (Cauchy–Pompeiu)]
\label{Z4.b}
\lean{Cohesive.dbar_cauchyTransform}
\leanok
Let F be a complex Banach space and h : ℂ → F real C¹ with compact support. Let u(w) = π\textsuperscript{−1} ∫\textsubscript{ℂ} t\textsuperscript{−1} h(w − t) dA(t). Then for every w ∈ ℂ, ½ (Du(w)(1) + i·Du(w)(i)) = h(w), i.e. ∂u/∂w̄ = h in the normalization of F3.
\par\smallskip{\footnotesize Source: blueprint.md §11 v1.14.}
\end{theorem}
\begin{proof}
\leanok
\uses{Z4.a}
See the Lean proof and the cited items.
\paragraph{Blueprint notes.}
(\texttt{Analysis/Pompeiu.lean}): for \texttt{h : ℂ → Λ} (Λ a complex Banach space), \texttt{ContDiff ℝ 1 h}, compact support, \texttt{u w = π⁻¹ ∫ t⁻¹ • h (w − t)} satisfies \texttt{antiPart (fderiv ℝ u w) (toConj 1) = h w}. Proof: \texttt{Du = π⁻¹(t⁻¹ ⋆ Dh)} (Mathlib convolution derivative); polar coordinates around w: the angular term integrates to 0 over each circle, the radial term to −h(w) along each ray, so ∫ ζ⁻¹ ∂̄h(w+ζ) dA = −π h(w). Proposed: port TauCeti \texttt{Pompeiu.lean} (Apache 2.0, with attribution in the module docstring) — needs the human's decision.
\end{proof}

\begin{theorem}[Z4.c — Parametric solver]
\label{Z4.c}
\lean{Cohesive.contDiff_cauchySolve, Cohesive.hasFDerivAt_cauchySolve, Cohesive.map_cauchySolve, Cohesive.dbar_cauchySolve}
\leanok
\uses{Z4.b}
Let E be a complex normed space and Λ a complex Banach space. For v ∈ E and F : E → Λ let u(z) = cauchySolve v F (z) = π\textsuperscript{−1} ∫\textsubscript{ℂ} t\textsuperscript{−1} F(z − t v) dA(t). Let v ≠ 0 and F have compact support.

\begin{itemize}
\item If F is real C\textsuperscript{∞}, then u is real C\textsuperscript{∞} on E.
\item If F is real C¹, then Du(z) = cauchySolve v (DF)(z) for every z.
\item If F is continuous and L : Λ → Λ′ is continuous ℂ-linear (Λ′ Banach), then L(u(z)) = cauchySolve v (L ∘ F)(z).
\item If F is real C¹, then ½ (Du(z)(v) + i·Du(z)(i v)) = F(z) for every z (∂̄ along v solves the equation).
\end{itemize}
\par\smallskip{\footnotesize Source: blueprint.md §11 v1.14.}
\end{theorem}
\begin{proof}
\leanok
See the Lean proof and the cited items.
\paragraph{Blueprint notes.}
(\texttt{Analysis/CauchySolve.lean}): \texttt{cauchySolve v F z := (π⁻¹ : ℂ) • ∫ t : ℂ, t⁻¹ • F (z − t • v)} for \texttt{F : E → Λ} C\textsuperscript{∞} with compact support: (i) C\textsuperscript{∞} on E with \texttt{fderiv (cauchySolve v F) = cauchySolve v (fderiv ℝ F)} (parametric convolution lemmas, g(z,ζ) = F(z + (ζ − ℓ z)•v), ℓ v = 1); (ii) \texttt{L (cauchySolve v F z) = cauchySolve v (L ∘ F) z} for ℂ-linear CLMs L; (iii) for v ≠ 0, \texttt{dbarPart (cauchySolve v F) z (toConj v) = F z} (slice through z, then Z4.b).
\end{proof}

\begin{theorem}[Z4.d — Contraction]
\label{Z4.d}
\lean{Cohesive.eq_zero_of_ctr_basis, Cohesive.dbarPart_comp_clm, Cohesive.ctr_dbar, Cohesive.ctr_dbar_zero, Cohesive.dbarPart_eq_zero_of_dbar_eq_zero}
\leanok
Let E, F be complex normed spaces. For v ∈ Ē, the contraction ι\textsubscript{v} : Λ\textsuperscript{0,q+1}(E; F) → Λ\textsuperscript{0,q}(E; F) is (ι\textsubscript{v} a)(w\textsubscript{1}, …, w\textsubscript{q}) = a(v, w\textsubscript{1}, …, w\textsubscript{q}).

\begin{itemize}
\item If ι\textsubscript{b\_i} a = 0 for every vector of a basis b of Ē, then a = 0.
\item D″(L ∘ ω)(z) = L ∘ D″ω(z) for continuous ℂ-linear maps L between fibres, ω real-differentiable at z.
\item For ω of degree q + 1, real-differentiable at z: ι\textsubscript{v} ∂̄ω(z) = D″ω(z)(v) − ∂̄(ι\textsubscript{v} ω)(z). For ω of degree 0: ι\textsubscript{v} ∂̄ω(z) = D″ω(z)(v).
\item Say a involves only dz̄\textsubscript{0}, …, dz̄\textsubscript{k−1} (for a basis b of Ē indexed by \{0, …, n−1\}) if ι\textsubscript{b\_j} a = 0 for all j ≥ k. If U ⊆ E is open, ω (degree q + 1) is real-differentiable on U, involves only dz̄\textsubscript{0}, …, dz̄\textsubscript{k−1} at every point of U and ∂̄ω = 0 on U, then D″ω(z)(b\textsubscript{j}) = 0 for all z ∈ U and j ≥ k.
\end{itemize}
\par\smallskip{\footnotesize Source: blueprint.md §11 v1.14.}
\end{theorem}
\begin{proof}
\leanok
\uses{F3}
See the Lean proof and the cited items.
\paragraph{Blueprint notes.}
(\texttt{Forms/Contraction.lean}): contraction = \texttt{curryLeft}; anticommutation, \texttt{curryLeft a (toConj v) = 0} on a basis ⇒ \texttt{a = 0}; \texttt{dbarPart} commutes with CLMs; B4 \texttt{curryLeft (dbar ω z) v = dbarPart ω z v − dbar (fun y ↦ curryLeft (ω y) v) z} (degree q + 1) and \texttt{curryLeft (dbar ω z) v = dbarPart ω z v} (degree 0); \texttt{InvolvesOnly b k a} (contractions with \texttt{b j}, j ≥ k, vanish); B5.
\end{proof}

\begin{theorem}[Z4.e — Flat lemma]
\label{Z4.e}
\lean{Cohesive.exists_dbar_eq, Cohesive.conjBasis, Cohesive.dbarPart_cauchySolve, Cohesive.dbarPart_cauchySolve_comm, Cohesive.involvesOnly_cauchySolve, Cohesive.exists_cutoff, Cohesive.contDiff_cutoff_smul}
\leanok
Let E be a finite-dimensional complex normed space with a fixed basis e\textsubscript{0}, …, e\textsubscript{n−1} (also a basis of Ē). Flat Dolbeault–Grothendieck lemma: let U ⊆ E be open, x ∈ U, and β : E → Λ\textsuperscript{0,q+1}(E; ℂ) real C\textsuperscript{∞} on U with ∂̄β = 0 on U. Then there are an open U′ with x ∈ U′ ⊆ U and η : E → Λ\textsuperscript{0,q}(E; ℂ), real C\textsuperscript{∞} on all of E, with ∂̄η = β on U′. Ingredients:

\begin{itemize}
\item for v ≠ 0 and F : E → Λ\textsuperscript{0,q}(E; ℂ) real C¹ with compact support, D″(cauchySolve v F)(z)(v) = F(z) and D″(cauchySolve v F)(z)(w) = cauchySolve v (D″F(·)(w))(z);
\item if every value of F involves only dz̄\textsubscript{0}, …, dz̄\textsubscript{k−1}, so does every value of cauchySolve v F;
\item cutoff: for x ∈ U and an index k there are a real C\textsuperscript{∞} ψ with compact support in U, ψ = 1 on an open W ∋ x, and an open slab S ⊇ W, stable under z ↦ z − t e\textsubscript{k} (t ∈ ℂ), on which the derivatives of ψ in the directions e\textsubscript{j}, i e\textsubscript{j} (j ≠ k) vanish; ψα is C\textsuperscript{∞} with compact support for every α that is C\textsuperscript{∞} on U.
\end{itemize}
\par\smallskip{\footnotesize Source: blueprint.md §11 v1.14.}
\end{theorem}
\begin{proof}
\leanok
\uses{F4, Z4.c, Z4.d}
See the Lean proof and the cited items.
\paragraph{Blueprint notes.}
(\texttt{Forms/FlatGrothendieck.lean}): cutoff ψ (S2), induction step (split on q = 0 / q + 1), \texttt{exists\_dbar\_eq} on open subsets of E with η globally C\textsuperscript{∞}.
\end{proof}

\begin{theorem}[Z4.f — Chart forms]
\label{Z4.f}
\lean{Cohesive.chartOpen, Cohesive.chartForm, Cohesive.chartExpr_chartForm, Cohesive.isDolbeault_of_chartExpr, Cohesive.isDolbeault_chartForm, Cohesive.isOfDegree_chartForm}
\leanok
Let E be a finite-dimensional complex normed space and X a complex manifold modelled on E. For x ∈ X, φ = chartAt(x), an open W ⊆ E and η : E → Λ\textsuperscript{0,q}(E; ℂ), the chart form of η is the form on U′ = φ.source ∩ φ\textsuperscript{−1}(W) whose value at y is η(φ(y)), placed in degree q and pulled back along the derivative of φ ∘ chartAt(y)\textsuperscript{−1} at chartAt(y)(y).

\begin{itemize}
\item If W ⊆ φ.target, the expression of the chart form in φ is η on W.
\item (Converse of F9.a.) If U ⊆ e.source for an atlas chart e and the expression of a form α on U in e is real C\textsuperscript{∞} at e(y) for all y ∈ U, then α is a Dolbeault form.
\item The chart form has degree q; if W ⊆ φ.target and η is real C\textsuperscript{∞} on W, it is a Dolbeault form.
\end{itemize}
\par\smallskip{\footnotesize Source: blueprint.md §11 v1.14.}
\end{theorem}
\begin{proof}
\leanok
\uses{F9}
See the Lean proof and the cited items.
\paragraph{Blueprint notes.}
(\texttt{Forms/ChartForm.lean}): from a flat form η on W ⊆ (chartAt x).target, a Dolbeault form on \texttt{(chartAt x).source ∩ (chartAt x)⁻¹' W} whose expression in \texttt{chartAt x} is η; injectivity of \texttt{chartExpr} at points of the chart. Needed private lemmas of \texttt{Forms/SheafBasic.lean}, \texttt{Forms/DolbeaultFree.lean} become public (visibility only).
\end{proof}

\begin{theorem}[Z4.g — H1: Dolbeault–Grothendieck lemma]
\label{Z4.g}
\lean{Cohesive.dolbeaultGrothendieck, Cohesive.flatExpr}
\leanok
Let E be a finite-dimensional complex normed space and X a complex manifold modelled on E (no compactness or Hausdorff assumption). Then X satisfies the Dolbeault–Grothendieck lemma H1. Proof outline: the flat expression of degree r of a form α in chartAt(x) (the degree-r component of its chart expression) is C\textsuperscript{∞} and ∂̄-closed near chartAt(x)(x); the flat lemma (Z4.e) solves ∂̄η = it, and the chart form of η (Z4.f) is the required local primitive.
\par\smallskip{\footnotesize Source: blueprint.md §11 v1.14.}
\end{theorem}
\begin{proof}
\leanok
\uses{H1, Z4.e, Z4.f}
See the Lean proof and the cited items.
\paragraph{Blueprint notes.}
(\texttt{Forms/DolbeaultGrothendieck.lean}): \texttt{theorem dolbeaultGrothendieck [IsManifold 𝓘(ℂ, E) ∞ X] : DolbeaultGrothendieck E X}.
\end{proof}

\begin{theorem}[Z4.h — Corollary]
\label{Z4.h}
\lean{Cohesive.blockPerf_of_localGaugeLemma, main_theorem_of_localGaugeLemma}
\leanok
\uses{H1}
Let E be a finite-dimensional complex normed space and X a compact Hausdorff complex manifold modelled on E, satisfying the local gauge lemma H2. Then J : Ho(P\textsubscript{A}) ⥤ D(O\textsubscript{X}), A = 𝒜\textsubscript{X}(X), is full and faithful with essential image D\textsubscript{perf}(X, O\textsubscript{X}), for the standard derived category and for any other \texttt{HasDerivedCategory} instance. (H1 holds by Z4.g.)
\par\smallskip{\footnotesize Source: blueprint.md §11 v1.14.}
\end{theorem}
\begin{proof}
\leanok
\uses{M, Z4.g}
See the Lean proof and the cited items.
\paragraph{Blueprint notes.}
(\texttt{MainH1.lean}): \texttt{MainStatement} without H1, as a theorem only (no new statement to freeze; a new main statement follows after Z5). \texttt{\#print axioms} checks.
\end{proof}

% Generated by scripts/blueprint_gen.py; do not edit.
\chapter{Proof of H2: the local gauge lemma}\label{chap:Z5}
\textbf{Target.} \texttt{theorem localGaugeLemma : LocalGaugeLemma E X} (\texttt{Cohesivemodules/Hypotheses.lean}, FROZEN), for every complex manifold \texttt{X} modelled on a finite-dimensional \texttt{E}. The statement: for every cohesive module \texttt{M} over \texttt{A = 𝒜\_X(X)} and \texttt{x ∈ X}, there are an open \texttt{U′ ∋ x}, a bounded complex \texttt{W} of finitely generated projective \texttt{O\_X(U′)}-modules, and an isomorphism in \texttt{Ho(P\_\{𝒜\_X(U′)\})} between \texttt{M|U′ = M.baseChange (dolbeaultRes E le\_top)} and \texttt{𝒜\_X(U′) ⊗\_\{O(U′)\} W = CohesiveModule.fromComplex (ρ := dolbeaultRhoU E U′) W …}.

\textbf{What we prove.} We prove the strong form, an actual isomorphism of cohesive modules with \texttt{W} free, as in Block's proof and in Bismut–Shen–Wei, Thm 5.3. Near \texttt{x}, in a chart and a frame, the cohesive operator is \texttt{𝕏 = ∂̄ + 𝒜}. Here \texttt{𝒜} is a smooth function with values in fibre operators, and flatness is the Maurer–Cartan equation \texttt{Σⱼ εⱼ ∂ⱼ𝒜 + 𝒜² = 0}. A gauge transformation, built one variable at a time following Donaldson–Kronheimer §2.2.2, makes \texttt{𝒜} free of forms and holomorphic. What remains is a complex of trivial holomorphic bundles, \texttt{W = (O\textasciicircum{}r, A₀)}.

\textbf{Five layers.}

\begin{enumerate}
\item \emph{Analytic core (C).} The one-variable gauge equation \texttt{∂̄\_k g = g·a} with parameters, solved in a finite-dimensional Banach algebra by an elementary Neumann series. It reuses the Z4 Cauchy transform.
\item \emph{Operator model (B).} Flat ∂̄-superconnections on an open subset of \texttt{E}, the gauge action, and the algebraic identities. All of this is pointwise algebra in \texttt{FibOp E r} together with derivative rules.
\item \emph{Induction (D).} Removing \texttt{dz̄₁, …, dz̄ₙ} one at a time.
\item \emph{Local freeness and frames (A, E).} The generator modules of \texttt{M|U′} are free near \texttt{x}, and the frame model connects \texttt{M|U′} to layer B.
\item \emph{Assembly (F).} \texttt{W}, the isomorphism in \texttt{HoP}, and \texttt{localGaugeLemma}.
\end{enumerate}
\textbf{Simplifications that drive the design.}

\begin{itemize}
\item The gauge equation is \emph{linear} (\texttt{∂̄\_k g = g a}), so no nonlinear fixed-point theory is needed.
\item With the operator formulation, holomorphy after each step follows by contracting the Maurer–Cartan equation with \texttt{ι\_i}. Block's p-degree separation argument is not needed (§1, B.4).
\item Smoothness comes from summable bounds on all derivatives of the Neumann series terms (\texttt{contDiff\_tsum}, in Mathlib). There is no Hölder or Sobolev space and no elliptic regularity; this closes the smoothness gap that all sources leave (§1, a5).
\item The global cutoff of Z4.e keeps every function globally defined. The slab property preserves holomorphy in the earlier variables.
\end{itemize}

\begin{theorem}[Z5.a — Kernel estimate]
\label{Z5.a}
\lean{Cohesive.integral_inv_norm_ball_zero, Cohesive.integrableOn_inv_norm_ball, Cohesive.integral_inv_norm_ball_le, Cohesive.norm_cauchySolve_le}
\leanok
\begin{itemize}
\item For ρ ≥ 0: ∫\textsubscript{B(0,ρ)} |t|\textsuperscript{−1} dA(t) = 2πρ.
\item For every c ∈ ℂ, t ↦ |t|\textsuperscript{−1} is integrable on B(c, ρ), and ∫\textsubscript{B(c,ρ)} |t|\textsuperscript{−1} dA(t) ≤ 3πρ for ρ ≥ 0.
\item Let E, Λ be complex normed spaces, v ∈ E, ℓ : E → ℂ continuous ℂ-linear with ℓ(v) = 1, and G : E → Λ with ‖G‖ ≤ C and G(w) ≠ 0 only where |ℓ(w) − c| ≤ ρ (ρ ≥ 0). Then ‖cauchySolve v G (z)‖ ≤ 3ρC for every z ∈ E.
\end{itemize}
\par\smallskip{\footnotesize Source: blueprint-ll.md §5.}
\end{theorem}
\begin{proof}
\leanok
\uses{Z4.c}
See the Lean proof and the cited items.
\paragraph{Blueprint notes.}
\texttt{∫\_\{B(c,ρ)\} \textbackslash{}|t\textbackslash{}|⁻¹ ≤ 3πρ} (polar coordinates), \texttt{norm\_cauchySolve\_le}

Files: \texttt{Analysis/KernelEstimate.lean}.
\end{proof}

\begin{theorem}[Z5.b — Derivatives of the solver]
\label{Z5.b}
\lean{Cohesive.pdbar, Cohesive.iteratedFDeriv_cauchySolve, Cohesive.pdbar_cauchySolve, Cohesive.pdbar_cauchySolve_comm}
\leanok
For complex normed spaces E, Λ, v ∈ E and f : E → Λ let ∂̄\textsubscript{v} f(z) = ½ (Df(z)(v) + i·Df(z)(i v)) (junk where f is not differentiable). Let E be finite-dimensional, Λ a Banach space, v ≠ 0 and F : E → Λ with compact support.

\begin{itemize}
\item If F is real C\textsuperscript{∞}: D\textsuperscript{m}(cauchySolve v F)(z) = cauchySolve v (D\textsuperscript{m} F)(z) for all m ∈ ℕ, z ∈ E.
\item If F is real C¹: ∂̄\textsubscript{v}(cauchySolve v F) = F, and ∂̄\textsubscript{w}(cauchySolve v F) = cauchySolve v (∂̄\textsubscript{w} F) for every w ∈ E.
\end{itemize}
\par\smallskip{\footnotesize Source: blueprint-ll.md §5.}
\end{theorem}
\begin{proof}
\leanok
\uses{Z4.c}
See the Lean proof and the cited items.
\paragraph{Blueprint notes.}
\texttt{pdbar} (+ link to \texttt{dbarPart}), \texttt{iteratedFDeriv\_cauchySolve}, \texttt{pdbar\_cauchySolve}, \texttt{pdbar\_cauchySolve\_comm} for any normed codomain; make Z4.e's \texttt{cauchySolve\_eq\_zero}, \texttt{fderiv\_cutoff\_smul*} public

Files: \texttt{Analysis/CauchySolveDeriv.lean}.
\end{proof}

\begin{theorem}[Z5.c — One-variable gauge solver]
\label{Z5.c}
\lean{Cohesive.exists_gauge_solution, Cohesive.coordCLM, Cohesive.exists_cutoff_radius, Cohesive.neumannTerm, Cohesive.gaugeSeries, Cohesive.norm_iteratedFDeriv_neumannTerm_le, Cohesive.contDiff_gaugeSeries}
\leanok
Let E be a finite-dimensional complex normed space with a fixed basis e\textsubscript{0}, …, e\textsubscript{n−1} and coordinate functionals c\textsubscript{k}, and 𝔄 a complex Banach algebra with ‖1‖ = 1.

\begin{itemize}
\item Cutoff: for a neighbourhood K of x and an index k, for all small r > 0 there is ψ as in Z4.e (with compact support in K), with 0 ≤ ψ ≤ 1 and |c\textsubscript{k}(w) − c\textsubscript{k}(x)| ≤ 2r wherever ψ(w) ≠ 0.
\item Neumann series: u\textsubscript{0} = 1, u\textsubscript{i+1} = cauchySolve v (u\textsubscript{i} · b), g = Σ\textsubscript{i} u\textsubscript{i}. If v ≠ 0, b is C\textsuperscript{∞} with compact support, ℓ(v) = 1, b vanishes outside \{|ℓ(w) − c| ≤ ρ\}, ‖b‖ ≤ B and θ = 3ρB < 1, then ‖D\textsuperscript{m} u\textsubscript{i}(z)‖ ≤ C\textsubscript{m} (i+1)\textsuperscript{m} θ\textsuperscript{i} for all m, i, z, and g is C\textsuperscript{∞} on E.
\item Main result (𝔄 finite-dimensional): let U ⊆ E be open, x ∈ U, a : E → 𝔄 real C\textsuperscript{∞} on U with values in a subalgebra S on U, and J a set of indices with k ∉ J and ∂̄\textsubscript{e\_j} a = 0 on U for j ∈ J. Then there are an open W with x ∈ W ⊆ U and g : E → 𝔄, real C\textsuperscript{∞}, with values in S and pointwise invertible, such that ∂̄\textsubscript{e\_k} g = g·a and ∂̄\textsubscript{e\_j} g = 0 (j ∈ J) on W.
\end{itemize}
\par\smallskip{\footnotesize Source: blueprint-ll.md §5.}
\end{theorem}
\begin{proof}
\leanok
\uses{Z4.e, Z5.a, Z5.b}
\textbf{C.0 Cutoff with radius control.} This is a variant of Z4.e \texttt{exists\_cutoff} exposing its construction \texttt{ψ(z) = χ₁(c\_k(z−x))·χ₂(P(z−x))}. Fix a compact neighbourhood \texttt{K ⊆ U} of \texttt{x} and \texttt{M\_K = sup\_K ‖a‖}. For every small \texttt{r > 0} there is \texttt{ψ} with:

\begin{itemize}
\item \texttt{0 ≤ ψ ≤ 1} and \texttt{tsupport ψ ⊆ K ∩ \{|c\_k(w−x)| ≤ 2r\}};
\item \texttt{ψ = 1} on an open \texttt{W ∋ x};
\item the slab \texttt{S\_r ⊇ W} is stable under \texttt{z ↦ z − t e\_k}, and \texttt{∂ⱼψ = 0} on it for \texttt{j ≠ k}.
Put \texttt{b := ψ·a} (extended by 0): \texttt{b ∈ C\_c\textasciicircum{}∞}, \texttt{‖b‖\_∞ ≤ M\_K}.

\end{itemize}
\textbf{C.3 Bounds.} Let \texttt{u₀ = 1} and \texttt{u\_\{i+1\} = T(u\_i b)}; each \texttt{u\_i} is \texttt{C\textasciicircum{}∞} with bounded derivatives (C.2, C.1).

\begin{itemize}
\item Notation: \texttt{θ = 6r M\_K}, \texttt{β\_l = sup ‖D\textasciicircum{}l b‖}, \texttt{M\_m(i) = sup ‖D\textasciicircum{}m u\_i‖}.
\item Recursion (Leibniz, \texttt{norm\_iteratedFDeriv\_mul\_le}): \texttt{M\_m(i+1) ≤ 6r Σ\_\{l≤m\} C(m,l) β\_\{m−l\} M\_l(i) ≤ θ M\_m(i) + B\_m (i+1)\textasciicircum{}\{m−1\} θ\textasciicircum{}i}.
\item Hence \texttt{M\_m(i) ≤ C\_m (i+1)\textasciicircum{}m θ\textasciicircum{}i}. Choose \texttt{r} with \texttt{θ < 1/2} first; the \texttt{β\_l} for \texttt{l ≥ 1} may then be large but are fixed. Only \texttt{β₀ ≤ M\_K} enters \texttt{θ}.
\end{itemize}
\textbf{C.4 Smoothness.} \texttt{g := Σ\_i u\_i} is \texttt{C\textasciicircum{}∞} by \texttt{contDiff\_tsum} with \texttt{v m i = C\_m (i+1)\textasciicircum{}m θ\textasciicircum{}i} (\texttt{summable\_pow\_mul\_geometric\_of\_norm\_lt\_one}).

\textbf{C.5 Equation.} \texttt{∂\_k g = Σ\_i ∂\_k u\_i = Σ\_i u\_i b = g b} (\texttt{fderiv\_tsum}, C.2, \texttt{tsum\_mul\_right}). On \texttt{W} (where \texttt{ψ = 1}), \texttt{∂\_k g = g a}.

\textbf{C.6 Subalgebra.} Let \texttt{Q : 𝔄 → 𝔄/S} be the quotient map (a CLM, finite dimension). \texttt{Q(T F) = T(Q∘F)} (\texttt{map\_cauchySolve}), so \texttt{u\_i ∈ S} by induction, and \texttt{Q(g) = Σ Q(u\_i) = 0} (\texttt{ContinuousLinearMap.map\_tsum}).

\textbf{C.7 Invertibility.} \texttt{‖g(z) − 1‖ ≤ Σ\_\{i≥1\} θ\textasciicircum{}i < 1}, so \texttt{g(z)} is a unit (\texttt{Units.oneSub}).

\textbf{C.8 Holomorphy.} For \texttt{j ∈ J}: on \texttt{S\_r}, \texttt{∂ⱼ b = (∂ⱼψ)a + ψ∂ⱼa = 0} (off \texttt{U}, \texttt{b ≡ 0} near the point). By induction \texttt{∂ⱼ u\_i = 0} on \texttt{S\_r}: \texttt{∂ⱼ u\_\{i+1\} = T(∂ⱼu\_i·b + u\_i·∂ⱼb)}, the argument vanishes on \texttt{S\_r}, and the lines \texttt{z − t e\_k} stay in \texttt{S\_r} (\texttt{cauchySolve\_eq\_zero}). Then \texttt{∂ⱼg = Σ ∂ⱼu\_i = 0} on \texttt{S\_r ⊇ W}.
\paragraph{Blueprint notes.}
cutoff with radius control (C.0), Neumann bounds (C.3), \texttt{exists\_gauge\_solution} (C.4–C.8)

Files: \texttt{Analysis/GaugeSolve.lean}.
\end{proof}

\begin{theorem}[Z5.d — Fibre operators]
\label{Z5.d}
\lean{Cohesive.dzbarTot, Cohesive.ctrTot, Cohesive.sgnTot, Cohesive.Fib, Cohesive.FibOp, Cohesive.compOp}
\leanok
Let E be a finite-dimensional complex normed space with a fixed basis e\textsubscript{0}, …, e\textsubscript{n−1}, and r ∈ ℕ. The fibre of the local model is Fib = Λ\textsuperscript{0,•}(E)\textsuperscript{r}; fibre operators are the continuous ℂ-linear endomorphisms of Fib (a finite-dimensional Banach algebra). Define: dz̄\textsubscript{j} ∈ Λ\textsuperscript{0,•}(E), the 1-forms dual to the basis; the contraction ι\textsubscript{j} of Λ\textsuperscript{0,•}(E) with e\textsubscript{j}; the grade involution σ = (−1)\textsuperscript{q} on Λ\textsuperscript{0,q}; and on Fib, componentwise, ε\textsubscript{j} = left multiplication by dz̄\textsubscript{j}, ι\textsubscript{j}, σ. Results:

\begin{itemize}
\item ε\textsubscript{iε\_j} = −ε\textsubscript{jε\_i}, ι\textsubscript{iι\_j} = −ι\textsubscript{jι\_i}, ι\textsubscript{jι\_j} = 0, ι\textsubscript{iε\_j} + ε\textsubscript{jι\_i} = δ\textsubscript{ij};
\item an operator commuting with every ε\textsubscript{j} commutes with multiplication by every t ∈ Λ\textsuperscript{0,•}(E); one anticommuting with every ε\textsubscript{j} satisfies T(t·s) = σ(t)·T(s);
\item an element of Λ\textsuperscript{0,•}(E) killed by every ι\textsubscript{j} has form degree 0.
\end{itemize}
For deg : \{0, …, r−1\} → ℤ, an element of Fib has total degree m if its i-th component has form degree m − deg(i); an operator has total degree d if it raises total degree by d.
\par\smallskip{\footnotesize Source: blueprint-ll.md §5.}
\end{theorem}
\begin{proof}
\leanok
\uses{F5}
\textbf{B.1 Coordinate formula.} For a section \texttt{s : V → Fib}, \texttt{∂̄s = Σⱼ εⱼ ∂ⱼ s}, componentwise \texttt{dbarTot}. Proof: on \texttt{Ē}, \texttt{antiPart L v = Σⱼ b\_j\textasciicircum{}*(v) antiPart L (b\_j)}, and \texttt{alternatizeUncurryFin (Σⱼ b\_j\textasciicircum{}* ⊗ cⱼ) = Σⱼ dz̄ⱼ ∧ cⱼ} (F6 wedge).
\paragraph{Blueprint notes.}
\texttt{Fib}, \texttt{FibOp}, \texttt{epsOp}, \texttt{iotaOp}, CAR relations, \texttt{IsOfTotDeg}, Λ-linearity, kernel lemma, B.1 coordinate formula

Files: \texttt{Forms/FibreOps.lean}.
\end{proof}

\begin{theorem}[Z5.e — Flat superconnections and gauges]
\label{Z5.e}
\lean{Cohesive.dbarFib, Cohesive.dbarFib_eq_sum, Cohesive.dbarFib_add_sq, Cohesive.gauge_conj, Cohesive.gaugeOp_odd, Cohesive.mcOp_gaugeOp}
\leanok
Let E be a finite-dimensional complex normed space, ∂̄\textsubscript{j} = ∂̄\textsubscript{e\_j} (Z5.b), and 𝒜, g : E → FibOp (Z5.d). For s : E → Fib, ∂̄s is taken componentwise (F9). Let mc(𝒜) = Σ\textsubscript{j} ε\textsubscript{j} ∂̄\textsubscript{j𝒜} + 𝒜² and g·𝒜 = g𝒜g\textsuperscript{−1} − (Σ\textsubscript{j} ε\textsubscript{j} ∂̄\textsubscript{j} g) g\textsuperscript{−1}. Call 𝒜 odd at z if ε\textsubscript{j} 𝒜(z) = −𝒜(z) ε\textsubscript{j} for all j.

\begin{itemize}
\item ∂̄s(z) = Σ\textsubscript{j} ε\textsubscript{j} ∂̄\textsubscript{j} s(z) if s is real-differentiable at z.
\item If 𝒜 is C¹ at z and odd at z, and s is C² at z: (∂̄ + 𝒜)((∂̄ + 𝒜)s)(z) = mc(𝒜)(z) s(z).
\item Let g be pointwise invertible and commute with every ε\textsubscript{j}. If g is C¹ at z and s differentiable at z: g(∂̄ + 𝒜)(g\textsuperscript{−1}s) = (∂̄ + g·𝒜)s at z. If 𝒜 is odd at z, so is g·𝒜. On an open U where g and 𝒜 are C\textsuperscript{∞} and 𝒜 is odd: mc(g·𝒜) = g mc(𝒜) g\textsuperscript{−1} on U.
\end{itemize}
\par\smallskip{\footnotesize Source: blueprint-ll.md §5.}
\end{theorem}
\begin{proof}
\leanok
\uses{F4, Z5.b, Z5.d}
\textbf{B.2 Curvature.} If \texttt{𝒜} is odd and Λ-linear, then \texttt{(∂̄ + 𝒜)² s = mcOp 𝒜 · s} with \texttt{mcOp 𝒜 = Σⱼ εⱼ ∂ⱼ𝒜 + 𝒜²}. This uses \texttt{∂̄² = 0} (F4), the product rule \texttt{∂ⱼ(𝒜 s) = (∂ⱼ𝒜)s + 𝒜 ∂ⱼ s}, and \texttt{𝒜εⱼ = −εⱼ𝒜}. Applied to constant sections, \texttt{𝕏² = 0 ⇔ mcOp 𝒜 = 0}.

\textbf{B.3 Gauge action.} For \texttt{g} smooth, pointwise invertible, even, Λ-linear and of degree 0, \texttt{g ∘ (∂̄ + 𝒜) ∘ g⁻¹ = ∂̄ + g·𝒜}, where \texttt{g·𝒜 := g𝒜g⁻¹ − (Σⱼ εⱼ ∂ⱼg) g⁻¹} (use \texttt{g ∂ⱼ(g⁻¹) = −(∂ⱼg)g⁻¹}).

\begin{itemize}
\item \texttt{g·𝒜} is again odd, Λ-linear and of degree 1.
\item \texttt{mcOp (g·𝒜) = g (mcOp 𝒜) g⁻¹}.
\item \texttt{(hg)·𝒜 = h·(g·𝒜)}.
\end{itemize}
\textbf{B.4 Contracted MC.} With \texttt{aᵢ := ιᵢ𝒜 + 𝒜ιᵢ} (even, degree 0): \texttt{ιᵢ mcOp 𝒜 − mcOp 𝒜 ιᵢ = ∂ᵢ𝒜 − Σⱼ εⱼ ∂ⱼaᵢ + aᵢ𝒜 − 𝒜aᵢ}.

\begin{itemize}
\item Check: \texttt{[ιᵢ, εⱼ∂ⱼ𝒜] = δᵢⱼ∂ⱼ𝒜 − εⱼ∂ⱼaᵢ} and \texttt{[ιᵢ, 𝒜²] = aᵢ𝒜 − 𝒜aᵢ}.
\item Corollary: if \texttt{mcOp 𝒜 = 0} and \texttt{aᵢ ≡ 0} on an open set, then \texttt{∂ᵢ𝒜 = 0} there.
\end{itemize}
\textbf{B.5 The step algebra.} Suppose \texttt{g} (as in B.3) commutes with \texttt{ιᵢ} for \texttt{i ∈ I ∪ \{k\}}, satisfies \texttt{∂ᵢg = 0} for \texttt{i ∈ I}, and \texttt{∂\_k g = g a\_k}. Then the coefficients of \texttt{𝒜′ = g·𝒜} are \texttt{a′ᵢ = g aᵢ g⁻¹ − (∂ᵢg)g⁻¹}; so \texttt{a′ᵢ = g aᵢ g⁻¹} for \texttt{i ∈ I}, and \texttt{a′\_k = 0}.

**B.6 Properties of \texttt{a\_k}.** If \texttt{aᵢ ≡ 0} on \texttt{U} for \texttt{i < k}, then on \texttt{U}:

\begin{itemize}
\item \texttt{a\_k} is even, Λ-linear (\texttt{[εⱼ, a\_k] = 0}) and of degree 0;
\item \texttt{[ιᵢ, a\_k] = 0} for \texttt{i ≤ k} (super Jacobi; \texttt{ι\_k² = 0});
\item \texttt{∂ᵢ a\_k = [ι\_k, ∂ᵢ𝒜]₊ = 0} for \texttt{i < k} (B.4).
\end{itemize}
\textbf{B.7 End of the induction.} If \texttt{aᵢ ≡ 0} for all \texttt{i}, then \texttt{𝒜(1 ⊗ f) ∈ Λ⁰ ⊗ F} (kernel lemma) and \texttt{𝒜(ω·s) = (−1)\textasciicircum{}\{|ω|\} ω·𝒜(s)}. So \texttt{𝒜 = σ ⊗ A₀} with \texttt{A₀ : F → F} of degree +1. By B.4, \texttt{∂ᵢA₀ = 0} for all \texttt{i}, so \texttt{A₀} is holomorphic; MC then gives \texttt{A₀² = 0}.
\paragraph{Blueprint notes.}
\texttt{mcOp}, \texttt{gaugeOp}, B.2–B.7

Files: \texttt{Forms/Superconnection.lean}.
\end{proof}

\begin{theorem}[Z5.f — Gauge induction]
\label{Z5.f}
\lean{Cohesive.gaugeSubalg, Cohesive.IsFlatSuperconn, Cohesive.IsGauge, Cohesive.exists_gauge_step, Cohesive.exists_gauge_of_aCoef, Cohesive.exists_gauge_holomorphic}
\leanok
Notation of Z5.d–Z5.e, with degrees deg : \{0, …, r−1\} → ℤ. A \emph{flat superconnection} on an open U ⊆ E is 𝒜 : E → FibOp, C\textsuperscript{∞} on U, of total degree 1, odd and with mc(𝒜) = 0 on U. A \emph{gauge} is g : E → FibOp, C\textsuperscript{∞} on E, pointwise invertible, of total degree 0, commuting with every ε\textsubscript{j}. Let a\textsubscript{i}(𝒜) = ι\textsubscript{i𝒜} + 𝒜ι\textsubscript{i} (the dz̄\textsubscript{i}-coefficient). Main result: if 𝒜 is a flat superconnection on U and x ∈ U, there are an open W with x ∈ W ⊆ U and a gauge g such that g·𝒜 is a flat superconnection on W with a\textsubscript{i}(g·𝒜) = 0 and ∂̄\textsubscript{e\_i}(g·𝒜) = 0 on W for every i. It follows by induction on k from:

\begin{itemize}
\item step: if a\textsubscript{i}(𝒜) = 0 on U for i < k, then near x there is a gauge g with a\textsubscript{i}(g·𝒜) = 0 for i ≤ k (g solves ∂̄\textsubscript{k} g = g·a\textsubscript{k} with values in the subalgebra of operators of total degree 0 commuting with every ε\textsubscript{j} and with ι\textsubscript{i} for i ≤ k; Z5.c).
\end{itemize}
\par\smallskip{\footnotesize Source: blueprint-ll.md §5.}
\end{theorem}
\begin{proof}
\leanok
\uses{Z5.c, Z5.e}
\paragraph{D. Induction (\texttt{exists\_gauge\_holomorphic})}
\begin{itemize}
\item **Claim P(k), \texttt{k = n, n−1, …, 0}.** For every open \texttt{U ∋ x} and every \texttt{𝒜} on \texttt{U} that is smooth, odd, Λ-linear, of degree 1 and MC, with \texttt{aᵢ ≡ 0} on \texttt{U} for \texttt{i < k}, there are \texttt{W ∋ x} and a gauge \texttt{g} (smooth, invertible, even, Λ-linear, degree 0) such that, on \texttt{W}, \texttt{g·𝒜 = σ ⊗ A₀} with \texttt{A₀} holomorphic and \texttt{A₀² = 0}.
\item \textbf{P(n):} take \texttt{g = 1} and B.7.
\item \textbf{P(k) from P(k+1):} apply C with:
\begin{itemize}
\item \texttt{a := a\_k};
\item \texttt{S := \{T | degree 0, [εⱼ,T] = 0 ∀j, [ιᵢ,T] = 0 ∀ i ≤ k\}};
\item \texttt{J := \{i < k\}}.
The hypotheses hold by B.6. C gives \texttt{W₁, g₁}. Then \texttt{𝒜₁ := g₁·𝒜} satisfies the hypotheses of P(k+1) on \texttt{W₁} (B.3, B.5). P(k+1) gives \texttt{W₂, g₂}; take \texttt{g := g₂g₁} (B.3 composition).

\end{itemize}
\item \texttt{exists\_gauge\_holomorphic} is P(0).
\end{itemize}
\paragraph{Blueprint notes.}
\texttt{exists\_gauge\_holomorphic} (D)

Files: \texttt{Forms/GaugeInduction.lean}.
\end{proof}

\begin{theorem}[Z5.g — Local freeness]
\label{Z5.g}
\lean{Cohesive.exists_free_baseChange, Cohesive.coef0, Cohesive.evalForm0, Cohesive.continuous_evalForm0, Cohesive.isUnit_of_evalForm0_ne_zero, Cohesive.exists_free_away, Cohesive.free_baseChange_of_isUnit}
\leanok
\begin{itemize}
\item Ring theory: a finitely generated projective module Q over a commutative ring R is free over a basic open containing a given prime 𝔪: Q[r\textsuperscript{−1}] is free over R[r\textsuperscript{−1}] for some r ∉ 𝔪. If Q[r\textsuperscript{−1}] is free and r becomes a unit in the R-algebra S, then S ⊗\textsubscript{R} Q is free over S.
\item Let E be a finite-dimensional complex normed space and X a complex manifold modelled on E. Evaluation at y ∈ U is a ring homomorphism 𝒜\textsuperscript{0,0}\textsubscript{X}(U) → ℂ; a degree-0 form is a continuous function on U, and a nowhere-vanishing one is a unit of 𝒜\textsuperscript{0,0}\textsubscript{X}(U).
\item Main result: for a cohesive module M over 𝒜\textsubscript{X}(X) with generators Q\textsuperscript{p} and x ∈ X, there is an open U′ ∋ x inside the source of chartAt(x) such that 𝒜\textsuperscript{0,0}\textsubscript{X}(U′) ⊗\_\{𝒜\textsuperscript{0,0}\textsubscript{X}(X)\} Q\textsuperscript{p} is a free 𝒜\textsuperscript{0,0}\textsubscript{X}(U′)-module for every p.
\end{itemize}
\par\smallskip{\footnotesize Source: blueprint-ll.md §5.}
\end{theorem}
\begin{proof}
\leanok
\uses{F12}
\textbf{A.1} A f.g. projective module \texttt{P} over \texttt{R = A⁰(X)} (degree-0 Dolbeault forms, i.e. smooth ℂ-valued functions) is the image of an idempotent \texttt{ε ∈ M\_N(R)}.

\textbf{A.2} \texttt{y ↦ ε(y) ∈ M\_N(ℂ)} is continuous and idempotent, so \texttt{rank ε(y) = tr ε(y)} is locally constant, equal to \texttt{r} near \texttt{x}. Choose \texttt{c₁…c\_r ∈ ℂ\textasciicircum{}N} with \texttt{ε(x)cᵢ} a basis of \texttt{im ε(x)}, and an \texttt{r × r} minor \texttt{μ(y)} of \texttt{[ε(y)c₁ … ε(y)c\_r]} with \texttt{det μ(x) ≠ 0}. On an open \texttt{U′ ∋ x}, \texttt{det μ ≠ 0} and \texttt{rank ε = r}.

\textbf{A.3} On \texttt{U′} the sections \texttt{σᵢ = ε cᵢ} form a basis of \texttt{A⁰(U′) ⊗\_R P}:

\begin{itemize}
\item they span, because \texttt{s = εs} gives \texttt{s(y) = Σ fᵢ(y) ε(y)cᵢ} with \texttt{f = μ⁻¹ ·} (the rows of \texttt{s}), which is smooth since \texttt{det μ ≠ 0} (Cramer; inverses of nonvanishing degree-0 forms are Dolbeault forms);
\item they are independent pointwise.
Shrink \texttt{U′} into \texttt{(chartAt x).source}. This gives \texttt{exists\_free\_baseChange}.

\end{itemize}
\begin{itemize}
\item \textbf{As built (blueprint §11 v1.17, recorded here by v1.18).} A.1–A.3 are replaced by localization (no idempotents, minors or Cramer): \texttt{m := ker (ev\_x : A⁰(X) → ℂ)} is prime; each \texttt{Q\textasciicircum{}p} has a free localization at some \texttt{r\_p ∉ m} (\texttt{Module.FinitePresentation.exists\_free\_localizedModule\_powers}); on \texttt{U′ := (chartAt x).source ∩ \{r\_p ≠ 0 for a ≤ p ≤ b\}} the \texttt{r\_p} are units, so \texttt{A⁰(U′) ⊗ Q\textasciicircum{}p} is free. Statement of \texttt{exists\_free\_baseChange} unchanged (\texttt{Forms/LocalFree.lean}).
\end{itemize}
\paragraph{Blueprint notes.}
A.1–A.3, \texttt{exists\_free\_baseChange}

Files: \texttt{Forms/LocalFree.lean}.
\end{proof}

\begin{theorem}[Z5.h — Frame model]
\label{Z5.h}
\lean{Cohesive.GradedFrame, Cohesive.GradedFrame.vec, Cohesive.GradedFrame.coefTot, Cohesive.GradedFrame.sum_coefTot, Cohesive.GradedFrame.coefTot_mem, Cohesive.GradedFrame.mat}
\leanok
Let A be a bounded graded-commutative dga (A1) over k and Q a graded A\textsuperscript{0}-module. A \emph{graded frame} of rank r is a degree function deg : \{0, …, r−1\} → ℤ and an A\textsuperscript{0}-basis (ẽ\textsubscript{i}) of ⨁\textsubscript{p} Q\textsuperscript{p} with ẽ\textsubscript{i} ∈ Q\textsuperscript{deg i} (coordinates of elements of Q\textsuperscript{p} vanish at frame vectors of other degrees).

\begin{itemize}
\item Every x ∈ A ⊗\textsubscript{A\textasciicircum{}0} ⨁\textsubscript{p} Q\textsuperscript{p} is uniquely Σ\textsubscript{i} x\textsubscript{i} (1 ⊗ ẽ\textsubscript{i}) with x\textsubscript{i} ∈ A; if x has total degree m, then x\textsubscript{i} ∈ A\textsuperscript{m − deg i}.
\item For a quasi-cohesive module N with a graded frame, the connection matrix ω\textsubscript{ji} ∈ A\textsuperscript{deg i + 1 − deg j} is the j-th coefficient of 𝕏(1 ⊗ ẽ\textsubscript{i}), and (𝕏x)\textsubscript{j} = d x\textsubscript{j} + Σ\textsubscript{i} σ(x\textsubscript{i}) ω\textsubscript{ji}.
\item Frame model: for N cohesive over 𝒜\textsubscript{X}(U) with a graded frame, x\textsubscript{0} ∈ U and φ = chartAt(x\textsubscript{0}), on V = φ(U ∩ φ.source) the chart expressions of the coefficients satisfy 𝕏 = ∂̄ + 𝒜 with 𝒜(z)(t)\textsubscript{j} = Σ\textsubscript{i} σ(t\textsubscript{i}) Ω\textsubscript{ji}(z) (Ω the chart expression of ω), and 𝒜 is a flat superconnection on V (Z5.f).
\end{itemize}
\par\smallskip{\footnotesize Source: blueprint-ll.md §5.}
\end{theorem}
\begin{proof}
\leanok
\uses{F10, F9, Z4.f, Z5.f}
\textbf{E.1 Frames.} On \texttt{U′ ⊆ chart source} with free generators (A), fix the frame \texttt{ẽᵢ} (degree \texttt{deg i}). Then \texttt{Tot(M|U′) = ⊕ᵢ DolbeaultAlg(U′) ẽᵢ} and \texttt{θ(ẽᵢ) = Σⱼ ωⱼᵢ ⊗ ẽⱼ} with \texttt{ωⱼᵢ} of degree \texttt{deg i + 1 − deg j}.

\textbf{E.2 Flat transport.}

\begin{itemize}
\item On \texttt{V = φ(U′)}, the chart expression of a section is \texttt{z ↦ (chartExpr sᵢ (z))ᵢ ∈ Fib}.
\item \texttt{𝒜(z)(t)\_j := Σᵢ σ(tᵢ)·chartExpr ωⱼᵢ(z)}, where \texttt{σ} is the degree sign on \texttt{TotFibre} and the product is in \texttt{TotFibre}.
\item Then \texttt{chartExpr(𝕏 s) = ∂̄(chartExpr s) + 𝒜·chartExpr s} (\texttt{chartExpr\_dbarSec\_eq}, \texttt{chartExpr\_mul}, B.1).
\item \texttt{𝒜} is smooth (F9.a), odd, Λ-linear and of degree 1; MC follows from \texttt{𝕏² = 0} via B.2, because every smooth \texttt{V → Fib} is the chart expression of a section (\texttt{chartForm}, Z4.f).
\item \emph{As built (v1.18 record):} MC is proved from \texttt{𝕏² = 0} on the frame vectors only, plus Λ-linearity of \texttt{mcOp}: \texttt{mcOp} is Λ-linear and even, so vanishing on the \texttt{r} constant generators of \texttt{Λ ⊗ ℂ\textasciicircum{}r} gives \texttt{mcOp = 0}. Same statement (\texttt{isFlatSuperconn\_frameOp}, \texttt{Forms/FrameModel.lean}); \texttt{chartForm} on all sections is not needed.
\end{itemize}
\paragraph{Blueprint notes.}
E.1–E.2 (frame, transported \texttt{𝒜}, MC from flatness)

Files: \texttt{Forms/FrameModel.lean}.
\end{proof}

\begin{theorem}[Z5.i — The local gauge lemma]
\label{Z5.i}
\lean{Cohesive.matComplex, Cohesive.exists_iso_fromComplex_of_frame, Cohesive.exists_iso_of_gauge, Cohesive.exists_local_iso}
\leanok
\begin{itemize}
\item For a commutative ring R, deg : \{0, …, r−1\} → ℤ and an r×r matrix f over R with f\textsubscript{ki} = 0 unless deg k = deg i + 1, and f² = 0: the complex W\textsuperscript{p} = R\textasciicircum{}\{\{i : deg i = p\}\}, d(w)\textsubscript{k} = Σ\textsubscript{i} f\textsubscript{ki} w\textsubscript{i}, a bounded complex of finitely generated projective R-modules.
\item Algebraic form: let A be a dga with curvature 0, ρ : R → A\textsuperscript{0} with d ∘ ρ = 0, N a cohesive module with graded frame of degrees deg and connection matrix ω, and G, H mutually inverse matrices over A with G\textsubscript{ij}, H\textsubscript{ij} ∈ A\textsuperscript{deg j − deg i} and dG\textsubscript{ij} + Σ\textsubscript{k} σ(G\textsubscript{kj}) ρ(f\textsubscript{ik}) = Σ\textsubscript{l} ω\textsubscript{lj} G\textsubscript{il}. Then N ≅ A ⊗\textsubscript{R} W in Ho(P\textsubscript{A}).
\item Geometric form: let X be a complex manifold modelled on a finite-dimensional E, x ∈ U ⊆ source of chartAt(x), N cohesive over 𝒜\textsubscript{X}(U) with a graded frame, g a gauge and B a flat superconnection on V = φ(U) with B = g·𝒜 (𝒜 the frame model), a\textsubscript{i}(B) = 0 and ∂̄\textsubscript{e\_i}B = 0 on V. Then N ≅ 𝒜\textsubscript{X}(U) ⊗\textsubscript{O\_X(U)} W in Ho(P\textsubscript{𝒜\_X(U)}), W a bounded complex of finitely generated projective O\textsubscript{X}(U)-modules.
\item Hence every cohesive module over 𝒜\textsubscript{X}(U′) with a graded frame is, near each x ∈ U′, isomorphic in Ho(P) to 𝒜\textsubscript{X}(U″) ⊗\textsubscript{O\_X(U″)} W after restriction to an open U″ ∋ x; with Z5.g this proves H2.
\end{itemize}
\par\smallskip{\footnotesize Source: blueprint-ll.md §5.}
\end{theorem}
\begin{proof}
\leanok
\uses{F12, H2, Z5.g, Z5.h}
**E.3 The complex \texttt{W}.**

\begin{itemize}
\item By D, on \texttt{W₀ ∋ φ x}: \texttt{g·𝒜 = σ ⊗ A₀}.
\item Let \texttt{U″ := chartOpen x W₀}. The entries of \texttt{A₀} are holomorphic on \texttt{W₀}; their \texttt{chartForm}s are ∂̄-closed degree-0 forms on \texttt{U″}, hence \texttt{ρ(fᵢⱼ)} for holomorphic \texttt{fᵢⱼ ∈ O(U″)} (F12).
\item Set \texttt{W\textasciicircum{}p := O(U″)\textasciicircum{}\{\#\{i : deg i = p\}\}} with differential \texttt{(fᵢⱼ)}. Then \texttt{d² = 0} (\texttt{ρ} injective, \texttt{A₀² = 0}), and \texttt{W} is bounded, f.g. and free (\texttt{IsBoundedFGProj}).
\item \emph{As built (v1.18 record):} this \texttt{ρ}-preimage construction is the one used (first option of the open question in §7; \texttt{exists\_hol\_of\_chartExpr}, \texttt{matComplex}, \texttt{Forms/LocalGauge.lean}).
\end{itemize}
\textbf{E.4 The isomorphism.} Let \texttt{G} be the \texttt{chartForm} transport of \texttt{g} (coefficients on \texttt{U″}).

\begin{itemize}
\item \texttt{Φ : fromComplex W → M|U″} via \texttt{fcExt}: \texttt{Φ(fcι wᵢ) = Σⱼ (G⁻¹)ⱼᵢ ⊗ ẽⱼ}. It is closed because \texttt{𝕏 G⁻¹ = G⁻¹(∂̄ + σ⊗A₀)} on generators, which is B.3 transported.
\item \texttt{Ψ : M|U″ → fromComplex W} on generators: \texttt{ẽⱼ ↦ Σᵢ Gᵢⱼ ⊗ fcι wᵢ} (\texttt{homEquivGen} / \texttt{famTot}-style extension).
\item \texttt{Ψ∘Φ = id} and \texttt{Φ∘Ψ = id} on generators, from \texttt{GG⁻¹ = 1} (\texttt{fc\_hom\_ext}, \texttt{hom\_ext} on generators).
\item The iso in \texttt{HoP} is built by the \texttt{DegreeZero} template. Base change from \texttt{X} to \texttt{U″} goes through \texttt{U′} (\texttt{bcCompIso}) or directly, freeness restricting along \texttt{U″ ≤ U′}.
\end{itemize}
\paragraph{F. Assembly}
\texttt{localGaugeLemma}: \texttt{intro M x}, then

\begin{enumerate}
\item A gives \texttt{U′}, frames;
\item E.1–E.2 give \texttt{𝒜} on \texttt{V};
\item D gives \texttt{W₀}, \texttt{g};
\item E.3–E.4 give \texttt{W} and the iso.
\end{enumerate}
The statement is exactly the frozen \texttt{LocalGaugeLemma E X} (\texttt{U′ := U″}).
\paragraph{Blueprint notes.}
E.3–E.4, F: \texttt{W}, \texttt{Φ}, \texttt{Ψ}, iso in \texttt{HoP}, \texttt{localGaugeLemma}

Files: \texttt{Forms/LocalGauge.lean}.
\end{proof}

\begin{theorem}[Z5.j — Main theorem]
\label{Z5.j}
\lean{Cohesive.MainStatement, main_theorem, Cohesive.blockPerf_unconditional}
\leanok
Let E be a finite-dimensional complex normed space and X a compact Hausdorff complex manifold modelled on E. Let A = 𝒜\textsubscript{X}(X) be the Dolbeault algebra of X and O\textsubscript{X} its sheaf of holomorphic functions. Then the functor J : Ho(P\textsubscript{A}) ⥤ D(O\textsubscript{X}), M ↦ α(M), into the (standard) derived category of sheaves of O\textsubscript{X}-modules is full and faithful, and its essential image is exactly D\textsubscript{perf}(X, O\textsubscript{X}). The same holds for any other \texttt{HasDerivedCategory} instance.

Deviation: D\textsubscript{perf} instead of D\textsuperscript{b\_}\{coh\}; the equivalence is stated as full, faithful and essential image D\textsubscript{perf}.
\par\smallskip{\footnotesize Source: blueprint-ll.md §5.}
\end{theorem}
\begin{proof}
\leanok
\uses{H1, H2, M, Z4.g, Z5.i}
See the Lean proof and the cited items.
\paragraph{Blueprint notes.}
Block's theorem with no hypotheses (theorem; a new frozen main statement needs your decision), \texttt{\#print axioms}

Files: \texttt{MainH2.lean}, \texttt{ChecksMain.lean}.
\end{proof}

% Generated by scripts/blueprint_gen.py; do not edit.
\chapter{The triangulated structure on Ho(P\textsubscript{A})}\label{chap:Z2}
\textbf{Targets.}

\begin{itemize}
\item \textbf{Z2.} For every bounded graded-commutative curved dga \texttt{𝒜} (the project's \texttt{GCDga}, curvature \texttt{c} allowed), \texttt{HoP 𝒜} gets:
\begin{itemize}
\item a shift \texttt{HasShift (HoP 𝒜) ℤ};
\item a zero object;
\item \texttt{Pretriangulated (HoP 𝒜)}, with distinguished triangles those isomorphic to cone triangles;
\item \texttt{IsTriangulated (HoP 𝒜)} (octahedron).
This is [B, Prop 2.4.1]: "P\textsubscript{A} is pretriangulated in the sense of Bondal–Kapranov, so Ho P\textsubscript{A} is triangulated". Block's proof is the single line "same as Prop 1 and Prop 2 of [BK]".

\end{itemize}
\item \textbf{K4} (blueprint §9, original form). For the unit module \texttt{𝒜} over the Dolbeault algebra (\texttt{Q = 𝒜⁰} in degree 0, \texttt{𝕏 = ∂̄}): \texttt{Hom\_\{HoP\}(𝒜, 𝒜⟦k⟧) ≅ H\textasciicircum{}\{0,k\}\_∂̄(X)} as ℂ-vector spaces, for every \texttt{k : ℤ}. Both sides vanish for \texttt{k < 0} and \texttt{k > dim X}.
\end{itemize}
\textbf{Route} (research of 1 Nov 2026). Direct construction on \texttt{HoP}, imitating Mathlib's \texttt{HomotopyCategory} files: \texttt{Shift} (320 lines), \texttt{MappingCone} (671), \texttt{Pretriangulated} (534) and \texttt{Triangulated} (229). The other options fail:

\begin{itemize}
\item Mathlib has no dg-categories or dg-modules.
\item The total-complex functor is not full, and curvature rules out ordinary complexes.
\item Transport along \texttt{J : HoP ≅ D\_perf} only works for the Dolbeault algebra, needs the main theorem, and gives a non-canonical shift. K4 would then not be a check of the shift.
\end{itemize}
Membership in \texttt{Hom\textasciicircum{}n}: \texttt{φ(α·x) = (−1)\textasciicircum{}\{n|α|\} α·φ(x)}. Differential: \texttt{δφ = 𝕏 ∘ φ − (−1)\textasciicircum{}n φ ∘ 𝕏} (A4, equal to Mathlib's \texttt{HomComplex.δ} by A4.d).

\begin{itemize}
\item \textbf{Shift.}
\begin{itemize}
\item \texttt{Q[a]\textasciicircum{}p := Q\textasciicircum{}\{p+a\}}.
\item \texttt{ι\_a : Tot(Q[a]) → Tot(Q)}, \texttt{α ⊗ e ↦ (−1)\textasciicircum{}\{a|α|\} α ⊗ e} (\texttt{GCDga.twist a} ⊗ reindexing by \texttt{Equiv.addRight a}). It is bijective and lies in \texttt{Hom\textasciicircum{}\{a\}}: an element of degree \texttt{m} goes to degree \texttt{m + a}, and \texttt{ι\_a(β x) = (−1)\textasciicircum{}\{a|β|\} β ι\_a(x)}.
\item \texttt{𝕏[a] := (−1)\textasciicircum{}a ι\_a⁻¹ ∘ 𝕏 ∘ ι\_a}. Checks:
\begin{itemize}
\item Leibniz: \texttt{ι\_a⁻¹(dβ · y) = (−1)\textasciicircum{}\{a(|β|+1)\} dβ ι\_a⁻¹ y}; the signs cancel to \texttt{𝕏[a](βx) = dβ·x + (−1)\textasciicircum{}\{|β|\} β 𝕏[a]x}.
\item Flatness: \texttt{𝕏[a]² = ι⁻¹ c ι = c}, since \texttt{c} has even degree.
\item On generators: \texttt{𝕏[a]\textasciicircum{}j = (−1)\textasciicircum{}\{a(j+1)\} 𝕏\textasciicircum{}j}, consistent with Mathlib's \texttt{n.negOnePow} on differentials.
\end{itemize}
\item Generic lemma (Z2.a): if \texttt{u : Tot' → Tot} is bijective, \texttt{u ∈ Hom\textasciicircum{}m} (semilinear of degree \texttt{m}), and \texttt{𝕏} is a cohesive operator, then \texttt{(−1)\textasciicircum{}m u⁻¹ 𝕏 u} is one. This is cleaner than going through \texttt{GenData}.
\item Morphisms: \texttt{φ ↦ ι\_a⁻¹ φ ι\_a}, with \texttt{δ(ι⁻¹hι) = (−1)\textasciicircum{}a ι⁻¹(δh)ι}, so boundaries go to boundaries.
\item \texttt{M[a][b] ≅ M[a+b]} is the identity on generators up to reindexing; the signs agree (\texttt{(−1)\textasciicircum{}\{(a+b)(j+1)\}}).
\end{itemize}
\item **Cone of \texttt{φ ∈ Z0(M, N)}.**
\begin{itemize}
\item \texttt{Q\_C\textasciicircum{}p := Q\_N\textasciicircum{}p × Q\_M\textasciicircum{}\{p+1\}}.
\item Structure maps, all A-semilinear and defined on generators via \texttt{homEquivGen} (A4.c): \texttt{inl ∈ Hom\textasciicircum{}\{−1\}(M, C)}, \texttt{inr ∈ Hom\textasciicircum{}0(N, C)}, \texttt{fst ∈ Hom\textasciicircum{}1(C, M)}, \texttt{snd ∈ Hom\textasciicircum{}0(C, N)}.
\item Identities: \texttt{fst∘inl = id}, \texttt{snd∘inr = id}, \texttt{fst∘inr = 0}, \texttt{snd∘inl = 0}, \texttt{inl∘fst + inr∘snd = id}.
\item Operator: \texttt{𝕏\_C := inr 𝕏\_N snd + inr φ fst − inl 𝕏\_M fst}. On \texttt{inr f + inl e} this is Block's matrix. Flatness from those of \texttt{M}, \texttt{N} and \texttt{δφ = 0}.
\item Then \texttt{δ(inl) = inr ∘ φ}, \texttt{δ(inr) = 0}, \texttt{δ(fst) = 0}, \texttt{δ(snd) = −φ ∘ fst} (Mathlib's \texttt{δ\_inl}, \texttt{δ\_snd}).
\end{itemize}
\item \textbf{Triangle.} \texttt{M →φ N →inr C(φ) →(−ι₁⁻¹∘fst) M⟦1⟧}, where \texttt{ι₁⁻¹ ∘ fst ∈ Z0(C, M[1])}. This is Mathlib's \texttt{Cocycle.homOf ((-fst φ).rightShift 1 0 \_)}.
\end{itemize}

\begin{theorem}[Z2.a — Shift]
\label{Z2.a}
\lean{Cohesive.shiftQ, Cohesive.shiftIota, Cohesive.QuasiCohesive.transport, Cohesive.CohesiveModule.shift, Cohesive.shiftIota_mem_hom, Cohesive.δ_shiftIota}
\leanok
Let A be a bounded graded-commutative curved dga over k and M = (Q, 𝕏) a cohesive module, a ∈ ℤ. The shifted generators are Q[a]\textsuperscript{p} = Q\textsuperscript{p+a}, and ι\textsubscript{a} : Tot Q[a] → Tot Q, α ⊗ e ↦ (−1)\textsuperscript{a|α|} α ⊗ e, is a k-linear bijection taking total degree m to m + a with ι\textsubscript{a}(βx) = (−1)\textsuperscript{a|β|} β ι\textsubscript{a}(x).

\begin{itemize}
\item Generally, if u : Tot Q′ → Tot Q is such a bijection of degree m, then (−1)\textsuperscript{m} u\textsuperscript{−1} 𝕏 u is a quasi-cohesive structure on Q′.
\item The shift M[a] = (Q[a], (−1)\textsuperscript{a} ι\textsubscript{a}\textsuperscript{−1} 𝕏 ι\textsubscript{a}) is a cohesive module, and ι\textsubscript{a} ∈ Hom\textsuperscript{a}(M[a], M) is closed.
\end{itemize}
Deviation: Block's 𝔼[1] = −𝔼 fails the Leibniz rule for left modules; the shift is transported along ι\textsubscript{a} instead.
\par\smallskip{\footnotesize Source: blueprint-tr.md §5.}
\end{theorem}
\begin{proof}
\leanok
\uses{A4}
See the Lean proof and the cited items.
\paragraph{Blueprint notes.}
\texttt{ι\_a}, generic transport lemma, \texttt{CohesiveModule.shift}

Files: \texttt{Algebra/Shift.lean}.
\end{proof}

\begin{theorem}[Z2.b — Shift functor on Ho(P\textsubscript{A})]
\label{Z2.b}
\lean{Cohesive.shiftMkCoreHoP, Cohesive.isShiftData_shiftIota, Cohesive.IsShiftData.trans, Cohesive.shiftCompZ0, Cohesive.shiftCompIso, Cohesive.conjL}
\leanok
Let A be a bounded graded-commutative curved dga over k. A \emph{shift datum} of degree m from M′ to M is a k-linear bijection v : Tot M′ → Tot M taking total degree n onto n + m, with v(βx) = (−1)\textsuperscript{m|β|} β v(x) and 𝕏′ = (−1)\textsuperscript{m} v\textsuperscript{−1} 𝕏 v; shift data compose (degrees add), and two shift data v₁, v₂ of the same degree onto M give an isomorphism v₂\textsuperscript{−1} ∘ v₁ in Ho(P\textsubscript{A}).

\begin{itemize}
\item Ho(P\textsubscript{A}) has a shift by ℤ: M⟦a⟧ = M[a] on objects, [φ]⟦a⟧ = [ι\textsubscript{a}\textsuperscript{−1} ∘ φ ∘ ι\textsubscript{a}] on morphisms, with M[0] ≅ M and M[a+b] ≅ M[a][b] the comparisons of the shift data ι\textsubscript{0}, id and ι\textsubscript{a+b}, ι\textsubscript{a} ∘ ι\textsubscript{b}.
\item Each shift functor is additive and k-linear.
\end{itemize}
\par\smallskip{\footnotesize Source: blueprint-tr.md §5.}
\end{theorem}
\begin{proof}
\leanok
\uses{A5, Z2.a}
See the Lean proof and the cited items.
\paragraph{Blueprint notes.}
shift on \texttt{Hom}/\texttt{Z0}/\texttt{H0}, \texttt{HasShift (HoP 𝒜) ℤ} (\texttt{hasShiftMk}), additivity, \texttt{k}-linearity

Files: \texttt{Algebra/ShiftHoP.lean}.
\end{proof}

\begin{theorem}[Z2.c — Zero object]
\label{Z2.c}
\lean{Cohesive.prodQ, Cohesive.dsumProdEquiv, Cohesive.zeroCohesive, Cohesive.zeroHoP, Cohesive.isZero_zeroHoP}
\leanok
Let A be a bounded graded-commutative curved dga over k.

\begin{itemize}
\item For families Q₁, Q₂ of A\textsuperscript{0}-modules, with (Q₁ × Q₂)\textsuperscript{p} = Q₁\textsuperscript{p} × Q₂\textsuperscript{p}, the map α ⊗ (e₁, e₂) ↦ (α ⊗ e₁, α ⊗ e₂) is an A-linear isomorphism Tot(Q₁ × Q₂) ≅ Tot Q₁ × Tot Q₂ preserving total degree componentwise.
\item The module with all Q\textsuperscript{p} = 0 and 𝕏 = 0 is cohesive and is a zero object of Ho(P\textsubscript{A}); in particular Ho(P\textsubscript{A}) has a zero object.
\end{itemize}
\par\smallskip{\footnotesize Source: blueprint-tr.md §5.}
\end{theorem}
\begin{proof}
\leanok
\uses{A5}
See the Lean proof and the cited items.
\paragraph{Blueprint notes.}
\texttt{Tot.prodEquiv} (v1.18, below), zero object

Files: \texttt{Algebra/TotProd.lean}, \texttt{Algebra/HoPZero.lean}.
\end{proof}

\begin{definition}[Z2.d — Cone]
\label{Z2.d}
\lean{Cohesive.coneQ, Cohesive.cone, Cohesive.coneX}
\leanok
\uses{Z2.a, Z2.c}
Let A be a bounded graded-commutative curved dga over k, M, N cohesive modules and φ ∈ Z\textsuperscript{0}(M, N) a closed degree-0 map. The cone C(φ) has generators Q\textsubscript{N}\textsuperscript{p} × Q\textsubscript{M}\textsuperscript{p+1}; through Z2.c and ι\textsubscript{1}, Tot C(φ) ≅ Tot M × Tot N, giving inl ∈ Hom\textsuperscript{−1}(M, C), inr ∈ Hom\textsuperscript{0}(N, C), fst ∈ Hom\textsuperscript{1}(C, M), snd ∈ Hom\textsuperscript{0}(C, N) with fst∘inl = id, snd∘inr = id, fst∘inr = 0, snd∘inl = 0, inl∘fst + inr∘snd = id. The operator is 𝕏\textsubscript{C} = inr 𝕏\textsubscript{N} snd + inr φ fst − inl 𝕏\textsubscript{M} fst, and C(φ) is a cohesive module with δ inl = inr ∘ φ, δ inr = 0, δ fst = 0, δ snd = −φ ∘ fst (Mathlib's mapping-cone signs).
\par\smallskip{\footnotesize Source: blueprint-tr.md §5.}
\end{definition}
\paragraph{Blueprint notes.}
cone, structure maps, identities, \texttt{δ} formulas

Files: \texttt{Algebra/Cone.lean}.

\begin{definition}[Z2.e — Maps of cones]
\label{Z2.e}
\lean{Cohesive.coneInrZ0, Cohesive.coneDesc, Cohesive.coneLift, Cohesive.coneMap}
\leanok
\uses{Z2.d}
In the setting of Z2.d (φ ∈ Z\textsuperscript{0}(M, N)), with T a cohesive module:

\begin{itemize}
\item inr : N → C(φ) is closed of degree 0;
\item (desc) for α ∈ Hom\textsuperscript{−1}(M, T) and closed β : N → T of degree 0 with δα = β ∘ φ, α ∘ fst + β ∘ snd : C(φ) → T is closed of degree 0;
\item (lift) for closed α ∈ Hom\textsuperscript{1}(T, M) and β ∈ Hom\textsuperscript{0}(T, N) with δβ = −φ ∘ α, inl ∘ α + inr ∘ β : T → C(φ) is closed of degree 0;
\item (map of cones) for φ\textsubscript{i} ∈ Z\textsuperscript{0}(M\textsubscript{i}, N\textsubscript{i}), closed a : M₁ → M₂, b : N₁ → N₂ of degree 0 and h ∈ Hom\textsuperscript{−1}(M₁, N₂) with δh = b ∘ φ₁ − φ₂ ∘ a, desc(inl₂ ∘ a + inr₂ ∘ h, inr₂ ∘ b) is a closed degree-0 map C(φ₁) → C(φ₂).
\end{itemize}
\par\smallskip{\footnotesize Source: blueprint-tr.md §5.}
\end{definition}
\paragraph{Blueprint notes.}
\texttt{lift}/\texttt{desc}, maps of cones from homotopies

Files: \texttt{Algebra/ConeMaps.lean}.

\begin{theorem}[Z2.f — Ho(P\textsubscript{A}) pretriangulated]
\label{Z2.f}
\lean{Cohesive.distinguishedTriangles, Cohesive.contractible_distinguished, Cohesive.distinguished_cocone_triangle, Cohesive.exists_coneTriangle_hom, Cohesive.complete_distinguished_triangle_morphism, Cohesive.rotateTriangleIso}
\leanok
Let A be a bounded graded-commutative curved dga over k (no further hypotheses). For φ ∈ Z\textsuperscript{0}(M, N) the \emph{cone triangle} is M →[φ] N →[inr] C(φ) →[w\textsubscript{φ}] M⟦1⟧ with w\textsubscript{φ} = −ι\textsubscript{1}\textsuperscript{−1} ∘ fst, and a triangle in Ho(P\textsubscript{A}) is \emph{distinguished} if it is isomorphic to a cone triangle. With the shift of Z2.b, this makes Ho(P\textsubscript{A}) a pretriangulated category:

\begin{itemize}
\item X →[id] X → 0 → X⟦1⟧ is distinguished, and every morphism is the first map of a distinguished triangle;
\item a commutative square between the first maps of two distinguished triangles extends to a morphism of triangles (for cone triangles, by a map of cones, Z2.e);
\item (cone triangle of φ).rotate ≅ cone triangle of inr\textsubscript{φ}, with components id, id and a homotopy equivalence M⟦1⟧ ≅ C(inr\textsubscript{φ}); distinguished triangles are closed under rotation.
\end{itemize}
\par\smallskip{\footnotesize Source: blueprint-tr.md §5.}
\end{theorem}
\begin{proof}
\leanok
\uses{Z2.a, Z2.b, Z2.e}
Notation: \texttt{ι\_a : Tot(M[a]) → Tot(M)} (Z2.a, \texttt{Hom\textasciicircum{}a}, closed), \texttt{w\_φ := −ι₁⁻¹ ∘ fst ∈ Z⁰(C(φ), M[1])} (P14: Mathlib's \texttt{triangle φ}.mor₃ with no extra sign on generators), \texttt{f[a] := ι\_a⁻¹ f ι\_a} (\texttt{conjZ0}, the shift on morphisms, Z2.b). Primes denote the structure maps of \texttt{C(inr\_φ)}. Mathlib's \texttt{leftShift} signs (\texttt{(a n' + a(a−1)/2)}) are absorbed into these formulas.

\begin{itemize}
\item \textbf{Triangle and set} (\texttt{Algebra/ConeTriangle.lean}, Z2.f1). \texttt{coneTriangle φ := (M →[φ] N →[inr] C(φ) →[w\_φ] M⟦1⟧)}; distinguished = isomorphic to some \texttt{coneTriangle φ} (\texttt{φ ∈ Z⁰}, representatives in \texttt{HoP}).
\begin{itemize}
\item contractible: \texttt{C(id\_M)} is zero: \texttt{δ(inl ∘ snd) = id} (\texttt{δ inl ∘ snd − inl ∘ δ snd = inr snd + inl fst}), \texttt{isZero\_of\_contraction};
\item cocone: choose a representative of \texttt{f};
\item morphisms: a square commuting in \texttt{HoP} gives \texttt{h ∈ Hom\textasciicircum{}\{−1\}} with \texttt{δh = bφ₁ − φ₂a} (\texttt{B⁰}); \texttt{coneMap} (Z2.e); \texttt{a[1] ∘ w₁ = w₂ ∘ coneMap} exactly (\texttt{fst₂ ∘ coneMap = a ∘ fst₁}).
\end{itemize}
\item \textbf{Rotation} (\texttt{Algebra/ConeRotate.lean}, Z2.f2; Mathlib \texttt{rotateHomotopyEquiv}). \texttt{hom := coneLift inr\_φ (φ ∘ ι₁) (−inl ∘ ι₁) : M[1] → C(inr\_φ)} (\texttt{δ(−inl ι₁) = −inr φ ι₁}), \texttt{inv := coneDesc inr\_φ 0 w\_φ : C(inr\_φ) → M[1]}; \texttt{inv ∘ hom = id} exactly, \texttt{hom ∘ inv − id = δ H} with \texttt{H := −inl' ∘ snd\_φ ∘ snd'}. Triangle iso \texttt{(coneTriangle φ).rotate ≅ coneTriangle inr\_φ} with components \texttt{(id, id, [hom])}: \texttt{w' ∘ hom = −φ[1]} exactly (\texttt{fst' ∘ hom = φ ι₁}), and \texttt{hom ∘ w\_φ − inr' = δ H₂}, \texttt{H₂ := −inl' ∘ snd\_φ}.
\item \textbf{Shift} (\texttt{Algebra/ConeShift.lean}, Z2.f3; Mathlib \texttt{shiftIso}, \texttt{shiftTriangleIso}). \texttt{C(φ)[n] ≅ C(φ[n])}: \texttt{hom := coneLift φ[n] ((−1)\textasciicircum{}n ι\_n⁻¹ fst ι\_n) (ι\_n⁻¹ snd ι\_n)}, \texttt{inv := coneDesc φ[n] ((−1)\textasciicircum{}n ι\_n⁻¹ inl ι\_n) (ι\_n⁻¹ inr ι\_n)}, inverse exactly; triangle iso \texttt{(Triangle.shiftFunctor n).obj (coneTriangle φ) ≅ coneTriangle φ[n]} with components \texttt{(id, (−1)\textasciicircum{}n, shiftIso)}; the third square involves \texttt{shiftFunctorComm} of \texttt{HoP} (Z2.b's \texttt{shiftFunctorAdd}), the main risk.
\item \textbf{Instance} (\texttt{Algebra/Pretriangulated.lean}): as Mathlib, \texttt{invRotate} via \texttt{invRotateIsoRotateRotateShiftFunctorNegOne} and shift-closure.
\end{itemize}
\paragraph{Blueprint notes.}
the pretriangulated axioms, shift of cones

Files: \texttt{Algebra/Pretriangulated.lean}.
\end{proof}

\begin{theorem}[Z2.g — Octahedron axiom]
\label{Z2.g}
\lean{Cohesive.compTriangle, Cohesive.compTriangle_distinguished, Cohesive.instIsTriangulatedHoP}
\leanok
Let A be a bounded graded-commutative curved dga over k. For closed degree-0 maps u : X₁ → X₂, v : X₂ → X₃, the triangle C(u) →[m₁] C(v ∘ u) →[m₂] C(v) →[w\textsubscript{v} ≫ inr\textsubscript{u}⟦1⟧] C(u)⟦1⟧, with m₁, m₂ the maps of cones of the squares (id, v) and (u, id) with zero homotopy, is distinguished. Hence the pretriangulated category Ho(P\textsubscript{A}) satisfies the octahedron axiom: Ho(P\textsubscript{A}) is triangulated, for every curvature c.
\par\smallskip{\footnotesize Source: blueprint-tr.md §5.}
\end{theorem}
\begin{proof}
\leanok
\uses{Z2.a, Z2.b, Z2.c, Z2.d, Z2.e, Z2.f}
Mathlib \texttt{HomotopyCategory/Triangulated.lean} (\texttt{mappingConeCompTriangle}, \texttt{mappingConeCompHomotopyEquiv}, \texttt{IsTriangulated.mk'}). For \texttt{u ∈ Z⁰(X₁, X₂)}, \texttt{v ∈ Z⁰(X₂, X₃)}, \texttt{uv := v ∘ u}. Conventions as in §3c (\texttt{δ inl = inr φ}, \texttt{δ snd = −φ fst}, \texttt{δ fst = δ inr = 0}, \texttt{δ(f∘g) = δf∘g + (−1)\textasciicircum{}\{|f|\} f∘δg}).

\begin{itemize}
\item \textbf{Comparison maps} (\texttt{Algebra/ConeComp.lean}). \texttt{m₁ := coneMap u uv id v 0 : C(u) → C(uv)}, \texttt{m₂ := coneMap uv v u id 0 : C(uv) → C(v)} (homotopy \texttt{0}: \texttt{v∘u − uv∘id = 0}, \texttt{id∘uv − v∘u = 0}). Triangle \texttt{C(u) →[m₁] C(uv) →[m₂] C(v) →[w\_v ≫ inr\_u⟦1⟧] C(u)⟦1⟧}.
\item \textbf{Homotopy equivalence} \texttt{C(v) ≃ C(m₁) =: D}. \texttt{hom := coneLift m₁ (inr\_u ∘ fst\_v) (inr\_\{uv\} ∘ snd\_v)} (\texttt{δ(inr\_\{uv\} snd\_v) = −inr\_\{uv\} v fst\_v = −m₁ inr\_u fst\_v} by \texttt{m₁ inr\_u = inr\_\{uv\} v}); \texttt{inv := coneDesc m₁ (inl\_v ∘ snd\_u) m₂} (\texttt{δ(inl\_v snd\_u) = inr\_v v snd\_u + inl\_v u fst\_u = m₂ ∘ m₁}, using \texttt{m₁ = inl\_\{uv\} fst\_u + inr\_\{uv\} v snd\_u}). \texttt{inv ∘ hom = id} exactly (\texttt{fst\_v inl\_v = id}, \texttt{snd\_v m₂ = snd\_\{uv\}}); \texttt{hom ∘ inv − id = δH} with \texttt{H := −inl\_D ∘ inl\_u ∘ fst\_\{uv\} ∘ snd\_D} (checked by hand: with \texttt{K = inl\_u fst\_\{uv\} snd\_D}, \texttt{δ(inl\_D K) = inr\_D inl\_\{uv\} fst\_\{uv\} snd\_D − inl\_D inr\_u u fst\_\{uv\} snd\_D + inl\_D inl\_u fst\_u fst\_D}, which is \texttt{−(hom ∘ inv − id)} after expanding \texttt{id\_D}, \texttt{id\_\{C(u)\}}, \texttt{id\_\{C(uv)\}}).
\item \textbf{Distinguished.} Triangle iso \texttt{(id, id, [hom])} to \texttt{coneTriangle m₁}: \texttt{inv ∘ inr\_D = m₂} exactly (so \texttt{m₂ ≫ [hom] = [inr\_D]} in \texttt{HoP}), and \texttt{w\_D ∘ hom = −ι₁⁻¹ inr\_u fst\_v = inr\_u[1] ∘ w\_v} exactly.
\item \textbf{Octahedron} (\texttt{Algebra/Triangulated.lean}). \texttt{IsTriangulated.mk'} with the three cone triangles of \texttt{u}, \texttt{v}, \texttt{uv}, \texttt{Iso.refl}s, and \texttt{Octahedron.mk [m₁] [m₂]}: \texttt{m₁ inr\_u = inr\_\{uv\} v}, \texttt{w\_\{uv\} m₁ = w\_u} (\texttt{fst m₁ = fst}), \texttt{m₂ inr\_\{uv\} = inr\_v}, \texttt{u[1] w\_\{uv\} = w\_v m₂} (\texttt{fst\_v m₂ = u fst\_\{uv\}}), all exact; the last field is the distinguished triangle above.
\end{itemize}
\paragraph{4. K4}
\begin{itemize}
\item \textbf{(i) Generic} (model: Mathlib \texttt{HomComplexCohomology.homAddEquiv}). \texttt{Hom\_\{HoP\}(M, N⟦k⟧) ≅ H\textasciicircum{}k(Hom\textasciicircum{}•(M,N))} via \texttt{φ ↦ ι\_k ∘ φ}, since \texttt{δ(ι\_k ∘ φ) = ± ι\_k ∘ δφ}. The cohomology is the quotient \texttt{(Hom\textasciicircum{}k ⊓ ker δ) ⧸ δ(Hom\textasciicircum{}\{k−1\})} (like \texttt{H0}), or the homology of A6's \texttt{homComplex}.
\item \textbf{(ii) Unit module.} \texttt{Hom\textasciicircum{}k(U, U) ≅ 𝒜\textasciicircum{}k} by evaluation at \texttt{1 ⊗ gen} (\texttt{homEquivGen}, \texttt{hom\_tmul}), with \texttt{δ ↔ d} and no sign (\texttt{θ(gen) = 0}). \texttt{unitCohesive} exists in \texttt{ChecksAlgebra.lean:338}; move it to a library file (it is not frozen).
\item \textbf{(iii) Dolbeault.} Define \texttt{H\textasciicircum{}k(𝒜, d) := (𝒜\textasciicircum{}k ⊓ ker d) ⧸ d(𝒜\textasciicircum{}m)} with \texttt{m + 1 = k} (ℤ-indexed, no ℕ subtraction) and \texttt{H\textasciicircum{}\{0,k\}\_∂̄(X) :=} that for \texttt{dolbeault E X} (\texttt{d = dbarL}). \texttt{K4 : ((⟨unitCohesive (dolbeault E X) rfl⟩ : HoP \_) ⟶ (⟨…⟩ : HoP \_)⟦k⟧) ≃ₗ[ℂ] H\textasciicircum{}\{0,k\}\_∂̄(X)}. Sanity: both sides are \texttt{0} for \texttt{k < 0} and \texttt{k > n}. For \texttt{k = 0} it is \texttt{O(X)} (F12).
\end{itemize}
\paragraph{5. Items, files, estimates (new files only; no FROZEN changes)}
\begin{itemize}
\item Item: Z2.a; File: \texttt{Algebra/Shift.lean}; Content: \texttt{ι\_a}, generic transport lemma, \texttt{CohesiveModule.shift}; Est. lines: 250–350
\item Item: Z2.b; File: \texttt{Algebra/ShiftHoP.lean}; Content: shift on \texttt{Hom}/\texttt{Z0}/\texttt{H0}, \texttt{HasShift (HoP 𝒜) ℤ} (\texttt{hasShiftMk}), additivity, \texttt{k}-linearity; Est. lines: 300–450
\item Item: K4.a; File: \texttt{Algebra/HomShift.lean}; Content: \texttt{Hom\_\{HoP\}(M, N⟦k⟧) ≃ₗ H\textasciicircum{}k(Hom\textasciicircum{}•(M,N))}; Est. lines: 100–150
\item Item: K4.b; File: \texttt{Algebra/UnitModule.lean}, \texttt{ChecksMain}/\texttt{ChecksForms}; Content: unit module, \texttt{Hom\textasciicircum{}•(U,U) ≅ (𝒜, d)}, Dolbeault cohomology, K4; Est. lines: 200–350
\item Item: Z2.c; File: \texttt{Algebra/TotProd.lean}, \texttt{Algebra/HoPZero.lean}; Content: \texttt{Tot.prodEquiv} (v1.18, below), zero object; Est. lines: 100–200
\item Item: Z2.d; File: \texttt{Algebra/Cone.lean}; Content: cone, structure maps, identities, \texttt{δ} formulas; Est. lines: 350–500
\item Item: Z2.e; File: \texttt{Algebra/ConeMaps.lean}; Content: \texttt{lift}/\texttt{desc}, maps of cones from homotopies; Est. lines: 200–300
\item Item: Z2.f; File: \texttt{Algebra/Pretriangulated.lean}; Content: the pretriangulated axioms, shift of cones; Est. lines: 400–600
\item Item: Z2.g; File: \texttt{Algebra/Triangulated.lean}; Content: octahedron, \texttt{IsTriangulated}; Est. lines: 250–400
\end{itemize}
Total: Z2 core ≈ 1,800–2,700 lines, K4 ≈ 300–500 lines. Order: Z2.a → Z2.b → K4, which needs no cones, then Z2.c–g. This can run in parallel with Z5 (disjoint files: Algebra layer vs Analysis and Forms).

Optional, not proposed now: Z2.h, \texttt{J.CommShift} and \texttt{J.IsTriangulated}. With them \texttt{Perf = J.essImage} is triangulated by Mathlib's \texttt{essImage} instance (+400–700 lines).

\paragraph{6. Risks}
\begin{itemize}
\item **Coherence of \texttt{HasShift}** (\texttt{M[a][b] ≅ M[a+b]}, \texttt{zero\_add}, \texttt{assoc}). This is dependent-type reindexing of \texttt{Q (p + a + b)} vs \texttt{Q (p + (a + b))}. Use a \texttt{reindex (h : f = g)} family with \texttt{subst}, as Mathlib's \texttt{XIsoOfEq}.
\item \textbf{Cone generators.} \texttt{⨁ (Q\_N × Q\_M[1])} is not a type-level biproduct of the \texttt{Tot}s, so the five identities need generator-level proofs (\texttt{homEquivGen}).
\item \textbf{Base-change elaboration traps} (PROGRESS) do not arise here; \texttt{HoP} is over a fixed \texttt{𝒜}.
\end{itemize}
\paragraph{7. Lean targets (to be compiled in \texttt{Scratch/P14\_TR.lean} before skeletons)}
\begin{verbatim}
def CohesiveModule.shift (M : CohesiveModule 𝒜) (a : ℤ) : CohesiveModule 𝒜
instance : HasShift (HoP 𝒜) ℤ
instance : HasZeroObject (HoP 𝒜)
instance (n : ℤ) : (shiftFunctor (HoP 𝒜) n).Additive
instance : Pretriangulated (HoP 𝒜)
instance : IsTriangulated (HoP 𝒜)
def homShiftEquiv (M N : HoP 𝒜) (k : ℤ) : (M ⟶ N⟦k⟧) ≃ₗ[k] HomCohomology M.as N.as k
def dolbeaultCohomology (k : ℤ) : Type
theorem K4 (k : ℤ) : Nonempty ((unitHoP (dolbeault E X) ⟶ (unitHoP (dolbeault E X))⟦k⟧) ≃ₗ[ℂ]
    dolbeaultCohomology E X k)
\end{verbatim}
\paragraph{Blueprint notes.}
octahedron, \texttt{IsTriangulated}

Files: \texttt{Algebra/Triangulated.lean}.
\end{proof}

% Generated by scripts/blueprint_gen.py; do not edit.
\chapter{The twisted theorem}\label{chap:Z3}
\textbf{Target.} For a compact Hausdorff complex manifold \texttt{X}, a \texttt{∂̄}-closed \texttt{h ∈ 𝒜\textasciicircum{}\{0,2\}(X)} and a twist datum \texttt{T} for \texttt{h}, the functor

\texttt{J\_T : Ho(P\_\{(𝒜\_X(X), ∂̄, h)\}) ⥤ D(TwMod O 𝒰 ȟ)}, \texttt{M ↦ ((α\textasciicircum{}h(M)|U\_α, 𝕏 − h⁰\_α), u\_αβ·)\_α}

is full and faithful with essential image the twisted perfect complexes \texttt{TwPerf}; and every closed \texttt{h} has a twist datum. Closest sources: [A, Thm 7.4.3] (first equivalence, compact \texttt{X}), [B, Thm 4.2.1] (stated without proof; \texttt{D\textasciicircum{}b\_coh} replaced by \texttt{D\_perf}, which is Z1).

\textbf{Shape of the proof.} On each \texttt{U\_α} the curved algebra is gauge-equivalent to the flat one (\texttt{𝕏 ↦ 𝕏 − h⁰\_α}), so every over-site lemma of C1/C2/C4/C7 and H2 on opens is reused on the cover pieces. Twisted sheaves have no global sections, so every global step (softness S3.b, globalization S10, lifting A11) acts on untwisted objects: Hom sheaves, a glued \texttt{𝒜⁰}-complex, global sections of a strictly perfect complex. New global work: the abelian category \texttt{TwMod} and its homological algebra, gluing of sheaves of modules (S12), twisted Hom sheaves, the glued curved model, A11 over the curved algebra (already general, §1.1).

\textbf{Size.} ≈ 12–21k lines (H2 was ≈ 6k); checked statement (statement closure, \texttt{exists\_twistData}, K8) ≈ 5–8k. Three stages, gates G1–G5, phases 8–10 (§5).

\begin{definition}[Z3.a — Curved dga/presheaf]
\label{Z3.a}
\lean{Cohesive.GCDga.withC, Cohesive.DgaPresheaf.withCurvature, Cohesive.DgaSheafData, Cohesive.PreResolvingData.toSheafData}
\leanok
\uses{C0}
Let A be a bounded graded-commutative curved dga (A1) with curvature c. For a closed element c′ ∈ A\textsuperscript{2} (dc′ = 0), A with curvature c′ is the curved dga with the same algebra, grading and differential and curvature c′. For a presheaf of dgas D on X and a closed h ∈ A(X)\textsuperscript{2}, the \emph{curved presheaf} D\textsubscript{h} has the same algebras, gradings and restriction maps, with curvature h|U on each open U. \emph{Sheaf data} on X are the fields of pre-resolving data (C0) except c = 0: a presheaf of possibly curved dgas with gluing and degree-locality of sections, 𝒜\textsuperscript{0} a fine sheaf of commutative rings, a structure sheaf O with ρ : O → 𝒜\textsuperscript{0} and d ∘ ρ = 0, and every 𝒜\textsuperscript{q} locally free over 𝒜\textsuperscript{0}. Every pre-resolving data are sheaf data, and sheaf data with D replaced by D\textsubscript{h} are again sheaf data.
\par\smallskip{\footnotesize Source: blueprint-tw.md §4.}
\end{definition}
\paragraph{Blueprint notes.}
\texttt{withC} (+ rfl API), \texttt{withCurvature}, \texttt{DgaSheafData}, \texttt{toSheafData}, bridges

Files: \texttt{Algebra/CurvedDga.lean}, \texttt{Comparison/CurvedPresheaf.lean}.

\begin{theorem}[Z3.b — Gauge]
\label{Z3.b}
\lean{Cohesive.CohesiveModule.gauge, Cohesive.QuasiCohesive.ungauge, Cohesive.gauge_hom, Cohesive.gauge_δ, Cohesive.gauge_Z0, Cohesive.gauge_B0}
\leanok
Let k be a commutative ring in which 2 is a unit, A a bounded graded-commutative curved dga over k with curvature c, β ∈ A\textsuperscript{1}, and c′ = c + dβ. For a (quasi-)cohesive module (Q, 𝕏) over (A, c′), its \emph{gauge} (Q, 𝕏 − β·) is a (quasi-)cohesive module over (A, c), with (𝕏 − β)\textsuperscript{2} = c′ − dβ + β\textsuperscript{2} = c; the inverse is 𝕏 ↦ 𝕏 + β·. The gauge does not change the Hom complexes: Hom\textsuperscript{n}(M₁, M₂) is the same submodule, δ agrees on it, hence so do the closed and exact degree-0 maps. Consequently the gauge is an equivalence Ho(P\textsubscript{(A, c + dβ)}) ≌ Ho(P\textsubscript{(A, c)}), the identity on morphisms.

Deviation: the sources write "2 ≠ 0"; Lean assumes 2 is a unit in k (true for k = ℂ).
\par\smallskip{\footnotesize Source: blueprint-tw.md §4.}
\end{theorem}
\begin{proof}
\leanok
\uses{Z3.a}
\textbf{A.2} gauge curvature.

\textbf{A.3} gauge on Hom: identity on \texttt{Hom\textasciicircum{}n}, \texttt{δ}, \texttt{Z0}, \texttt{B0}, composition.

\textbf{A.4} gauge commutes with base change.
\paragraph{Blueprint notes.}
A.2–A.4, \texttt{gaugeHoP} (explicit \texttt{@}-projections, P15)

Files: \texttt{Algebra/Gauge.lean}, \texttt{Algebra/GaugeBaseChange.lean}.
\end{proof}

\begin{theorem}[Z3.c — Curved Dolbeault, twist data]
\label{Z3.c}
\lean{Cohesive.DolbeaultClosed2, Cohesive.curvedDolbeaultPresheaf, Cohesive.curvedDolbeault, Cohesive.TwistCocycle, Cohesive.TwistCocycle.trivial, Cohesive.exactTwistDatum}
\leanok
Let X be a complex manifold modelled on a finite-dimensional complex normed space E.

\begin{itemize}
\item A \emph{closed (0,2)-form} is a global Dolbeault form h ∈ 𝒜\textsubscript{X}(X) of form degree 2 with ∂̄h = 0. The curved Dolbeault presheaf is U ↦ (𝒜\textsubscript{X}(U), ∂̄, h|U), and the curved Dolbeault algebra is (𝒜\textsubscript{X}(X), ∂̄, h), whose curvature is h.
\item A \emph{twist cocycle} for a sheaf of commutative rings R is a finite open cover (U\textsubscript{α}) of X with units ȟ\textsubscript{αβγ} ∈ R(U\textsubscript{α} ∩ U\textsubscript{β} ∩ U\textsubscript{γ})\textasciicircum{}×, equal to 1 when two indices coincide, satisfying ȟ\textsubscript{βγδ} ȟ\textsubscript{αγδ}\textsuperscript{−1} ȟ\textsubscript{αβδ} ȟ\textsubscript{αβγ}\textsuperscript{−1} = 1 on every open contained in U\textsubscript{α}, U\textsubscript{β}, U\textsubscript{γ}, U\textsubscript{δ}. The trivial cocycle is the cover \{X\} with ȟ = 1.
\item A \emph{twist datum} for a closed h ∈ 𝒜(X)\textsuperscript{2} over sheaf data S is a twist cocycle ȟ of S.O together with primitives h\textsuperscript{0}\textsubscript{α} ∈ 𝒜\textsuperscript{1}(U\textsubscript{α}), dh\textsuperscript{0}\textsubscript{α} = h|U\textsubscript{α}, and units u\textsubscript{αβ} ∈ 𝒜\textsuperscript{0}(U\textsubscript{α} ∩ U\textsubscript{β})\textasciicircum{}× with du\textsubscript{αβ} = (h\textsuperscript{0}\textsubscript{β} − h\textsuperscript{0}\textsubscript{α}) u\textsubscript{αβ}, u\textsubscript{αα} = 1, u\textsubscript{βα} = u\textsubscript{αβ}\textsuperscript{−1}, and ρ(ȟ\textsubscript{αβγ}) = u\textsubscript{βγ} u\textsubscript{αβ} u\textsubscript{αγ}\textsuperscript{−1}. For h = dβ with β ∈ 𝒜\textsuperscript{1}(X) there is a twist datum on the cover \{X\} with h\textsuperscript{0} = β, u = 1, ȟ = 1 (for β = 0: the trivial datum for h = 0).
\end{itemize}
Deviation: multiplicative data (h\textsuperscript{0}, u, ȟ) on a finite cover instead of Antweiler's (h\textsuperscript{0}, h\textsuperscript{1}, h\textsuperscript{2}, exp); the normalization and cocycle condition of ȟ are imposed as fields.
\par\smallskip{\footnotesize Source: blueprint-tw.md §4.}
\end{theorem}
\begin{proof}
\leanok
\uses{F14, Z3.a}
\begin{itemize}
\item \texttt{withC (i : GCDga 𝒜) (c') (hc : c' ∈ 𝒜 2) (hd : i.d c' = 0) : GCDga 𝒜 := \{ i with c := c', … \}}, \texttt{@[reducible]}, never an instance. \texttt{d}, \texttt{N}, \texttt{commRing0}, \texttt{Tot}, \texttt{Hom} agree by \texttt{rfl} (P15).
\item \texttt{DgaPresheaf.withCurvature D h hh hd}: same \texttt{A}, \texttt{𝒜}, grading; \texttt{dga U := withC (D.dga U) (D.res le\_top h) …}; \texttt{res} the same algebra maps with \texttt{map\_c} from \texttt{res\_comp} (built with explicit \texttt{@DgaHom.mk}, P15).
\item \texttt{DgaSheafData X} (u): \texttt{PreResolvingData} without \texttt{c\_eq\_zero}; \texttt{PreResolvingData.toSheafData}, \texttt{DgaSheafData.withCurvature}.
\item \textbf{Gauge} (A.2–A.4). For \texttt{M} over \texttt{withC i c'} and \texttt{β ∈ 𝒜¹} with \texttt{c' = c + dβ}: \texttt{gauge M β := (M.Q, 𝕏 − β·)} over \texttt{i}. \texttt{(𝕏 − β)² = c' − dβ + β² = c} by Leibniz \texttt{𝕏(βx) = dβ·x − β𝕏x}, and \texttt{β² = 0} since \texttt{β² = −β²} (graded commutativity). \warn{} Generic \texttt{k} needs no 2-torsion: hypothesis \texttt{IsUnit (2 : k)} in the generic gauge lemma (true for ℂ); deviation from A/B, which wrote "2 ≠ 0". Same \texttt{Hom\textasciicircum{}n}, same \texttt{δ} (the \texttt{β}-terms cancel since \texttt{φ(βx) = (−1)\textasciicircum{}n βφ(x)}), so \texttt{gaugeHoP : HoP (withC i c') ≌ HoP i}, identity on \texttt{H0}. \texttt{(gauge M β).baseChange f = gauge (M.baseChange f) (f β)} on generators. Dictionary: [B, §2.7] \texttt{(id, ω = β)}, sign reversed for left modules.
\item \texttt{DolbeaultClosed2 E X := \{h : DolbeaultAlg E ⊤ // h ∈ dolbeaultGrade E ⊤ 2 ∧ ∂̄h = 0\}}; \texttt{curvedDolbeaultPresheaf E X h := (dolbeaultPresheaf E X).withCurvature h.1 …}; \texttt{curvedDolbeault E X h := (curvedDolbeaultPresheaf E X h).𝒜 ⊤}.
\end{itemize}
\texttt{TwistCocycle R} (u): \texttt{ι : Type}, \texttt{[Fintype ι]}, \texttt{U : ι → Opens X}, \texttt{⨆ U = ⊤}, \texttt{ch α β γ : (R(U\_αβγ))ˣ}, normalized, Čech cocycle on quadruple intersections. \texttt{TwistDatum S h extends TwistCocycle S.O} (u): \texttt{h0 α ∈ 𝒜¹(U\_α)} with \texttt{d h0\_α = h|U\_α}; \texttt{u α β : 𝒜⁰(U\_αβ)ˣ} with \texttt{d u\_αβ = (h0\_β − h0\_α) u\_αβ}, \texttt{u\_αα = 1}, \texttt{u\_βα = u\_αβ⁻¹}; \texttt{ρ(ch\_αβγ) = u\_βγ u\_αβ u\_αγ⁻¹} on \texttt{U\_αβγ}. \texttt{TwistData E X h := TwistDatum (dolbeaultPreData E X) .toSheafData h.1}. Normalization and cocycle of \texttt{ch} follow from (4), (3) and injectivity of \texttt{ρ} (B correction 8). Special data: \texttt{trivialTwistDatum} (\texttt{h = 0}: \texttt{ι = Unit}, \texttt{h0 = 0}, \texttt{u = 1}, \texttt{ch = 1}), \texttt{exactTwistDatum h β} (\texttt{h = dβ}: \texttt{h0 = β}, \texttt{u = 1}, \texttt{ch = 1}).
\paragraph{Blueprint notes.}
§2.1–2.2, \texttt{trivialTwistDatum}, \texttt{exactTwistDatum}

Files: \texttt{Forms/CurvedDolbeault.lean}, \texttt{Sheaf/Twisted/CocycleDef.lean}, \texttt{Comparison/TwistDatumDef.lean}.
\end{proof}

\begin{theorem}[Z3.d — Exp, existence]
\label{Z3.d}
\lean{Cohesive.expForm, Cohesive.expUnit, Cohesive.dbar_expForm, Cohesive.exists_cover_primitives, Cohesive.twistData_of_primitives, Cohesive.exists_twistData}
\leanok
Let X be a compact Hausdorff complex manifold modelled on a finite-dimensional complex normed space E, and h ∈ 𝒜\textsuperscript{0,2}(X) with ∂̄h = 0. Then h has a twist datum. More precisely:

\begin{itemize}
\item there are a finite open cover (U\textsubscript{α}) of X, primitives h\textsuperscript{0}\textsubscript{α} ∈ 𝒜\textsuperscript{1}(U\textsubscript{α}) of h|U\textsubscript{α}, and functions η\textsubscript{αβ} ∈ 𝒜\textsuperscript{0}(U\textsubscript{α} ∩ U\textsubscript{β}) with ∂̄η\textsubscript{αβ} = h\textsuperscript{0}\textsubscript{β} − h\textsuperscript{0}\textsubscript{α}, η\textsubscript{αα} = 0 and η\textsubscript{βα} = −η\textsubscript{αβ};
\item for any such data, u\textsubscript{αβ} = exp(η\textsubscript{αβ}) (a unit of 𝒜\textsuperscript{0} with ∂̄ exp(η) = exp(η)·∂̄η) and the holomorphic units ȟ\textsubscript{αβγ} with ρ(ȟ\textsubscript{αβγ}) = u\textsubscript{βγ} u\textsubscript{αβ} u\textsubscript{αγ}\textsuperscript{−1} form a twist datum for h.
\end{itemize}
\par\smallskip{\footnotesize Source: blueprint-tw.md §4.}
\end{theorem}
\begin{proof}
\leanok
\uses{F12, H1, Z3.c, Z4.d, Z4.g, Z5.g}
\textbf{B.1} H1 (\texttt{q = 1}) + finite subcover: \texttt{U\_α}, \texttt{h0\_α}.

\textbf{B.2} \texttt{ω\_αβ := h0\_β − h0\_α} closed, antisymmetric.

\textbf{B.3} refinement: shrink (\texttt{exists\_iUnion\_eq\_closure\_subset}), at each \texttt{x} intersect H1 (\texttt{q = 0}) neighbourhoods in \texttt{U\_αβ} or complements of \texttt{cl U'\_α ∩ cl U'\_β}; finitely many \texttt{W\_i}; cover \texttt{U'\_α ⊓ W\_i} (index \texttt{ι × Fin m}): \texttt{ω} exact on all double intersections; symmetrized primitives \texttt{η'}.

\textbf{B.4} \texttt{u := exp ∘ η'} (\texttt{HasFDerivAt.cexp}, \texttt{ContDiff.cexp}; \texttt{∂̄} passes the holomorphic factor through, Z4.d \texttt{dbarPart\_comp\_clm}).

\textbf{B.5} \texttt{ζ = u\_bc u\_ab u\_ac⁻¹} is \texttt{∂̄}-closed, so \texttt{ζ = ρ(ch)} (F12, \texttt{ρ} injective).

\textbf{B.6} cocycle, normalization, cyclic symmetry \texttt{ch\_bca = ch\_abc} (used in C.4).
\paragraph{Blueprint notes.}
B.1–B.6

Files: \texttt{Forms/ExpForm.lean}, \texttt{Forms/TwistExists.lean}.
\end{proof}

\begin{theorem}[Z3.e — H2 on opens, LL\textsuperscript{h}]
\label{Z3.e}
\lean{Cohesive.curvedRes, Cohesive.curvedRestrict, Cohesive.curvedLL, Cohesive.localGaugeLemmaOpen, Cohesive.exists_free_baseChange_open}
\leanok
\uses{Z5.g, Z5.h, Z5.i}
Let X be a complex manifold modelled on a finite-dimensional complex normed space E.

\begin{itemize}
\item (H2 over an open.) For an open U₀ ⊆ X, a cohesive module N over 𝒜\textsubscript{X}(U₀) and x ∈ U₀, there are an open U′ ⊆ U₀ containing x, a bounded complex W of finitely generated projective O\textsubscript{X}(U′)-modules and an isomorphism N|U′ ≅ 𝒜\textsubscript{X}(U′) ⊗\textsubscript{O\_X(U′)} W in Ho(P\textsubscript{𝒜\_X(U′)}).
\item (Curved local lifting.) Let X moreover be compact Hausdorff, h ∈ 𝒜\textsuperscript{0,2}(X) closed, T a twist datum for h and M a cohesive module over (𝒜\textsubscript{X}(X), ∂̄, h). For x ∈ U\textsubscript{α} the same conclusion holds for the gauged restriction gauge\_\{h\textsuperscript{0}\textsubscript{α}\}(M|U\textsubscript{α}), a cohesive module over the flat algebra 𝒜\textsubscript{X}(U\textsubscript{α}) with operator 𝕏 − h\textsuperscript{0}\textsubscript{α}.
\end{itemize}
Deviation: Antweiler states the curved lemma for M|U′ itself; it is stated after the gauge by h\textsuperscript{0}\textsubscript{α}, which is an equivalence of homotopy categories (Z3.b).
\par\smallskip{\footnotesize Source: blueprint-tw.md §4.}
\end{theorem}
\begin{proof}
\leanok
\uses{Z3.b, Z3.c}
See the Lean proof and the cited items.
\paragraph{Blueprint notes.}
Z5.g–i with base \texttt{U₀}; \texttt{localGaugeLemma} re-derived at \texttt{⊤}; \texttt{curvedLL}

Files: \texttt{Forms/LocalFreeOpen.lean}, \texttt{Forms/LocalGaugeOpen.lean}, \texttt{Forms/CurvedLL.lean}.
\end{proof}

\begin{definition}[Z3.f — \texttt{restrOver}]
\label{Z3.f}
\lean{Cohesive.restrOver, Cohesive.restrOverAdj, Cohesive.restr_restrOver}
\leanok
\uses{S1}
Let R be a sheaf of rings on X and V ⊆ U opens. Restriction M ↦ M|V from sheaves of R|U-modules on the over-site of U to sheaves of R|V-modules on the over-site of V is Mathlib's pushforward along Over V → Over U. It is left adjoint to the pushforward along the other direction (which therefore preserves injective objects), it also has a left adjoint, so it is exact, and restricting to U and then to V is isomorphic to restricting to V.
\par\smallskip{\footnotesize Source: blueprint-tw.md §4.}
\end{definition}
\paragraph{Blueprint notes.}
P17 promoted: adjoints, exactness, \texttt{restr\_restrOver}, injectives

Files: \texttt{Sheaf/RestrOver.lean}.

\begin{definition}[Z3.g — \texttt{TwMod}]
\label{Z3.g}
\lean{Cohesive.sec, Cohesive.toSec, Cohesive.secRes, Cohesive.TwMod, Cohesive.TwMod.Hom, Cohesive.TwMod.comp}
\leanok
\uses{S1}
Let R be a sheaf of commutative rings on X, (U\textsubscript{α})\textsubscript{α∈ι} a family of opens and ȟ\textsubscript{αβγ} ∈ R(U\textsubscript{α} ∩ U\textsubscript{β} ∩ U\textsubscript{γ})\textasciicircum{}×. An \emph{ȟ-twisted sheaf of R-modules} consists of sheaves of modules M\textsubscript{α} on the over-sites of U\textsubscript{α}, and for every open W ⊆ U\textsubscript{α} ∩ U\textsubscript{β} an additive map φ\textsubscript{αβ} : M\textsubscript{α}(W) → M\textsubscript{β}(W), R(W)-linear and compatible with restriction, with φ\textsubscript{αα} = id and φ\textsubscript{βγ} ∘ φ\textsubscript{αβ} = ȟ\textsubscript{αβγ}|W · φ\textsubscript{αγ} on every W ⊆ U\textsubscript{α} ∩ U\textsubscript{β} ∩ U\textsubscript{γ}. Morphisms are families of morphisms M\textsubscript{α} → N\textsubscript{α} commuting with the gluing maps; composition is componentwise, and M ↦ M\textsubscript{α} is the component functor.

Deviation: the gluing is given on sections over every W ⊆ U\textsubscript{α} ∩ U\textsubscript{β} instead of as isomorphisms of restricted sheaves; φ\textsubscript{αα} = id is a field.
\par\smallskip{\footnotesize Source: blueprint-tw.md §4.}
\end{definition}
\paragraph{Blueprint notes.}
P16 promoted: \texttt{sec}, \texttt{toSec}, \texttt{secRes}, structure, \texttt{Hom}, \texttt{Category}, \texttt{comp}, \texttt{φ\_self}, \texttt{φ\_symm}

Files: \texttt{Sheaf/Twisted/Def.lean}.

\begin{theorem}[Z3.h — Abelian]
\label{Z3.h}
\lean{Cohesive.homOfSec, Cohesive.TwMod.Hom.glue_comm, Cohesive.TwMod.isIso_coimageImageComparison, Cohesive.TwMod.kerIsLimit, Cohesive.TwMod.cokIsColimit, Cohesive.sec_injective_of_mono}
\leanok
For every sheaf of commutative rings R on X, family of opens (U\textsubscript{α}) and family of units ȟ\textsubscript{αβγ} (no cocycle condition needed), the category of ȟ-twisted sheaves of R-modules is abelian. Kernels and cokernels are computed componentwise, with gluing maps induced from those of the source and target, and the component functors M ↦ M\textsubscript{α} preserve them. Supporting facts: a compatible family of R(W)-linear section maps over all W ⊆ V defines a morphism of restricted sheaves on V, and a monomorphism of sheaves of modules is injective on sections.
\par\smallskip{\footnotesize Source: blueprint-tw.md §4.}
\end{theorem}
\begin{proof}
\leanok
\uses{Z3.f, Z3.g}
\textbf{C.2} \texttt{Abelian}: componentwise kernels/cokernels (sections commute with finite limits; \texttt{restrOver} exact), \texttt{Abelian.ofCoimageImageComparisonIsIso} with \texttt{PreservesCoimageImageComparison.iso}; quasi-isomorphisms componentwise.

\textbf{C.2′} (refinement, 3 Nov 2026, P19): sections do not commute with cokernels (sheafification), so \texttt{φ} on a cokernel is not defined section by section. Bridge: a family \texttt{φ α β} over all \texttt{W ≤ U\_α ∩ U\_β}, \texttt{R(W)}-linear and natural, is the same as a morphism of sheaves \texttt{glueHom α β : restrOver inf\_le\_left (M α) ⟶ restrOver inf\_le\_right (M β)} on \texttt{Over (U\_α ∩ U\_β)} (sections of \texttt{restrOver} are the original sections by \texttt{rfl}, P19). Kernels: componentwise, \texttt{φ} restricted (sections commute with kernels). Cokernels: componentwise; \texttt{glueHom} is the map induced on cokernels (\texttt{restrOver} preserves cokernels); \texttt{φ\_self} and the cocycle are equalities of sheaf morphisms on \texttt{Over W}-type over-sites, checked after precomposing with the epimorphism \texttt{restrOver (coker.π)}; then read on sections. Implementation (3 Nov 2026): \texttt{homOfSec} turns a linear natural section family on \texttt{W' ≤ W} into a sheaf morphism on \texttt{Over W}; \texttt{sec\_eq\_of\_epi} cancels \texttt{restrOver p} (cocycle of a quotient, with the scalar \texttt{ȟ} as a section family, so no scalar endomorphism of sheaves is needed); sub-objects use injectivity of monos on sections; \texttt{TwMod.ofMono}/\texttt{ofEpi} package this once for kernels and cokernels; biproducts use joint injectivity of \texttt{fst, snd} on sections.

\textbf{C.3} products, coproducts componentwise.
\paragraph{Blueprint notes.}
C.2–C.3

Files: \texttt{Sheaf/Twisted/SecHom.lean}, \texttt{SubQuot.lean}, \texttt{Preadditive.lean}, \texttt{Kernels.lean}, \texttt{Cokernels.lean}, \texttt{Biprod.lean}, \texttt{Abelian.lean}.
\end{proof}

\begin{theorem}[Z3.i — Existence of injectives]
\label{Z3.i}
\lean{Cohesive.pushCounit, Cohesive.pushHomEquiv, Cohesive.TwistCocycle.unitRes_cyclic, Cohesive.TwMod.mono_of_mono_f, Cohesive.TwMod.comp_preservesMonomorphisms, Cohesive.push_preservesInjectiveObjects}
\leanok
Let R be a sheaf of commutative rings on X and ȟ a twist cocycle of R on a finite cover (U\textsubscript{α}).

\begin{itemize}
\item For each α, Toda's twisted pushforward N ↦ push\textsubscript{α}(N), from sheaves of R|U\textsubscript{α}-modules to ȟ-twisted sheaves (sections over W ⊆ U\textsubscript{β} are N(W ∩ U\textsubscript{α}), gluing by ȟ\textsubscript{αβγ}), is right adjoint to the component functor M ↦ M\textsubscript{α}.
\item The component functors preserve monomorphisms, and a morphism with mono components is mono; hence push\textsubscript{α} preserves injective objects.
\item Consequently the category of ȟ-twisted sheaves has enough injectives.
\item The cocycle satisfies the cyclic symmetry ȟ\textsubscript{βγα} = ȟ\textsubscript{αβγ}.
\end{itemize}
\par\smallskip{\footnotesize Source: blueprint-tw.md §4.}
\end{theorem}
\begin{proof}
\leanok
\uses{Z3.c, Z3.f, Z3.h}
\textbf{C.4} enough injectives: \texttt{push\_α N} (\texttt{β}-component \texttt{W ↦ N(W ⊓ U\_α)}, gluing \texttt{ch\_αβγ·}), \texttt{comp\_α ⊣ push\_α} (uses B.6), \texttt{push\_α} preserves injectives, \texttt{M ↪ ∏\_α push\_α I\_α} (components have enough injectives: \texttt{IsGrothendieckAbelian}).

\textbf{C.4′} (detail, 3 Nov 2026, P20). Fix a twist cocycle \texttt{T} (normalized, Čech cocycle). \textbf{(a)} Cyclic symmetry: the cocycle at \texttt{(α, β, γ, α)} with \texttt{ȟ\_αγα = ȟ\_αβα = 1} gives \texttt{ȟ\_βγα = ȟ\_αβγ} on every \texttt{W} in the three opens. \textbf{(b)} \texttt{j\_α N := overPullback (U\_α ≤ ⊤) N} on \texttt{X}; sections over \texttt{W} are \texttt{N(P\_α W)}, \texttt{P\_α W := pullback (W ≤ ⊤) (U\_α ≤ ⊤)} (a chosen pullback in \texttt{Opens X}, equal to \texttt{W ⊓ U\_α} but not definitionally; \texttt{P\_α W ≤ W}, \texttt{P\_α W ≤ U\_α}, \texttt{W ⊓ U\_α ≤ P\_α W}); restriction is restriction in \texttt{N} and \texttt{r ∈ R(W)} acts by \texttt{r|P\_α W} (P20, all \texttt{rfl}). \textbf{(c)} \texttt{push\_α N}: components \texttt{(push\_α N)\_β := restrOver (U\_β ≤ ⊤) (j\_α N)}, so sections over \texttt{W ≤ U\_β} are \texttt{N(P\_α W)} for every \texttt{β} (definitionally, independent of \texttt{β}: this is why one pushes to \texttt{X} first); gluing \texttt{φ\_βγ x := ȟ\_αβγ|P\_α W · x}. Linear (commutativity of \texttt{R}), natural (restriction of units), \texttt{φ\_ββ = id} (\texttt{ȟ\_αββ = 1}), cocycle \texttt{ȟ\_αγδ ȟ\_αβγ = ȟ\_βγδ ȟ\_αβδ} (Čech cocycle at \texttt{(α, β, γ, δ)}). On morphisms \texttt{push\_α k} has components \texttt{restrOver (j\_α k)} (commutes with scalars). \textbf{(d)} Counit \texttt{ε\_N : (push\_α N)\_α ⟶ N}, restriction along \texttt{W ≤ P\_α W} (\texttt{W ≤ U\_α}); an iso. \textbf{(e)} \texttt{Hom(M, push\_α N) ≃ Hom(M\_α, N)}: \texttt{F ↦ F\_α ≫ ε\_N}; inverse \texttt{g ↦ G}, \texttt{G\_β(x) := g(φ\_βα(x|P\_α W))} for \texttt{x ∈ M\_β(W)}. \texttt{G} commutes with the gluing: \texttt{G\_γ(φ\_βγ x) = g(φ\_γα φ\_βγ (x|P)) = g(ȟ\_βγα φ\_βα(x|P)) = ȟ\_αβγ G\_β(x)} by (a). Left inverse: \texttt{F.comm} at \texttt{(β, α)} and \texttt{ȟ\_αβα = 1} give \texttt{F\_α(φ\_βα y) = F\_β(y)}, so \texttt{G\_β x = ε(F\_β(x)|P)} and \texttt{P\_α(P\_α W) = P\_α W} up to restrictions (functoriality of restriction). Right inverse: \texttt{ε(G\_α x) = g(x)} by \texttt{φ\_αα = id} and naturality of \texttt{g}. Naturality in \texttt{M} is associativity, in \texttt{N} naturality of \texttt{ε}. Hence \texttt{comp\_α ⊣ push\_α} (\texttt{Adjunction.mkOfHomEquiv}). \textbf{(f)} \texttt{comp\_α} preserves monomorphisms (it preserves kernels), so \texttt{push\_α} preserves injectives (\texttt{Functor.preservesInjectiveObjects\_of\_adjunction\_of\_preservesMonomorphisms}). \textbf{(g)} Given \texttt{M}, choose monos \texttt{ι\_α : M\_α ↪ I\_α}, \texttt{I\_α} injective; \texttt{g := Pi.lift (fun α => (homEquiv)⁻¹ ι\_α) : M ⟶ ∏\_α push\_α I\_α} (finite product, injective). \texttt{g} is mono because each component \texttt{g.f β} is: \texttt{(g ≫ π\_β).f β ≫ ε = ι\_β} is mono. (A morphism of twisted sheaves with mono components is mono, by \texttt{hom\_ext}.) So \texttt{EnoughInjectives (TwMod R T.U T.ch)} for finite covers.

\textbf{C.5} restriction preserves injectives (exact left adjoints). (3 Nov 2026: for \texttt{comp\_α} the left adjoint is \texttt{extTw α}, built in Z3.l, so the \texttt{comp\_α} part of C.5 moves to Z3.l; \texttt{restrOver} part done in Z3.f.)
\paragraph{Blueprint notes.}
C.4–C.5

Files: \texttt{Sheaf/Twisted/Pushforward.lean}, \texttt{Injective.lean}.
\end{proof}

\begin{theorem}[Z3.j — Perfectness]
\label{Z3.j}
\lean{Cohesive.glueNatIso, Cohesive.IsPerfectOver, Cohesive.IsStrictlyPerfect.restrOver, Cohesive.TwIsPerfect, Cohesive.IsPerfectOverIn, Cohesive.IsPerfectOverIn.exactAt}
\leanok
Let R be a sheaf of commutative rings on X.

\begin{itemize}
\item A complex K of sheaves of R|V-modules on the over-site of an open V is \emph{perfect} if every x ∈ V has an open neighbourhood W ⊆ V with quasi-isomorphisms K|W ← L → P, P strictly perfect. If moreover all such P can be chosen concentrated in degrees [a, b], then the homology sheaves of K vanish outside [a, b]; restrictions of strictly perfect complexes are strictly perfect.
\item A complex K of ȟ-twisted sheaves is \emph{twisted perfect} if each component complex K\textsubscript{α} is perfect on the over-site of U\textsubscript{α}. D\textsubscript{perf}(X, ȟ) is the class of objects of the derived category of ȟ-twisted sheaves isomorphic to Q(K) for a twisted perfect K.
\item If ȟ\textsubscript{αβα} = 1 for all α, β, the gluing maps assemble into a natural isomorphism M\textsubscript{α}|U\textsubscript{α} ∩ U\textsubscript{β} ≅ M\textsubscript{β}|U\textsubscript{α} ∩ U\textsubscript{β}.
\end{itemize}
Deviation: twisted perfectness is componentwise (the gluing maps play no role), and local roofs of quasi-isomorphisms replace isomorphisms in the derived category.
\par\smallskip{\footnotesize Source: blueprint-tw.md §4.}
\end{theorem}
\begin{proof}
\leanok
\uses{S5, S9, Z3.c, Z3.h}
\textbf{C.6} \texttt{TwIsPerfect}, invariance, twisted S9.d (compactness). Detail of C.6 (3 Nov 2026, Z3.j). Restrictions compose definitionally: \texttt{restrOver h' ∘ restrOver h = restrOver (h'.trans h)} and \texttt{restrOver (le\_refl V) = 𝟭}, also on complexes (\texttt{Iso.refl} checks). \textbf{(a)} \texttt{restrOver} preserves summands of finite free modules (it commutes with \texttt{free}, as \texttt{restr} in S5.a) and strict bounds (additive), hence strict perfectness; it is exact, so preserves quasi-isomorphisms. \textbf{(b)} Over-site S9.c: \texttt{IsPerfectOver} is invariant under quasi-isomorphism (compose roofs; homotopy pullback of roofs through \texttt{D(R|W)}), verbatim as S9.c. \textbf{(c)} \texttt{IsPerfectOverIn K a b}: models in \texttt{[a, b]}. Restriction: a roof on \texttt{W ∋ x} restricts to \texttt{W ⊓ V'} by (a). Locality: if every \texttt{x ∈ V} has \texttt{V' ∋ x} with \texttt{IsPerfectOverIn (K|V') a b}, then \texttt{IsPerfectOverIn K a b} (models of \texttt{K|V'} are models of \texttt{K}, since \texttt{K|V'|W = K|W}). A single roof on \texttt{V} gives \texttt{IsPerfectOverIn K a b} (take \texttt{W = V}). \textbf{(d)} Homology of \texttt{K} with \texttt{IsPerfectOverIn K a b} vanishes outside \texttt{[a, b]}: verbatim S9.d, with locality of zero sheaves on \texttt{Over V} (sections over \texttt{W ≤ V} are determined on a cover of \texttt{W}). \textbf{(e)} \texttt{comp\_α} is exact (C.2), so it preserves quasi-isomorphisms and homology; a twisted sheaf with all components zero is zero. Hence \texttt{TwIsPerfect} is invariant under quasi-isomorphism and \texttt{TwPerf (Q K) ↔ TwIsPerfect K} (as S9.c), and \texttt{K.ExactAt i ↔ ∀ α, K\_α.ExactAt i}. \textbf{(f)} For a normalized \texttt{ch}, the gluing maps give an isomorphism of functors \texttt{comp\_α ⋙ resL α β ≅ comp\_β ⋙ resR α β} (\texttt{glueHom} is natural in \texttt{M} and has inverse built from \texttt{φ\_βα}, \texttt{φ\_symm}); so \texttt{K\_α|W ≅ K\_β|W} for every \texttt{W ≤ U\_α ⊓ U\_β}, as complexes. \textbf{(g)} Twisted S9.d: let \texttt{X} be compact, the cover cover \texttt{X}, \texttt{K} twisted perfect. For each \texttt{x} choose \texttt{β(x)} with \texttt{x ∈ U\_β(x)} and a roof for \texttt{K\_β(x)} on \texttt{W\_x ∋ x}, models in \texttt{[a\_x, b\_x]}; a finite subcover \texttt{W\_\{x\_1\}, …} gives one range \texttt{[a, b]}. For any \texttt{α} and \texttt{y ∈ U\_α}, \texttt{y ∈ W\_\{x\_i\}}; with \texttt{W' = W\_\{x\_i\} ⊓ U\_α}, the roof restricted to \texttt{W'} (c) and transported along \texttt{K\_β|W' ≅ K\_α|W'} (f) gives \texttt{IsPerfectOverIn (K\_α|W') a b}; locality (c) gives \texttt{IsPerfectOverIn K\_α a b}. So \texttt{TwIsPerfectIn K a b}, and by (d), (e) the homology of \texttt{K} vanishes outside \texttt{[a, b]}. \textbf{(h)} Sanity: \texttt{IsPerfectOver (K|⊤) ↔ IsPerfect K} (\texttt{restr\_restrOver}).
\paragraph{Blueprint notes.}
C.6

Files: \texttt{Sheaf/Twisted/PerfectDef.lean}, \texttt{PerfectOver.lean}, \texttt{Perfect.lean}.
\end{proof}

\begin{theorem}[Z3.k — Gluing]
\label{Z3.k}
\lean{Cohesive.glueEval, Cohesive.glueCoeval, Cohesive.ofSheafGlueIso, Cohesive.toGlueOfSheaf, Cohesive.rescale, Cohesive.isSheaf_gluePre}
\leanok
Let R be a sheaf of commutative rings on X and (U\textsubscript{α}) a family of opens.

\begin{itemize}
\item For a twisted sheaf M with trivial cocycle, the compatible families (s\textsubscript{α}) ∈ ∏\textsubscript{α} M\textsubscript{α}(W ∩ U\textsubscript{α}) form a sheaf glue(M) of R-modules on X, and M is isomorphic to the twisted sheaf of restrictions of glue(M) (no cover hypothesis).
\item If ⋃\textsubscript{α} U\textsubscript{α} = X, gluing is an equivalence between twisted sheaves with trivial cocycle and sheaves of R-modules on X, with inverse N ↦ (N|U\textsubscript{α}, identity gluing).
\item (Rescaling.) For units g\textsubscript{αβ} ∈ R(U\textsubscript{α} ∩ U\textsubscript{β})\textasciicircum{}× with g\textsubscript{αα} = 1 and families ȟ, ȟ′ with ȟ′\textsubscript{αβγ} g\textsubscript{αγ} = g\textsubscript{βγ} g\textsubscript{αβ} ȟ\textsubscript{αβγ}, replacing φ\textsubscript{αβ} by g\textsubscript{αβ} φ\textsubscript{αβ} is an equivalence between ȟ-twisted and ȟ′-twisted sheaves, the identity on components and morphisms.
\end{itemize}
\par\smallskip{\footnotesize Source: blueprint-tw.md §4.}
\end{theorem}
\begin{proof}
\leanok
\uses{S8, Z3.i}
\textbf{C.7} S12 gluing for \texttt{ch = 1}: \texttt{glue(N)(W) = \{(s\_α) : ψ\_αβ s\_α = s\_β\}}, \texttt{glue|U\_α ≅ N\_α} (cocycle only), \texttt{glueEquiv : TwMod R U 1 ≌ Mod R} (sheaf condition, \texttt{TopCat.Sheaf.existsUnique\_gluing}); \texttt{rescale g : TwMod R U ch ≌ TwMod R U (ch·ďg)}; tensor with summands of finite free \texttt{𝒜⁰}-modules commutes with \texttt{glue}. Detail of C.7 (3 Nov 2026, Z3.k; prototype P21). \textbf{(a)} Sections of a sheaf on \texttt{Over U} glue uniquely along open covers of \texttt{W ≤ U} (sheaf condition for the sieve generated by the cover, pullbacks are \texttt{⊓}). \textbf{(b)} For \texttt{M : TwMod R U 1}, \texttt{glue M (W) := \{(s\_α) ∈ ∏\_α M\_α(W ⊓ U\_α) : φ\_αβ(s\_α|W') = s\_β|W' for all W' ≤ W ⊓ U\_α ⊓ U\_β\}}, an \texttt{R(W)}-module via \texttt{R(W) → R(W ⊓ U\_α)}; restriction componentwise. (The condition is quantified over all \texttt{W'}, so no intersections are rewritten; by \texttt{φ\_res} it is the usual one.) Sheaf: given compatible families on a cover \texttt{W\_i} of \texttt{W}, for each \texttt{α} the components \texttt{s\textasciicircum{}i\_α} on \texttt{W\_i ⊓ U\_α} agree on overlaps, so glue (a) to \texttt{s\_α} on \texttt{W ⊓ U\_α}; the compatibility of \texttt{(s\_α)} holds locally on \texttt{W' ∩ W\_i}, hence by separatedness. \textbf{(c)} \texttt{ofSheaf (glue M) ≅ M}, no cover needed: on \texttt{W ≤ U\_γ}, \texttt{(s\_α) ↦ s\_γ|W} (\texttt{W ≤ W ⊓ U\_γ}), inverse \texttt{t ↦ (φ\_γα(t|W ⊓ U\_α))\_α}, compatible by the cocycle with \texttt{ȟ = 1} and \texttt{φ\_res}; inverse laws by \texttt{φ\_γγ = id} and compatibility of \texttt{(s\_α)} at \texttt{W' = W ⊓ U\_α}; it commutes with the gluing maps (cocycle again); natural in \texttt{M}. \textbf{(d)} For \texttt{⨆ U\_α = ⊤}, \texttt{glue (ofSheaf N) ≅ N}: \texttt{t ↦ (t|W ⊓ U\_α)\_α}, bijective by the sheaf condition of \texttt{N} for the cover \texttt{W ⊓ U\_α} of \texttt{W}; natural in \texttt{N}. \textbf{(e)} \texttt{glueEquiv : TwMod R U 1 ≌ Mod R} from (c), (d) (\texttt{Equivalence.mk}). \textbf{(f)} \texttt{rescale g}: for units \texttt{g\_αβ ∈ R(U\_αβ)ˣ} with \texttt{g\_αα = 1} and \texttt{ȟ'\_αβγ g\_αγ = g\_βγ g\_αβ ȟ\_αβγ} on every \texttt{W ≤ U\_αβγ}, \texttt{φ' := g\_αβ φ} defines \texttt{TwMod R U ȟ ⥤ TwMod R U ȟ'} (identity on components and morphisms); \texttt{g⁻¹} gives the inverse, unit and counit have identity components. With \texttt{g = 1} this is transport along \texttt{ȟ = ȟ'}. \textbf{(g)} "Tensor with summands of finite free \texttt{𝒜⁰}-modules commutes with \texttt{glue}" needs the twisted base change (Z3.l \texttt{BaseChange.lean}): moved to Z3.l.

\textbf{C.7′} (Z3.l′ detail, 5 Nov 2026, for review at G5; prototype P24). Fix a map of sheaves of commutative rings \texttt{ψ : R ⟶ S} on \texttt{X}; \texttt{S ⊗ −} is S8.a's \texttt{sheafBC} on an over-site, \texttt{η} the unit of \texttt{sheafBC ⊣ restrictScalars} (\texttt{η(m) = 1 ⊗ m}, as C7.b's \texttt{bcη}). \textbf{(a)} \emph{Base change commutes with restriction between over-sites}: for \texttt{V ≤ U} and \texttt{M} on \texttt{Over U}, \texttt{(S|U ⊗ M)|V ≅ S|V ⊗ M|V} (\texttt{restrOver}), natural in \texttt{M}, and on sections over \texttt{W ≤ V} it sends \texttt{η(m)} to \texttt{η(m)}. Proof: S8.a's \texttt{sheafBCOverIso} with \texttt{Over.map (V ≤ U)} in place of \texttt{Over.forget} (sectionwise the same presheaf; Mathlib's compatibility of pushforward along a continuous functor with sheafification, which applies to \texttt{Over.map} by P24). \textbf{(b)} \emph{Uniqueness}: two \texttt{S(W′)}-linear families on sections of \texttt{S|U ⊗ M} over \texttt{W′ ≤ W ≤ U}, natural in \texttt{W′}, with values in a sheaf of \texttt{S}-modules, are equal as soon as they agree on all \texttt{η(m)}, \texttt{m ∈ M(W′)}. Proof: as morphisms on \texttt{Over W} (\texttt{homOfSec}), transported by (a), they correspond under \texttt{sheafBC ⊣ restrictScalars} to their composites with \texttt{η}. \textbf{(c)} *\texttt{twBC ψ : TwMod R U ȟ ⥤ TwMod S U ψ(ȟ)}*: components \texttt{S|U\_α ⊗ M\_α}; gluing over \texttt{W ≤ U\_α ∩ U\_β} given by the sections of \texttt{(a)⁻¹ ∘ (S ⊗ g\_αβ) ∘ (a)}, where \texttt{g\_αβ} is \texttt{φ\_αβ} as a morphism on \texttt{Over (U\_α ∩ U\_β)} (\texttt{homOfSec}, C.2′). Linearity and naturality hold because this is a morphism of sheaves. Key formula: \texttt{φ′\_αβ(η m) = η(φ\_αβ m)} (from (a) and naturality of \texttt{η}). \texttt{φ′\_αα = id} and the cocycle \texttt{φ′\_βγ φ′\_αβ = ψ(ȟ\_αβγ)·φ′\_αγ} hold by (b): on \texttt{η(m)} both sides are \texttt{η(φ\_βγ φ\_αβ m) = η(ȟ\_αβγ·φ\_αγ m) = ψ(ȟ\_αβγ)·η(φ\_αγ m)} (\texttt{η} is \texttt{R}-linear into the restriction of scalars). On morphisms \texttt{S ⊗ f\_α} componentwise, which commutes with the gluing by (b). Functorial and additive (\texttt{sheafBC} is), and \texttt{comp\_α ∘ twBC = sheafBC ∘ comp\_α} by \texttt{rfl}. \textbf{(d)} \emph{Scope change}: the statement "tensor with summands of finite free \texttt{𝒜⁰}-modules commutes with \texttt{glue}" (C.7 (g)) follows from C7.a (\texttt{T} is a sheaf for every \texttt{F}) and C.7 (c). It is proved where it is used (Z3.t, G.5: assembling the global curved connection from local ones), not as a separate file of Z3.l′.
\paragraph{Blueprint notes.}
C.7 (a)–(f)

Files: \texttt{Sheaf/Twisted/GlueSec.lean}, \texttt{Glue.lean}, \texttt{GlueEquiv.lean}, \texttt{Rescale.lean}.
\end{proof}

\begin{theorem}[Z3.l — Generators]
\label{Z3.l}
\lean{Cohesive.extXRestrIso, Cohesive.extTwLiftSec, Cohesive.extTwLift, Cohesive.extTwUnit, Cohesive.extX_locally_zero, Cohesive.isFlat_extX}
\leanok
Let R be a sheaf of commutative rings on X and ȟ a twist cocycle of R on a finite cover (U\textsubscript{α}).

\begin{itemize}
\item For a sheaf F of R|V-modules, its extension by zero j\_!F is a sheaf of R-modules on X with (j\_!F)|V ≅ F, whose sections vanish near every point outside V; if F is flat, so is j\_!F.
\item The twisted extension by zero N ↦ ext\textsubscript{α}(N) (components (j\_!N)|U\textsubscript{β}, gluing by ȟ\textsubscript{αβγ}) is left adjoint to the component functor M ↦ M\textsubscript{α}.
\item Every twisted sheaf is a quotient of one with flat components, and every complex of twisted sheaves strictly zero above degree b receives a quasi-isomorphism from such a complex with componentwise flat terms.
\end{itemize}
\par\smallskip{\footnotesize Source: blueprint-tw.md §4.}
\end{theorem}
\begin{proof}
\leanok
\uses{S8, Z3.i, Z3.k}
Step C.5: see Z3.i.

\textbf{C.8} generators \texttt{extTw α (extZero\_V R)} with \texttt{Hom(G, M) ≅ M\_α(V)}; bounded-above resolutions by finite sums of them (\texttt{exists\_quasiIso\_of\_enough}), components flat. (v1.20, blueprint §11, approved 3 Nov 2026: explicit extension by zero \texttt{extX V F} on \texttt{X}, \texttt{extTw α N} with components \texttt{restr (U β) (extX (U α) N)} and gluing \texttt{ȟ\_αβγ ·}, \texttt{extTw α ⊣ comp\_α}, flatness by the S8.b argument, resolutions by \texttt{⊕\_α extTw α F\_α}; base change and C.7 (g) split off as Z3.l′.)
\paragraph{Blueprint notes.}
C.5 (\texttt{comp\_α}), C.8

Files: \texttt{Sheaf/Twisted/ExtZero.lean}, \texttt{ExtTw.lean}, \texttt{ExtTwAdj.lean}, \texttt{FlatResolution.lean}; \texttt{isFlat\_of\_locally\_zero} in \texttt{Sheaf/Flat.lean}.
\end{proof}

\begin{theorem}[Z3.l′ — Base change]
\label{Z3.lprime}
\lean{Cohesive.sheafBCRestrOverIso, Cohesive.sec_eq_of_bcUnit, Cohesive.bcφ_bcUnit}
\leanok
Let ψ : R → S be a morphism of sheaves of commutative rings on X.

\begin{itemize}
\item For V ⊆ U and a sheaf M of R|U-modules, base change commutes with restriction: (S|U ⊗ M)|V ≅ S|V ⊗ M|V.
\item Section maps out of S ⊗ M that are linear and natural are determined by their values on the elements 1 ⊗ m.
\item Componentwise base change is a functor from ȟ-twisted R-modules to ψ(ȟ)-twisted S-modules, whose gluing maps satisfy φ′\textsubscript{αβ}(1 ⊗ m) = 1 ⊗ φ\textsubscript{αβ}(m).
\end{itemize}
\par\smallskip{\footnotesize Source: blueprint-tw.md §4.}
\end{theorem}
\begin{proof}
\leanok
\uses{S8, Z3.h}
Step C.7: see Z3.k.

Step C.7′: see Z3.k.
\paragraph{Blueprint notes.}
C.7 (g), componentwise \texttt{sheafBC}; before Z3.u

Files: \texttt{Sheaf/Twisted/BaseChange.lean}, \texttt{GlueTensor.lean}.
\end{proof}

\begin{theorem}[Z3.m — C1 over \texttt{DgaSheafData}]
\label{Z3.m}
\lean{Cohesive.DgaSheafData.A0, Cohesive.DgaSheafData.𝒜X, Cohesive.DgaSheafData.r, Cohesive.DgaSheafData.R, Cohesive.DgaSheafData.TotSec, Cohesive.DgaSheafData.deg0}
\leanok
\uses{C1}
Let S be sheaf data on X (possibly curved), U₀ an open and M a cohesive module over 𝒜(U₀). The constructions of C1 are defined over S word for word: the sheaf of rings 𝒜\textsuperscript{0}, the global algebra 𝒜(X) (with the curvature of S), the restrictions r\textsubscript{V}, the ring sheaf O, the total module Tot(M|V) for V ⊆ U₀, and its elements of total degree n, an O(V)-module α(M)\textsuperscript{n}(V). These form a sheaf α(M)\textsuperscript{n} of O-modules with the operator 𝕏 : α(M)\textsuperscript{n} → α(M)\textsuperscript{n+1}, and α acts on closed degree-0 maps and homotopies. For pre-resolving data these agree definitionally with C1.

Deviation: no complex α(M) and no functor are built here, since 𝕏\textsuperscript{2} = c· need not vanish.
\par\smallskip{\footnotesize Source: blueprint-tw.md §4.}
\end{theorem}
\begin{proof}
\leanok
\uses{Z3.a}
\textbf{D.6} local identification \texttt{(J\_T M)\_α|W' ≅ sheafOfOver\_\{W'\}(gauge\_\{h0\_α\}(M|U\_α)|W')} (C1.c, A.2/A.4).
\paragraph{Blueprint notes.}
copies + rfl bridges (done 3 Nov 2026: AlphaDef, C1.a, C1.b section and degree maps; no complex/\texttt{J}, they need \texttt{c = 0}; \texttt{AlphaOver} not copied, D.6 applies it to the flat gauged module)

Files: \texttt{Comparison/SheafDataAlphaDef.lean}, \texttt{SheafDataAlpha.lean}, \texttt{SheafDataMap.lean}.
\end{proof}

\begin{theorem}[Z3.n — \texttt{J\_T}]
\label{Z3.n}
\lean{Cohesive.DgaSheafData.smul0L, Cohesive.DgaSheafData.gaugeDV, Cohesive.DgaSheafData.gaugeDV_gaugeDV, Cohesive.DgaSheafData.gaugeDV_smul0L, Cohesive.DgaSheafData.bcHomL_gauge_homotopy, Cohesive.TwistDatum.twTerm}
\leanok
Let S be sheaf data on X, h ∈ 𝒜(X)\textsuperscript{2} closed, T a twist datum for h over S, and let M be a cohesive module over the curved algebra (𝒜(X), h).

\begin{itemize}
\item For any sheaf data, β ∈ 𝒜\textsuperscript{1}(V) and M over 𝒜(U₀), V ⊆ U₀, the gauged operator 𝕏 − β· on α(M)(V) satisfies (𝕏 − β)\textsuperscript{2} = 0 if dβ is the curvature over V (h|V for the curved data); if du = (β′ − β)u for u ∈ 𝒜\textsuperscript{0}(V), then (𝕏 − β′)(u·x) = u·(𝕏 − β)x; and homotopies of maps give homotopies for the gauged operators.
\item J\textsubscript{T}(M) is the complex of ȟ-twisted O-modules with components J\textsubscript{T}(M)\textsuperscript{n\_α} = α(M)\textsuperscript{n}|U\textsubscript{α}, gluing maps u\textsubscript{αβ}· and differential 𝕏 − h\textsuperscript{0}\textsubscript{α} on U\textsubscript{α}.
\item M ↦ J\textsubscript{T}(M) defines functors J\textsubscript{T} : Ho(P\textsubscript{(𝒜(X), h)}) ⥤ K(TwMod(O, U, ȟ)) and, for any derived category structure, Ho(P\textsubscript{(𝒜(X), h)}) ⥤ D(TwMod(O, U, ȟ)).
\end{itemize}
Deviation: the gluing is by the units u\textsubscript{αβ} instead of Antweiler's exp(h\textsuperscript{1}\textsubscript{αβ}).
\par\smallskip{\footnotesize Source: blueprint-tw.md §4.}
\end{theorem}
\begin{proof}
\leanok
\uses{C1, Z3.i, Z3.m}
\textbf{D.1} terms (curved C1.a copy).

\textbf{D.2} \texttt{d\_α = 𝕏 − h0\_α}, \texttt{d\_α² = h − dh0\_α = 0}.

\textbf{D.3} gluing \texttt{u\_αβ·}, chain map by A.5, cocycle by (4).

\textbf{D.4} morphisms and homotopies (C1, \texttt{k} commutes with \texttt{u}).
\paragraph{Blueprint notes.}
D.1–D.4

Files: \texttt{Comparison/TwistedFunctorDef.lean}, \texttt{TwistedFunctor.lean}.
\end{proof}

\begin{theorem}[Z3.o — Local identification]
\label{Z3.o}
\lean{Cohesive.twLocalIso, Cohesive.curvedResM, Cohesive.gaugedRes, Cohesive.twCmp, Cohesive.twCmp_X}
\leanok
Let P be pre-resolving data on X, h ∈ 𝒜(X)\textsuperscript{2} closed, T a twist datum for h over P and M a cohesive module over (𝒜(X), h). For each α, the restriction M|U\textsubscript{α} is a cohesive module over (𝒜(U\textsubscript{α}), h|U\textsubscript{α}), and its gauge gauge\_\{h\textsuperscript{0}\textsubscript{α}\}(M|U\textsubscript{α}) (Z3.b, h|U\textsubscript{α} = dh\textsuperscript{0}\textsubscript{α}) is a cohesive module over the flat algebra 𝒜(U\textsubscript{α}). The comparison Tot(M|V) → Tot(gauge\_\{h\textsuperscript{0}\textsubscript{α}\}(M|U\textsubscript{α})|V), V ⊆ U\textsubscript{α}, intertwines 𝕏 − h\textsuperscript{0}\textsubscript{α} with the operator of the gauged module, and the component complex J\textsubscript{T}(M)\textsubscript{α} on the over-site of U\textsubscript{α} is isomorphic to α\textsubscript{U\_α}(gauge\_\{h\textsuperscript{0}\textsubscript{α}\}(M|U\textsubscript{α})) (C1).
\par\smallskip{\footnotesize Source: blueprint-tw.md §4.}
\end{theorem}
\begin{proof}
\leanok
\uses{Z3.b, Z3.m, Z3.n}
Step D.6: see Z3.m.
\paragraph{Blueprint notes.}
D.6 (generic \texttt{P : PreResolvingData}; \texttt{gaugedRes}, \texttt{twCmp})

Files: \texttt{Comparison/TwistedLocal.lean}; \texttt{Scratch/P23\_TwLocal.lean}.
\end{proof}

\begin{theorem}[Z3.p — Image of J]
\label{Z3.p}
\lean{Cohesive.PreResolvingData.restrOverSheafOfIso, Cohesive.isPerfectOver_sheafOfOver, Cohesive.twIsPerfect_twistedSheafOf}
\leanok
\uses{C1, H2}
Let X be a compact Hausdorff complex manifold modelled on a finite-dimensional complex normed space E.

\begin{itemize}
\item For pre-resolving data, opens U′ ⊆ U and a cohesive module N over 𝒜(U), α\textsubscript{U}(N)|U′ ≅ α\textsubscript{U′}(N|U′) as complexes on the over-site of U′.
\item For every open U and every cohesive module N over 𝒜\textsubscript{X}(U), α\textsubscript{U}(N) is perfect on U.
\item For a closed h ∈ 𝒜\textsuperscript{0,2}(X), a twist datum T for h and a cohesive module M over (𝒜\textsubscript{X}(X), ∂̄, h), the complex J\textsubscript{T}(M) is twisted perfect.
\end{itemize}
\par\smallskip{\footnotesize Source: blueprint-tw.md §4.}
\end{theorem}
\begin{proof}
\leanok
\uses{C2, Z3.e, Z3.j, Z3.o, Z4.g}
\textbf{E.1} LL\textsuperscript{h} = H2 on opens (Z3.e) for \texttt{gauge\_\{h0\_α\}(M|U\_α)}.

\textbf{E.2} C2's over-site core (base generalized \texttt{⊤ → U₀}, V3) + D.6.
\paragraph{Blueprint notes.}
E.1–E.2 (H2 over opens for any \texttt{N} over \texttt{𝒜(U)}, so no \texttt{curvedLL} bridge is needed)

Files: \texttt{Comparison/AlphaOverRestr.lean} (C1.c between opens), \texttt{Comparison/TwistedPerfect.lean}.
\end{proof}

\begin{theorem}[Z3.q — C3 curved]
\label{Z3.q}
\lean{Cohesive.DgaSheafData.ℋ}
\leanok
Let S be sheaf data on X (possibly curved, curvature c) and M₁, M₂ cohesive modules over 𝒜(X). The \emph{Hom complex of sheaves} ℋ(M₁, M₂) is the complex of sheaves of O-modules whose degree-n sections over V are the degree-n maps χ : Tot(M₁) → Tot(M₂|V) with χ(a·x) = (−1)\textsuperscript{n|a|} a|V·χ(x), with differential δχ = 𝕏₂ ∘ χ − (−1)\textsuperscript{n} χ ∘ 𝕏₁; δ\textsuperscript{2} = 0 because 𝕏₁\textsuperscript{2} and 𝕏₂\textsuperscript{2} are both multiplication by c. Its global sections are Hom\textsuperscript{•\_}\{P\textsubscript{A}\}(M₁, M₂) (a bijection compatible with δ), and if X is compact Hausdorff its terms are soft.
\par\smallskip{\footnotesize Source: blueprint-tw.md §4.}
\end{theorem}
\begin{proof}
\leanok
\textbf{F.2} curved ℋ\textsuperscript{h} (C3 over \texttt{DgaSheafData}); \texttt{𝓗om\textasciicircum{}tw(K, I)} = glued over-site Hom sheaves (conjugation, \texttt{ch} cancels); sections = \texttt{HomComplex} cochains of \texttt{TwMod}; \texttt{Ψ\textasciicircum{}tw = r\_α ∘ sheafOfMapV}.
\paragraph{Blueprint notes.}
F.2 (ℋ\textsuperscript{h})

Files: \texttt{Comparison/CurvedHomSheaf*.lean}, \texttt{CurvedHomGauge.lean}.
\end{proof}

\begin{theorem}[Z3.r — \texttt{𝓗om\textasciicircum{}tw}]
\label{Z3.r}
\lean{Cohesive.secHomOf, Cohesive.twHomComplex, Cohesive.twHomD_restr, Cohesive.hom_ext_of_restr, Cohesive.twHomRestrIso, Cohesive.twHomRestrOverIso}
\leanok
Let R be a sheaf of commutative rings on X, (U\textsubscript{α}) a cover of X, ȟ a family of units with ȟ\textsubscript{αβα} = 1, and K, L complexes of ȟ-twisted sheaves.

\begin{itemize}
\item For twisted sheaves K, L, the Hom sheaves Hom(K\textsubscript{α}, L\textsubscript{α}) with gluing s ↦ φ\textsuperscript{L\_αβ} ∘ s ∘ φ\textsuperscript{K\_βα} form a twisted sheaf with trivial cocycle; its gluing has global sections Hom(K, L).
\item 𝓗om\textsuperscript{tw}(K, L) is the complex of sheaves of R-modules on X with 𝓗om\textsuperscript{tw}(K, L)\textsuperscript{n} = ∏\textsubscript{p+n=q} glue(Hom(K\textsuperscript{p}, L\textsuperscript{q})), Mathlib's sign for the differential.
\item For V ⊆ U\textsubscript{α}, 𝓗om\textsuperscript{tw}(K, L)|V ≅ 𝓗om(K\textsubscript{α}|V, L\textsubscript{α}|V) as complexes on the over-site of V, and morphisms of sheaves on X agreeing on every U\textsubscript{α} agree.
\end{itemize}
\par\smallskip{\footnotesize Source: blueprint-tw.md §4.}
\end{theorem}
\begin{proof}
\leanok
\uses{S7, Z3.k}
Step F.2: see Z3.q.
\paragraph{Blueprint notes.}
F.2

Files: \texttt{Sheaf/Twisted/HomSheaf.lean}, \texttt{HomSections.lean}.
\end{proof}

\begin{theorem}[Z3.s — Full faithfulness]
\label{Z3.s}
\lean{Cohesive.bijective_twistedJ_map, Cohesive.full_twistedJ_dolbeault, Cohesive.faithful_twistedJ_dolbeault, Cohesive.TwistDatum.Ψtw, Cohesive.TwistDatum.Ψtw_top, Cohesive.LLOpen, Cohesive.jKtw, Cohesive.evLtw, Cohesive.ιtw}
\leanok
Let D be resolving data (C0) on a compact Hausdorff space X satisfying the local gauge lemma over every open (LL\textsubscript{open}: each cohesive module over 𝒜(U₀) is near each x ∈ U₀ isomorphic to 𝒜(U′) ⊗\textsubscript{O(U′)} W, W bounded finitely generated projective). Let h ∈ 𝒜(X)\textsuperscript{2} be closed and T a twist datum for h. Then:

\begin{itemize}
\item there is a morphism of complexes Ψ\textsuperscript{tw} : ℋ\textsuperscript{h}(M₁, M₂) → 𝓗om\textsuperscript{tw}(J\textsubscript{T} M₁, J\textsubscript{T} M₂), which on global sections sends a closed degree-0 map φ to J\textsubscript{T}(φ), and becomes a quasi-isomorphism after composing with a bounded below injective resolution of J\textsubscript{T}(M₂);
\item J\textsubscript{T} : Ho(P\textsubscript{(𝒜(X), h)}) ⥤ D(TwMod(O, U, ȟ)) is bijective on morphisms.
\end{itemize}
For the Dolbeault data of a compact Hausdorff complex manifold LL\textsubscript{open} holds, so J\textsubscript{T} is full and faithful for every closed h ∈ 𝒜\textsuperscript{0,2}(X) and every twist datum T.
\par\smallskip{\footnotesize Source: blueprint-tw.md §4.}
\end{theorem}
\begin{proof}
\leanok
\uses{A10, C2, C3, C4, S3.b, S4, S5, S6, S7, Z3.e, Z3.p, Z3.r}
\textbf{F.1} injective resolution \texttt{r : J\_T M₂ → I} (C.4).

Step F.2: see Z3.q.

\textbf{F.3} locally \texttt{Ψ\textasciicircum{}tw} is C4's \texttt{Ψ\_\{U'\}} for the gauged flat modules (A.3, C.7, D.6), quasi-iso by C4 step 3 + LL\textsuperscript{h}; globalize (S6).

\textbf{F.4} softness (S7.a′, S3.b).

\textbf{F.5} \texttt{H⁰} dictionary (\texttt{Qh\_map\_bijective\_of\_isKInjective}, \texttt{isKInjective\_of\_injective}). \textbf{F′} (detail, 4 Nov 2026, approved at G4 5 Nov 2026; refines F.1–F.5, no decision changed). Fix sheaf data, \texttt{h}, a twist datum \texttt{T} (normalized \texttt{ȟ}), \texttt{S\_h} the curved data, \texttt{M₁, M₂} cohesive over \texttt{(𝒜(X), h)}. \textbf{(a)} \emph{Curved C3} (Z3.q): C3's section data over \texttt{S : DgaSheafData} word for word (\texttt{HomSec}, \texttt{homRes}, \texttt{homδV}, \texttt{ℋ}, \texttt{homTop}). Only \texttt{δ² = 0} changes: \texttt{δ²χ = 𝕏₂²χ − χ𝕏₁² = c·χ − χ(c·) = 0} (\texttt{c} of degree 2, so \texttt{χ(c·y) = c·χ(y)}). \texttt{ℋ\textasciicircum{}h(M₁, M₂)(X) ≅ Hom\textasciicircum{}•(M₁, M₂)} as in C3. Softness as C3.a(ii). \textbf{(b)} \emph{Hom sections on over-sites.} For \texttt{P, Q} on \texttt{Over U} and \texttt{W ≤ U}, morphisms \texttt{P.over (W → U) ⟶ Q.over (W → U)} (sections of \texttt{homSheaf P Q} over \texttt{W}) are the same as \texttt{R(W′)}-linear natural section families on \texttt{W′ ≤ W} (\texttt{homOfSec₀}-type bridge, P19 rfl sections). \textbf{(c)} *\texttt{homTw K L}* (Z3.r), for \texttt{K, L : TwMod R U ȟ}: the untwisted (\texttt{ȟ = 1}) twisted sheaf with components \texttt{homSheaf (K\_α) (L\_α)} on \texttt{Over U\_α} and gluing by conjugation on sections over \texttt{W ≤ U\_α ∩ U\_β}: \texttt{s ↦ φ\textasciicircum{}L\_αβ ∘ s ∘ φ\textasciicircum{}K\_βα} (by (b), on \texttt{W′ ≤ W}). Linear (target action commutes with \texttt{φ\textasciicircum{}L}), natural, \texttt{φ\_self} from \texttt{φ\textasciicircum{}K\_αα = φ\textasciicircum{}L\_αα = id}; cocycle: \texttt{φ\textasciicircum{}K\_βα φ\textasciicircum{}K\_γβ = ȟ\_αβγ⁻¹ φ\textasciicircum{}K\_γα} (inverse of the cocycle, \texttt{φ\_symm}) and \texttt{φ\textasciicircum{}L\_βγ φ\textasciicircum{}L\_αβ = ȟ\_αβγ φ\textasciicircum{}L\_αγ}, so the scalars cancel (\texttt{s} is \texttt{R}-linear). Functorial: \texttt{homTwPost g}, \texttt{homTwPre f} (componentwise \texttt{homPost}, \texttt{homPre}). \textbf{(d)} *\texttt{𝓗om\textasciicircum{}tw(K, L)}* for complexes of \texttt{TwMod R U ȟ}: degree \texttt{n} is \texttt{∏\_\{p + n = q\} glue(homTw(Kᵖ, L\textasciicircum{}q))} in \texttt{Mod R} (\texttt{glue}, C.7 (b); \texttt{⨆ U = ⊤}), differential with Mathlib's sign as S7 (\texttt{homComplexD}), \texttt{δ² = 0} formally from functoriality of \texttt{homTwPost}/\texttt{homTwPre}; functorial in both variables (\texttt{twHomMap₂}, \texttt{twHomMap₁}). \textbf{(e)} \emph{Global sections}: \texttt{Γ(𝓗om\textasciicircum{}tw(K, L))ⁿ ≅ Cochain\_\{TwMod\}(K, L, n)} compatibly with \texttt{δ} and with \texttt{twHomMap₂} (post-composition): a global section of \texttt{glue(homTw(Kᵖ, L\textasciicircum{}q))} is a family of morphisms \texttt{Kᵖ\_α → L\textasciicircum{}q\_α} compatible with conjugation, i.e. a morphism of twisted sheaves. \textbf{(f)} \emph{Restriction}: for \texttt{V ≤ U\_α}, \texttt{𝓗om\textasciicircum{}tw(K, L)|V ≅ 𝓗om\_\{Over V\}(K\_α|V, L\_α|V)} as complexes (\texttt{glue(N)|U\_α ≅ N\_α}, C.7 (c); restriction preserves products; \texttt{restrOver} commutes with \texttt{homSheaf}), natural in \texttt{L} (and in \texttt{K}). \textbf{(g)} \emph{Twisted S5.c}: for \texttt{I} injective in \texttt{TwMod}, every restriction \texttt{(I\_α)|V} is injective (\texttt{I} is a summand of \texttt{∏\_β push\_β J\_β}, C.4 (g); the components of \texttt{push\_β J} are restrictions of the injective \texttt{j\_β J} on \texttt{X}, \texttt{overPullback} preserves injectives; \texttt{restr\_restrOver}; S5.c). \textbf{(h)} \emph{Softness, bounds}: \texttt{𝒜⁰(X)} acts on each \texttt{J\_T(M₁)\textasciicircum{}p} by morphisms of twisted sheaves (multiplication commutes with \texttt{u\_αβ·}), hence on the terms of \texttt{𝓗om\textasciicircum{}tw(J\_T M₁, I)} by precomposition, unitally and vanishing over \texttt{V} where \texttt{a|V = 0} (S3.b(i)); bounded below for \texttt{I} bounded below (\texttt{J\_T M₁} is bounded above). \textbf{(i)} *\texttt{Ψ\textasciicircum{}tw : ℋ\textasciicircum{}h(M₁, M₂) ⟶ 𝓗om\textasciicircum{}tw(J\_T M₁, J\_T M₂)}*: \texttt{χ ∈ ℋ\textasciicircum{}\{h,n\}(V)} goes, in component \texttt{α}, over \texttt{W ≤ V ⊓ U\_α}, degree \texttt{p}, to C4's \texttt{extAlong (χ|W)}. It commutes with the gluing (\texttt{u\_αβ} has degree \texttt{0}) and with \texttt{δ} (\texttt{h⁰\_α} has degree 1: \texttt{χ(h⁰\_α y) = (−1)ⁿ h⁰\_α χ(y)}, so the \texttt{h⁰\_α}-terms cancel). \textbf{(j)} \emph{Local} (F.3): \texttt{x ∈ U\_α}; LL over opens (Z3.e) for \texttt{gM₁ := gauge\_\{h⁰\_α\}(M₁|U\_α)} gives \texttt{U′ ∋ x}, \texttt{U′ ≤ U\_α}, \texttt{W} and a homotopy equivalence \texttt{φ : 𝒜(U′) ⊗ W → gM₁|U′}; with C2's \texttt{L, ε, u} (C4 step 3), \texttt{ι := u ≫ α\_\{U′\}(φ) ≫ (α\_\{U\_α\}(gM₁)|U′ ≅ α\_\{U′\}(gM₁|U′))⁻¹ ≫ (twLocalIso|U′)⁻¹ : L → (J\_T M₁)\_α|U′}. \texttt{Φ\textasciicircum{}tw := Ψ\textasciicircum{}tw|U′ ≫ (f) ≫ ι\textasciicircum{}*} is a sectionwise quasi-isomorphism: over \texttt{V ≤ U′} it is \texttt{ℋ\textasciicircum{}h(V) ≅ Hom\_along(gM₁|U′, gM₂|V)} (gauge is the identity on \texttt{Hom}, A.3; base change, as C4's \texttt{jK}) \texttt{→φ\textasciicircum{}*→ Hom\_along(𝒜(U′) ⊗ W, gM₂|V) ≅ Hom\_\{O(U′)\}(W, Tot(gM₂|V))} (A10.a) \texttt{≅ Cochain(L|V, α\_\{U\_α\}(gM₂)|V)} (S4.c, C4's \texttt{evL} with base \texttt{U\_α}) \texttt{≅ Cochain(L|V, (J\_T M₂)\_α|V)} (D.6). Then \texttt{r\_*} (S7.b′; \texttt{r\_α|U′} a quasi-iso, \texttt{comp\_α}, \texttt{restrOver} exact) and \texttt{ι\textasciicircum{}*} (S7.c′, (g)): \texttt{(Ψ\textasciicircum{}tw ≫ r\_*)|U′} is a quasi-isomorphism; S6. \textbf{(k)} \emph{Global} (F.4–F.5): \texttt{Γ(Ψ\textasciicircum{}tw ≫ r\_*)} is a quasi-isomorphism (S3.b(iv), (h)); \texttt{H⁰Γ ℋ\textasciicircum{}h = Hom\_\{Ho(P)\}(M₁, M₂)} ((a)), \texttt{H⁰Γ 𝓗om\textasciicircum{}tw(J\_T M₁, I) = K(TwMod)(J\_T M₁, I) = D(TwMod)(J\_T M₁, I)} ((e); \texttt{I} K-injective); the composite is \texttt{f ↦ J\_T(f) ≫ Q(r)}. Stated for resolving data \texttt{D} (compact Hausdorff \texttt{X}) with LL over opens as a hypothesis (\texttt{LLOpen}); for the Dolbeault data it is proved (\texttt{llOpen\_dolbeault}, from Z3.e), so the Dolbeault theorems carry no such hypothesis.
\paragraph{Blueprint notes.}
F.1–F.5

Files: \texttt{Comparison/TwistedFullyFaithful*.lean}.
\end{proof}

\begin{theorem}[Z3.t — Curved model]
\label{Z3.t}
\lean{Cohesive.twModel_TD_comm, Cohesive.isRescale_untwist, Cohesive.untwistRestrIso, Cohesive.untwist_compat, Cohesive.untwist_ext, Cohesive.πTot}
\leanok
Let P be pre-resolving data on X, h ∈ 𝒜(X)\textsuperscript{2} closed, T a twist datum for h and K′ a complex of ȟ-twisted O-modules.

\begin{itemize}
\item The cocycle ρ(ȟ) of 𝒜\textsuperscript{0}-units, rescaled by u\textsuperscript{−1}, is trivial; hence untwist(M) = glue(rescale\_\{u\textsuperscript{−1}\}(𝒜\textsuperscript{0} ⊗ M)) is a sheaf of 𝒜\textsuperscript{0}-modules on X with untwist(M)|U\textsubscript{α} ≅ 𝒜\textsuperscript{0}|U\textsubscript{α} ⊗ M\textsubscript{α}, whose sections are determined by their components, which satisfy s\textsubscript{β} = u\textsubscript{αβ}\textsuperscript{−1}·φ\textsubscript{αβ}(s\textsubscript{α}).
\item For X compact and connections θ\textsubscript{α} on 𝒜\textsuperscript{0}|U\textsubscript{α} ⊗ K′\textsubscript{α} (C7), the differentials of the local models commute with the gluing, giving a twisted model complex of K′.
\item For F = untwist(K′), the maps Tot(𝒜(X), F(X)) → Tot(𝒜(U\textsubscript{α}), F\textsubscript{α}(U\textsubscript{α})) jointly determine elements, and there is a curved operator θ\textsuperscript{h} with (θ\textsuperscript{h})\textsuperscript{2} = h restricting to θ\textsubscript{α} + h\textsuperscript{0}\textsubscript{α}.
\end{itemize}
\par\smallskip{\footnotesize Source: blueprint-tw.md §4.}
\end{theorem}
\begin{proof}
\leanok
\uses{C7, Z3.k, Z3.lprime, Z3.o}
\textbf{G.5} curved model: \texttt{∇\textasciicircum{}h\_α := ∇\_α + h0\_α ⊗ −} glue to \texttt{∇\textasciicircum{}h} with \texttt{(∇\textasciicircum{}h)² = h}; \texttt{N := (Γ(F\textasciicircum{}•), 𝕏\textasciicircum{}h)} quasi-cohesive over \texttt{(𝒜(X), h)} (curved C7.c copy).

\textbf{G.7} \texttt{e\_α ∈ Z0(gauge(M|U\_α), N\_α)}.
\paragraph{Blueprint notes.}
G.5, G.7

Files: \texttt{Comparison/CurvedConn.lean}, \texttt{TwistedModel.lean}.
\end{proof}

\begin{theorem}[Z3.u — Essential surjectivity]
\label{Z3.u}
\lean{Cohesive.exists_twistedJ_obj_iso, Cohesive.TwIsPerfect.exists_flat_model, Cohesive.isKFlat_twComp, Cohesive.exists_roof_sheafBC_restrOver, Cohesive.untwFRestrIso, Cohesive.isPerfect_untwF}
\leanok
Let D be resolving data on a compact space X, h ∈ 𝒜(X)\textsuperscript{2} closed and T a twist datum for h. Then every object K of D(TwMod(O, U, ȟ)) in D\textsubscript{perf}(X, ȟ) is isomorphic to J\textsubscript{T}(M) for some M in Ho(P\textsubscript{(𝒜(X), h)}). Steps proved on the way:

\begin{itemize}
\item every twisted perfect complex receives a quasi-isomorphism from a strictly bounded above complex K′ with componentwise flat terms, whose components K′\textsubscript{α} are then K-flat;
\item for such K′ (over pre-resolving data), F = untwist(K′) satisfies F|U\textsubscript{α} ≅ 𝒜\textsuperscript{0}|U\textsubscript{α} ⊗ K′\textsubscript{α} and is a perfect complex of 𝒜\textsuperscript{0}-modules on X;
\item base change along R → S of a bounded above flat complex preserves the existence of local roofs to strictly perfect complexes.
\end{itemize}
Deviation: K-flat variant (no flatness of 𝒜 over O assumed); X is not assumed Hausdorff here.
\par\smallskip{\footnotesize Source: blueprint-tw.md §4.}
\end{theorem}
\begin{proof}
\leanok
\uses{A11, A3, C7, S10, S3, S4, S8, S9, Z3.j, Z3.l, Z3.t}
\textbf{G.1} \texttt{D ≅ Q K}, bounded (C.6), \texttt{K' → K} componentwise flat (C.8).

\textbf{G.2} \texttt{F := glue(𝒜⁰ ⊗ K'\_α, u⁻¹(1 ⊗ φ))} ([A, Rem 7.2.4]).

\textbf{G.3} \texttt{F} globally bounded perfect on \texttt{X}.

\textbf{G.4} S10: \texttt{P̃ → F}, \texttt{P := Γ(P̃)}.

\textbf{G.6} A11 (general \texttt{c}, §1.2): \texttt{M}, \texttt{e ∈ Z0(M, N)}.

\textbf{G.8} \texttt{Ξ} glued, componentwise quasi-iso by C7.f.

\textbf{G.9} \texttt{ιT : K' → T\textasciicircum{}tw(K')} (C7.d).

\textbf{G.10} \texttt{Q(J\_T M) ≅ D}. \textbf{G′} (detail, 5 Nov 2026, for review at G5; refines G.1–G.10, no decision changed). Fix flat resolving data \texttt{P}, \texttt{h} closed of degree 2, a twist datum \texttt{T} (\texttt{h⁰\_α}, \texttt{u\_αβ}, \texttt{ȟ}), and \texttt{K} twisted perfect in \texttt{TwMod O U ȟ} with \texttt{D ≅ Q K}. \texttt{B\_α := 𝒜⁰|U\_α ⊗ −}, \texttt{η(k) = 1 ⊗ k} as in C.7′. \textbf{(a)} \emph{Flat model} (G.1): \texttt{K} has homology in \texttt{[a, b]} (C.6 (g)); \texttt{K' → K} a quasi-isomorphism from a bounded above complex with flat components (Z3.l, applied to \texttt{truncLE b}). \texttt{K'} is twisted perfect (C.6 (e)) and each \texttt{K'\_α} is K-flat on \texttt{Over U\_α} (S8.e). \textbf{(b)} \emph{Base change, untwisting} (G.2): \texttt{B := twBC ρ K'} (C.7′, cocycle \texttt{ρ(ȟ)}, components \texttt{B\_α(K'\_α) = F\_\{K'\_α\}} of C7.b by \texttt{rfl}). By B.5, \texttt{ρ(ȟ\_αβγ) = u\_βγ u\_αβ u\_αγ⁻¹}, so \texttt{g\_αβ := u\_αβ⁻¹} is a rescaling \texttt{ρ(ȟ) → 1} (\texttt{IsRescale}: \texttt{g\_αγ = g\_βγ g\_αβ ρ(ȟ\_αβγ)} ✓). \texttt{F := glue(rescale\_g B)}, a complex of \texttt{𝒜⁰}-modules on \texttt{X}, with gluing \texttt{φ\textasciicircum{}F\_αβ = u\_αβ⁻¹·φ\textasciicircum{}B\_αβ} and \texttt{φ\textasciicircum{}B(η k) = η(φ k)} (\texttt{bcφ\_bcUnit}). \texttt{F|U\_α ≅ F\_\{K'\_α\}} (C.7 (c)), compatible with the differentials. \textbf{(c)} \emph{Perfectness} (G.3): \texttt{F\_\{K'\_α\}} is perfect on \texttt{Over U\_α}: S8.g's \texttt{exists\_roof\_sheafBC} on over-sites, with C.7′ (a) in place of \texttt{sheafBCOverIso}. Hence \texttt{F} is perfect on \texttt{X} (local roofs through \texttt{F|U\_α ≅ F\_\{K'\_α\}} and \texttt{restr\_restrOver}), globally bounded (S9.d, \texttt{X} compact). \textbf{(d)} \emph{S10} (G.4): \texttt{Pt} strictly perfect over \texttt{𝒜⁰}, \texttt{f⁰ : Pt → F} a quasi-isomorphism; \texttt{P := Γ(Pt)} is a bounded complex of finitely generated projective \texttt{𝒜⁰(X)}-modules (S4.c) and \texttt{Γ(f⁰)} is a quasi-isomorphism (S3, \texttt{𝒜⁰} fine). \textbf{(e)} \emph{Local connections}: \texttt{θ\_α} the C7.b connection on \texttt{F\_\{K'\_α\}} over \texttt{Over U\_α} (\texttt{exists\_conn}; \texttt{θ\_α(η k) = 1 ⊗ η(d k)}), \texttt{T\_α := θ\_α.TC}, \texttt{ιT\_α : K'\_α → T\_α} a quasi-isomorphism (C7.d, \texttt{quasiIso\_ιT}, K-flat). \textbf{(f)} \emph{Twisted model} (G.9): \texttt{T\textasciicircum{}tw ∈ TwMod O U ȟ} (complexes) with components \texttt{T\_α} and gluing \texttt{a ⊗ s ↦ a ⊗ φ\textasciicircum{}B\_αβ(s)} on \texttt{Tot(𝒜(W), F\_\{K'\_α\}(W))} (\texttt{totMapAlong} with the identity of \texttt{𝒜(W)}; \texttt{φ\textasciicircum{}B} is \texttt{𝒜⁰(W)}-linear). Cocycle \texttt{ȟ} (acting through \texttt{ρ}). Chain maps: \texttt{θ\_β φ\textasciicircum{}B} and \texttt{φ\textasciicircum{}B θ\_α} are both \texttt{𝒜⁰}-connections, so their difference is \texttt{𝒜⁰}-linear out of \texttt{F\_\{K'\_α\}}, and it vanishes on \texttt{η k} (\texttt{d\_\{K'\}} commutes with \texttt{φ}), hence everywhere (C.7′ (b)). \texttt{ιT\textasciicircum{}tw : K' → T\textasciicircum{}tw} has components \texttt{ιT\_α} (commutes with the gluing: \texttt{φ\textasciicircum{}B η = η φ}) and is a quasi-isomorphism componentwise (C.6 (e)). \textbf{(g)} *Sections of the model of \texttt{F}* (C.7 (g)): for \texttt{V} open, the map \texttt{Tot(𝒜(V), F(V)) → ∏\_α Tot(𝒜(V ⊓ U\_α), F\_\{K'\_α\}(V ⊓ U\_α))}, restriction followed by the component \texttt{α} of \texttt{F|U\_α ≅ F\_\{K'\_α\}} (\texttt{π\_α}), is injective, and its image consists of the families with \texttt{y\_β = u\_αβ⁻¹·φ\textasciicircum{}T\_αβ(y\_α)} on the overlaps. Proof: \texttt{V ↦ Tot(𝒜(V), F(V))} is a sheaf (C7.a, for any \texttt{F}), and (b). \textbf{(h)} \emph{Curved connection} (G.5): for \texttt{s ∈ F\textasciicircum{}p(X)}, \texttt{θ\textasciicircum{}h(s)} is the element of (g) with components \texttt{θ\_α(s\_α) + h⁰\_α·s\_α}. Compatibility on \texttt{W ≤ U\_α ∩ U\_β}: on \texttt{s\_α = η k}, \texttt{θ\_β(u⁻¹η(φk)) + h⁰\_β u⁻¹η(φk) = u⁻¹(η(φ dk) + h⁰\_α η(φk))}, since \texttt{d(u\_αβ⁻¹) = −u\_αβ⁻¹(h⁰\_β − h⁰\_α)} (from \texttt{d u\_αβ = (h⁰\_β − h⁰\_α)u\_αβ}); in general both sides are connections (C.7′ (b)). \texttt{θ\textasciicircum{}h} is generator data over \texttt{(𝒜(X), h)} (A3.a): degree 1; flat, since in each component \texttt{(θ\_α + h⁰\_α)² = θ\_α² + d h⁰\_α = h|U\_α} (\texttt{h⁰\_α} odd); its form-degree-0 part is \texttt{d\_F} (\texttt{h⁰\_α} has form degree 1). \texttt{N\textasciicircum{}h := (F\textasciicircum{}•(X), θ\textasciicircum{}h)} is quasi-cohesive over \texttt{(𝒜(X), h)} with \texttt{(0,0)}-complex \texttt{Γ(F)}. Only \texttt{V = X} is needed: no curved copy of C7's \texttt{Conn} for all \texttt{V}. \textbf{(i)} \emph{Lifting} (G.6): A11 (general \texttt{c}) with \texttt{e⁰ := Γ(f⁰)} gives \texttt{M} cohesive over \texttt{(𝒜(X), h)} on \texttt{𝒜(X) ⊗ P} and \texttt{e ∈ Z0(M, N\textasciicircum{}h)} with \texttt{(0,0)}-component \texttt{e⁰} (curved copy of \texttt{exists\_cohesive\_gr0\_quasiIso}, steps (A)–(C)). \textbf{(j)} \emph{Localization} (G.7): \texttt{e\_α ∈ Z0(gauge\_\{h⁰\_α\}(M|U\_α), N\_α)}, \texttt{N\_α := θ\_α.N(U\_α)}, \texttt{x ↦ π\_α(extAlong(res) e (x))}. Closed: \texttt{π\_α ∘ res ∘ θ\textasciicircum{}h = (θ\_α + h⁰\_α) ∘ π\_α ∘ res} by (h), and the gauged differential is \texttt{𝕏 − h⁰\_α} (A.2). \textbf{(k)} \emph{Comparison} (G.8): \texttt{Ξ\textasciicircum{}tw : J\_T(M) → T\textasciicircum{}tw} with components \texttt{twLocalIso} (D.6, \texttt{W′ = U\_α}) followed by \texttt{Ξ\_\{e\_α\}} (C7.e). It commutes with the gluing: \texttt{J\_T}'s gluing is \texttt{u\_αβ·}, and \texttt{φ\textasciicircum{}T(π\_α y) = u\_αβ·π\_β(y)} by (g). Quasi-isomorphism componentwise by C7.f: \texttt{gr0Ξ\_\{e\_α\}} corresponds to \texttt{f⁰|U\_α} under \texttt{G0(gauge(M|U\_α)) ≅ Pt|U\_α} and \texttt{F|U\_α ≅ F\_\{K'\_α\}} (gauging changes only the form-degree-1 component), as in C5's step (D) with base \texttt{U\_α} in place of \texttt{⊤}. In detail: the generators of \texttt{gauge(M|U\_α)} are \texttt{𝒜⁰(U\_α) ⊗\_\{𝒜⁰(X)\} Γ(Pt\textasciicircum{}p)}, and \texttt{φ\_p : β ⊗ q ↦ β·q|U\_α} onto \texttt{Pt\textasciicircum{}p(U\_α)} is bijective because \texttt{Pt\textasciicircum{}p} is a summand of a finite free \texttt{𝒜⁰}-module (true for \texttt{𝒜⁰\textasciicircum{}m}, and preserved by retracts); \texttt{φ} intertwines \texttt{X0comp} with the differential of \texttt{Pt|U\_α} (gauging does not change \texttt{X0comp}). Then C5 (D) with base \texttt{U\_α} (\texttt{quasiIso\_gr0Ξ\_of\_equiv}, \texttt{Sc := Pt|U\_α}, \texttt{g := f⁰|U\_α} followed by \texttt{F|U\_α ≅ F\_\{K'\_α\}}) gives \texttt{gr0Ξ\_\{e\_α\}} a quasi-isomorphism; \texttt{hg} is \texttt{ε ∘ π\_α = π\_α ∘ ε} on \texttt{e(1 ⊗ q)} and the \texttt{(0,0)}-component \texttt{Γ(f⁰)} of \texttt{e} from (i). \textbf{(l)} \emph{Assembly} (G.10): \texttt{Q(J\_T M) ≅ Q(T\textasciicircum{}tw) ≅ Q(K') ≅ Q(K) ≅ D}. \textbf{(m)} \emph{Split and risks}: Z3.t = (b), (e)–(h), (j); Z3.u = (a), (c), (d), (i), (k), (l). Risks: (k) generalizing C5's step (D) (\texttt{quasiIso\_gr0Ξ\_abstract}, private, at \texttt{⊤}) to base \texttt{U\_α}; (g) the transport of \texttt{T(F)} sections along \texttt{F|U\_α ≅ F\_\{K'\_α\}}; curved/flat instance traps between \texttt{(𝒜(X), h)} and the flat \texttt{𝒜(U\_α)} in (h)–(j) (P15 notes; \texttt{cmpTot} as in \texttt{twCmp}).
\paragraph{Blueprint notes.}
G.1–G.4, G.6, G.8–G.10

Files: \texttt{Comparison/TwistedEssSurj*.lean}.
\end{proof}

\begin{theorem}[Z3.v — Main theorem twisted]
\label{Z3.v}
\lean{Cohesive.twistedBlockFunctor, Cohesive.MainStatementTwisted, Cohesive.main_theorem_twisted, Cohesive.essImage_twistedBlockFunctor}
\leanok
Let E be a finite-dimensional complex normed space and X a compact Hausdorff complex manifold modelled on E. Let h ∈ 𝒜\textsuperscript{0,2}(X) be ∂̄-closed, A = (𝒜\textsubscript{X}(X), ∂̄, h) the curved Dolbeault algebra and T a twist datum for h with cover (U\textsubscript{α}) and cocycle ȟ. Then the functor J\textsubscript{T} : Ho(P\textsubscript{A}) ⥤ D(TwMod(O\textsubscript{X}, U, ȟ)), M ↦ J\textsubscript{T}(M), into the derived category of ȟ-twisted sheaves of O\textsubscript{X}-modules is full and faithful, and its essential image is exactly D\textsubscript{perf}(X, ȟ) (the essential image statement holds for any choice of derived category).

Deviation: stated for each twist datum T (independence of T is not proved); D\textsubscript{perf} instead of D\textsuperscript{b\_}\{coh\}; twisted O\textsubscript{X}-modules only, curvature c = h for left modules, no factor 2πi.
\par\smallskip{\footnotesize Source: blueprint-tw.md §4.}
\end{theorem}
\begin{proof}
\leanok
\uses{Z3.s, Z3.u}
\textbf{H.1} \texttt{twistedComparison} for \texttt{CurvedResolvingData} (u).

\textbf{H.2} Dolbeault instance, \texttt{main\_theorem\_twisted}.

\textbf{H.3} K8: at \texttt{(0, trivialTwistDatum)} the source \texttt{HoP} is literally \texttt{HoP (dolbeault E X)} (P15); \texttt{glueEquiv : TwMod O \{⊤\} 1 ≌ Mod O}, \texttt{twistedBlockFunctor E X 0 T₀ ≅ blockFunctor ⋙ glue.mapDerivedCategory}-type isomorphism (\texttt{Equivalence.mapDerivedCategory} ≈ 40 lines from \texttt{mapDerivedCategoryCompIso}/\texttt{IdIso}/\texttt{NatTrans.mapDerivedCategory}), \texttt{TwPerf ↔ Perf}.
\paragraph{Blueprint notes.}
H.1–H.3

Files: \texttt{Comparison/TwistedEquivalence.lean}, \texttt{TwistedBlockFunctorDef.lean}, \texttt{StatementsTwisted.lean}, \texttt{MainTwisted.lean}, \texttt{ChecksTwisted.lean}.
\end{proof}


\begin{thebibliography}{TauCeti}
\bibitem[A]{A} P.~Antweiler, \emph{Curved $\infty$-local systems and projectively flat
Riemann--Hilbert correspondence}, arXiv:2411.19595v1, 2024.

\bibitem[B]{B} J.~Block, \emph{Duality and equivalence of module categories in noncommutative
geometry}, in: A celebration of the mathematical legacy of Raoul Bott, CRM Proc.\ Lecture
Notes~50, Amer.\ Math.\ Soc., 2010; arXiv:math/0509284v2 (numbering of the arXiv version).

\bibitem[BK]{BK} A.~Bondal, M.~Kapranov, \emph{Enhanced triangulated categories}, Math.\ USSR
Sb.~70 (1991), 93--107.

\bibitem[BSW]{BSW} J.-M.~Bismut, S.~Shen, Z.~Wei, \emph{Coherent sheaves, superconnections, and
Riemann--Roch--Grothendieck}, Progr.\ Math.~347, Birkhäuser, 2023; arXiv:2102.08129.

\bibitem[Că]{Cal} A.~Căldăraru, \emph{Derived categories of twisted sheaves on Calabi--Yau
manifolds}, PhD thesis, Cornell University, 2000.

\bibitem[CHL]{CHL} J.~Chuang, J.~Holstein, A.~Lazarev, \emph{Maurer--Cartan moduli and theorems
of Riemann--Hilbert type}, Appl.\ Categ.\ Structures~29 (2021); arXiv:1802.02549v3 (numbering of
the arXiv version).

\bibitem[DK]{DK} S.~K.~Donaldson, P.~B.~Kronheimer, \emph{The geometry of four-manifolds}, Oxford
University Press, 1990.

\bibitem[GH]{GH} P.~Griffiths, J.~Harris, \emph{Principles of algebraic geometry}, Wiley, 1978.

\bibitem[Hö]{Hor} L.~Hörmander, \emph{An introduction to complex analysis in several variables},
3rd ed., North-Holland, 1990.

\bibitem[Mathlib]{Mathlib} The mathlib Community, \emph{The Lean mathematical library}, in: Proc.\
CPP 2020, ACM, 2020, 367--381; version \texttt{v4.35.0}.

\bibitem[TauCeti]{TauCeti} The TauCeti project, Lean formalization library,
\url{https://github.com/TauCetiProject/TauCeti} (branch \texttt{main}, fetched 1 November 2026;
Apache License~2.0).

\bibitem[Toda]{Toda} Y.~Toda, \emph{Deformations and Fourier--Mukai transforms}, J.\ Differential
Geom.~81 (2009), 197--224.
\end{thebibliography}

\end{document}
